% Limitation de la fréquence de certaines maladies géniques et/ou chromosomiques au sein des familles — SVT Tle D — Cours signé OnBuch+
% Programme officiel MINESEC. Compiler avec : tectonic N06-hérédité-humaine-prevision-genetique.tex
\newcommand{\DOCMATIERE}{SVT}
\newcommand{\DOCNIVEAU}{Tle D}
\newcommand{\DOCMODULE}{Module I : Le Monde Vivant — Prévision en génétique humaine}
\newcommand{\DOCLECON}{N06}
\newcommand{\DOCTITRE}{Limitation de la fréquence de certaines maladies géniques et/ou chromosomiques au sein des familles}
% OnBuch+ — préambule « cours signés » ENRICHI (compilé avec tectonic / XeLaTeX)
% Système visuel « Pop » d'OnBuch+ : fond crème (jamais blanc), bordures encre
% épaisses, ombres « dures » décalées, titres Archivo Black en capitales.
% Version MATHÉMATIQUES : graphes (pgfplots/tikz), boîtes pédagogiques
% (définition, propriété, méthode, attention, le savais-tu, exercice résolu,
% exercice, TP, bilan), page de garde et signature OnBuch+ ancrée en bas.
%
% Macros attendues dans main.tex AVANT \input{preamble} :
%   \DOCMATIERE  \DOCNIVEAU  \DOCMODULE  \DOCLECON (numéro)  \DOCTITRE
\documentclass[11pt]{article}

\usepackage{fontspec}
\usepackage[french]{babel}
\usepackage[a4paper,margin=2cm,top=3.4cm,bottom=2.8cm,headheight=1.4cm,footskip=1.3cm]{geometry}
\usepackage{amsmath,amssymb,amsfonts}
\usepackage{mathtools} % \xmapsto, \coloneqq... souvent utilisés par les modèles
\usepackage{cancel} % simplifications barrées (bilans de forces)
% \abs{z} (module, valeur absolue) et \norm{u} (norme), utilisés par les modèles
\providecommand{\abs}{}\let\abs\relax\DeclarePairedDelimiter{\abs}{\lvert}{\rvert}
\providecommand{\norm}{}\let\norm\relax\DeclarePairedDelimiter{\norm}{\lVert}{\rVert}
\usepackage[table]{xcolor} % [table] : fournit aussi \rowcolors (alternance de lignes)
\usepackage{pagecolor}
\usepackage{tikz}
\usetikzlibrary{arrows.meta,positioning,calc,shapes.geometric,shapes.misc,shapes.arrows,decorations.pathmorphing,decorations.pathreplacing,decorations.markings,patterns,shadows,mindmap,trees,fit,backgrounds,matrix,intersections,angles,quotes}
\usepackage{pgfplots}
\usepackage[version=4]{mhchem} % équations nucléaires \ce{^{235}_{92}U}
\usepackage[european,siunitx]{circuitikz} % schémas électriques
\pgfplotsset{compat=1.18}
\usepgfplotslibrary{fillbetween,groupplots} % aires entre courbes (calcul intégral)
\usepackage{enumitem}
\usepackage{fancyhdr}
% [most] charge toutes les bibliothèques tcolorbox (listings, minted, raster,
% documentation...) et gonfle inutilement la mémoire TeX sur un document long
% (~300 boîtes) : on ne charge que ce qui est réellement utilisé.
\usepackage[skins,breakable]{tcolorbox}
\usepackage{graphicx}
\usepackage{colortbl}
\usepackage{booktabs}
\usepackage{tabularx}
\usepackage{diagbox} % case d'en-tête diagonale des tableaux à double entrée
% Colonnes extensibles centrée / gauche / droite (C, L, R), souvent utilisées par les modèles
\newcolumntype{C}{>{\centering\arraybackslash}X}
\newcolumntype{L}{>{\raggedright\arraybackslash}X}
\newcolumntype{R}{>{\raggedleft\arraybackslash}X}
\usepackage{array}
\usepackage{multicol}
\usepackage{siunitx}
% parse-numbers=false : accepte aussi \SI{2,5 \times 10^{-3}}{...}
\sisetup{locale=FR,output-decimal-marker={,},group-minimum-digits=4,parse-numbers=false}
\DeclareSIUnit\ohm{\ensuremath{\symup{\Omega}}} % U+2126 absent de la police
\DeclareSIUnit\division{div} % réglages d'oscilloscope (ms/div, V/div)
\DeclareSIUnit\tr{tr} % tours (vitesse de rotation, ex. \SI{1500}{\tr\per\minute})
\DeclareSIUnit\molar{M} % concentration molaire (ex. \SI{3}{\molar}, \SI{1}{\micro\molar})
\DeclareSIUnit\an{an} % année (le modèle confond parfois avec \year, un compteur TeX) — ex. \SI{5}{\kilo\meter\per\an}
\DeclareSIUnit\FCFA{FCFA} % franc CFA (ex. \SI{500}{\FCFA\per\kilo\gram})
\DeclareSIUnit\degreeBrix{\textdegree Brix} % teneur en sucre d'un sirop/jus (réfractométrie)
\DeclareSIUnit\month{mois} % le modèle utilise parfois \month, un compteur TeX, comme unité de durée
\usepackage{qrcode}
% \degree : commande fréquemment utilisée par les modèles pour un angle ou une
% température en dehors de \SI{}{} (non fournie par les paquets chargés).
\providecommand{\degree}{^{\circ}}
% \qed (fin de démonstration), utilisé par les modèles sans amsthm
\providecommand{\qed}{\hfill$\square$}
% \barre : double barre des valeurs interdites dans un tableau de variations (tkz-tab)
\providecommand{\barre}{\ensuremath{\|}}
% environnement proof (amsthm non chargé) utilisé par les modèles
\ifdefined\proof\else\newenvironment{proof}[1][Démonstration]{\par\noindent\textit{#1.}\ }{\qed\par}\fi
\usepackage{lastpage}
\usepackage{hyperref}
\hypersetup{hidelinks}

% ============================================================
% Palette « Pop » OnBuch+ — mêmes hex que la classe P (plus_ui.dart)
% ============================================================
\definecolor{poporange}{HTML}{F59321}
\definecolor{popdark}{HTML}{E07A0C}
\definecolor{popcream}{HTML}{FFF6E8}
\definecolor{popink}{HTML}{1C1714}
\definecolor{poppurple}{HTML}{7A5AE0}
\definecolor{popgreen}{HTML}{2E9E6A}
\definecolor{poppink}{HTML}{E0457B}
\definecolor{popblue}{HTML}{2F7DE1}
\definecolor{popgold}{HTML}{C98A12}
\definecolor{popmuted}{HTML}{8A7F76}
\definecolor{popline}{HTML}{EADFD2}
\definecolor{popgray}{HTML}{8A7F76}
\colorlet{popgrey}{popgray}
% Alias tolérants (le modèle nomme parfois les couleurs différemment)
\colorlet{popwhite}{white}
\colorlet{popblack}{popink}
\colorlet{popred}{poppink}
\colorlet{popyellow}{popgold}
% Teintes claires (fonds de tableaux/encadrés) — le modèle nomme parfois
% les couleurs avec un suffixe L (ex. popblueL) sans les avoir définies.
\colorlet{poporangeL}{poporange!20}
\colorlet{popdarkL}{popdark!20}
\colorlet{poppurpleL}{poppurple!20}
\colorlet{popgreenL}{popgreen!20}
\colorlet{poppinkL}{poppink!20}
\colorlet{popblueL}{popblue!20}
\colorlet{popgoldL}{popgold!20}
\colorlet{popredL}{popred!20}
\colorlet{popgrayL}{popgray!20}
% Teintes claires pour fonds de tableaux / zones
\colorlet{popblueL}{popblue!10!white}
\colorlet{poporangeL}{poporange!14!white}
\colorlet{popgreenL}{popgreen!12!white}
\colorlet{poppurpleL}{poppurple!12!white}
\colorlet{poppinkL}{poppink!12!white}
\colorlet{popgoldL}{popgold!14!white}

\pagecolor{popcream}

% ============================================================
% Typographie
% ============================================================
\defaultfontfeatures{Ligatures=TeX}
\setmainfont{PlusJakartaSans-Medium.ttf}[Path=../../../fonts/,BoldFont=PlusJakartaSans-Bold.ttf]
% Maths en sans-serif (Fira Math) avec chiffres et lettres droites de Plus
% Jakarta, pour que les formules se fondent dans le texte.
\usepackage[math-style=ISO,bold-style=ISO]{unicode-math}
\setmathfont{FiraMath-Regular.otf}[Path=../../../fonts/]
\setmathfont{PlusJakartaSans-Medium.ttf}[Path=../../../fonts/,range={up/num,up/{Latin,latin},\mathup}]
\usepackage{icomma} % virgule décimale sans espace en mode math
% Caractères mathématiques Unicode absents des polices de texte (Plus Jakarta,
% Archivo Black) : rendus par leur équivalent mathématique, en texte comme en
% formule — sinon ils s'affichent en carré vide (ex. « Arithmétique dans ℤ »).
\usepackage{newunicodechar}
\newunicodechar{ℝ}{\ensuremath{\mathbb{R}}}
\newunicodechar{ℤ}{\ensuremath{\mathbb{Z}}}
\newunicodechar{ℂ}{\ensuremath{\mathbb{C}}}
\newunicodechar{ℕ}{\ensuremath{\mathbb{N}}}
\newunicodechar{ℚ}{\ensuremath{\mathbb{Q}}}
\newunicodechar{𝔻}{\ensuremath{\mathbb{D}}}
\newunicodechar{α}{\ensuremath{\alpha}}
\newunicodechar{β}{\ensuremath{\beta}}
\newunicodechar{γ}{\ensuremath{\gamma}}
\newunicodechar{δ}{\ensuremath{\delta}}
\newunicodechar{ε}{\ensuremath{\varepsilon}}
\newunicodechar{θ}{\ensuremath{\theta}}
\newunicodechar{λ}{\ensuremath{\lambda}}
\newunicodechar{μ}{\ensuremath{\mu}}
\newunicodechar{σ}{\ensuremath{\sigma}}
\newunicodechar{τ}{\ensuremath{\tau}}
\newunicodechar{φ}{\ensuremath{\varphi}}
\newunicodechar{ψ}{\ensuremath{\psi}}
\newunicodechar{ω}{\ensuremath{\omega}}
\newunicodechar{Σ}{\ensuremath{\Sigma}}
% majuscules produites par \MakeUppercase dans les titres (α-aminés -> Α)
\newunicodechar{Α}{\ensuremath{\alpha}}
\newunicodechar{Β}{\ensuremath{\beta}}
\newunicodechar{Γ}{\ensuremath{\Gamma}}
\newunicodechar{Δ}{\ensuremath{\Delta}}
\newunicodechar{Ε}{\ensuremath{\varepsilon}}
\newunicodechar{Θ}{\ensuremath{\Theta}}
\newunicodechar{Λ}{\ensuremath{\Lambda}}
\newunicodechar{Μ}{\ensuremath{\mu}}
\newunicodechar{Τ}{\ensuremath{\tau}}
\newunicodechar{Φ}{\ensuremath{\Phi}}
\newunicodechar{Ψ}{\ensuremath{\Psi}}
\newunicodechar{Ω}{\ensuremath{\Omega}}
\newunicodechar{↦}{\ensuremath{\mapsto}}
\newunicodechar{∈}{\ensuremath{\in}}
\newunicodechar{∉}{\ensuremath{\notin}}
\newunicodechar{∧}{\ensuremath{\wedge}}
\newunicodechar{⇔}{\ensuremath{\Leftrightarrow}}
\newunicodechar{⟺}{\ensuremath{\Longleftrightarrow}}
\newunicodechar{⇒}{\ensuremath{\Rightarrow}}
\newunicodechar{⟹}{\ensuremath{\Longrightarrow}}
\newunicodechar{∘}{\ensuremath{\circ}}
\newunicodechar{∪}{\ensuremath{\cup}}
\newunicodechar{∩}{\ensuremath{\cap}}
\newunicodechar{≡}{\ensuremath{\equiv}}
\newunicodechar{⊂}{\ensuremath{\subset}}
\newunicodechar{⊥}{\ensuremath{\perp}}
\newunicodechar{∃}{\ensuremath{\exists}}
\newunicodechar{∀}{\ensuremath{\forall}}
\newunicodechar{∛}{\ensuremath{\sqrt[3]{\,}}}
\newunicodechar{∜}{\ensuremath{\sqrt[4]{\,}}}
\newunicodechar{⃗}{\ensuremath{{}^{\to}}}
\newunicodechar{ⁿ}{\ifmmode^{n}\else\textsuperscript{\ensuremath{n}}\fi}
\newunicodechar{ˣ}{\ifmmode^{x}\else\textsuperscript{\ensuremath{x}}\fi}
\newunicodechar{ᵇ}{\ifmmode^{b}\else\textsuperscript{\ensuremath{b}}\fi}
\newunicodechar{ʳ}{\ifmmode^{r}\else\textsuperscript{\ensuremath{r}}\fi}
\newunicodechar{ᵘ}{\ifmmode^{u}\else\textsuperscript{\ensuremath{u}}\fi}
\newunicodechar{ʸ}{\ifmmode^{y}\else\textsuperscript{\ensuremath{y}}\fi}
\newunicodechar{ᵃ}{\ifmmode^{a}\else\textsuperscript{\ensuremath{a}}\fi}
\newunicodechar{ᵉ}{\ifmmode^{e}\else\textsuperscript{\ensuremath{e}}\fi}
\newunicodechar{ᵏ}{\ifmmode^{k}\else\textsuperscript{\ensuremath{k}}\fi}
\newunicodechar{ᵖ}{\ifmmode^{p}\else\textsuperscript{\ensuremath{p}}\fi}
\newunicodechar{ᴺ}{\ifmmode^{N}\else\textsuperscript{\ensuremath{N}}\fi}
\newunicodechar{ᶜ}{\ifmmode^{c}\else\textsuperscript{\ensuremath{c}}\fi}
\newunicodechar{ʰ}{\ifmmode^{h}\else\textsuperscript{\ensuremath{h}}\fi}
\newunicodechar{ᵠ}{\ifmmode^{q}\else\textsuperscript{\ensuremath{q}}\fi}
\newunicodechar{ₙ}{\ifmmode_{n}\else\textsubscript{\ensuremath{n}}\fi}
\newunicodechar{ₐ}{\ifmmode_{a}\else\textsubscript{\ensuremath{a}}\fi}
\newunicodechar{ᵢ}{\ifmmode_{i}\else\textsubscript{\ensuremath{i}}\fi}
\newunicodechar{ᵣ}{\ifmmode_{r}\else\textsubscript{\ensuremath{r}}\fi}
\newunicodechar{ₖ}{\ifmmode_{k}\else\textsubscript{\ensuremath{k}}\fi}
\newunicodechar{ᵦ}{\ifmmode_{\beta}\else\textsubscript{\ensuremath{\beta}}\fi}
\newunicodechar{ₓ}{\ifmmode_{x}\else\textsubscript{\ensuremath{x}}\fi}
\newfontfamily\popdisplay{ArchivoBlack-Regular.ttf}[Path=../../../fonts/]
\newfontfamily\poplogo{SpaceGrotesk-Bold.ttf}[Path=../../../fonts/]
\newfontfamily\popsemi{PlusJakartaSans-SemiBold.ttf}[Path=../../../fonts/]
\newfontfamily\popxbold{PlusJakartaSans-ExtraBold.ttf}[Path=../../../fonts/]
\newfontfamily\popxboldls{PlusJakartaSans-ExtraBold.ttf}[Path=../../../fonts/,LetterSpace=3.0]

\setlist[enumerate]{leftmargin=1.7em,itemsep=5pt,topsep=4pt}
\setlist[itemize]{leftmargin=1.7em,itemsep=4pt,topsep=4pt,label={\color{poporange}\textbullet}}
\setlength{\parindent}{0pt}
\setlength{\parskip}{5pt}
\renewcommand{\arraystretch}{1.25}

% pgfplots : style Pop par défaut
\pgfplotsset{
  every axis/.append style={axis line style={popink,line width=1pt},tick style={popink},
    grid style={popline,line width=0.5pt},label style={font=\small},tick label style={font=\footnotesize}},
  popcurve/.style={poporange,line width=1.8pt,no markers,smooth},
}
% Même style disponible dans un \draw TikZ simple (hors pgfplots)
\tikzset{popcurve/.style={poporange,line width=1.8pt}}
\tikzset{popgrid/.style={popline,very thin}}
\colorlet{popcurve}{poporange} % le nom du style est parfois utilisé comme couleur

% ============================================================
% Pill / squiggle
% ============================================================
\newcommand{\poppill}[3]{%
  \tikz[baseline=-0.55ex]\node[draw=popink,line width=1.2pt,rounded corners=2.6pt,%
    fill=#2,inner xsep=6.5pt,inner ysep=3pt]{%
    \popxboldls\fontsize{7}{7}\selectfont\color{#3}\MakeUppercase{#1}};%
}
\newcommand{\popsquiggle}[1]{%
  \tikz[baseline=-0.4ex]{
    \draw[#1,line width=2.6pt,line cap=round]
      plot [smooth] coordinates {(0,0.09) (0.34,0.01) (0.68,0.21) (1.02,0.01) (1.36,0.10)};
  }%
}
% Mot-clé surligné (marqueur orange sous le texte)
\newcommand{\cle}[1]{{\popxbold\color{popdark}#1}}

% ============================================================
% En-tête / pied
% ============================================================
\pagestyle{fancy}
\fancyhf{}
\renewcommand{\headrulewidth}{0pt}
\renewcommand{\footrulewidth}{0pt}
\lhead{\raisebox{-0.15ex}{\poplogo\fontsize{13}{13}\selectfont\color{popink}OnBuch\color{poporange}+}}
\rhead{\poppill{\DOCMATIERE\ \textbullet\ \DOCNIVEAU\ \textbullet\ Leçon \DOCLECON}{white}{popink}}
\lfoot{\popxbold\fontsize{7.6}{7.6}\selectfont\color{popmuted}\MakeUppercase{Cours signé OnBuch+}\ \textbullet\ {\popsemi\DOCTITRE}}
\rfoot{\tikz[baseline=-0.6ex]\node[rounded corners=8.5pt,fill=popink,minimum height=17pt,minimum width=17pt,inner xsep=5pt,inner sep=0]{\popxbold\fontsize{8}{8}\selectfont\color{popcream}\thepage\,/\,\pageref*{LastPage}};}

% ============================================================
% Titres
% ============================================================
\newcommand{\chaptitre}[2]{%
  \begin{tcolorbox}[enhanced,colback=poporange,colframe=popink,boxrule=2.6pt,arc=11pt,
      drop shadow=popink,left=15pt,right=15pt,top=12pt,bottom=11pt,before skip=3pt,after skip=18pt]
    \poppill{#2}{popink}{popcream}\par\vspace{9pt}
    {\raggedright\hyphenpenalty=10000\exhyphenpenalty=10000\popdisplay\fontsize{22}{24}\selectfont\color{popcream}\MakeUppercase{#1}\par}
  \end{tcolorbox}
}
% Section numérotée : pastille orange avec le numéro + titre Archivo
\newcounter{popsec}
\newcommand{\coursec}[1]{%
  \stepcounter{popsec}\setcounter{popsub}{0}%
  \par\vspace{16pt}\noindent
  \begin{minipage}{\linewidth}
  \tikz[baseline=-0.7ex]\node[draw=popink,line width=1.6pt,fill=poporange,rounded corners=4pt,minimum size=22pt,inner sep=0]{\popdisplay\fontsize{12}{12}\selectfont\color{popcream}\thepopsec};%
  \hspace{8pt}{\popdisplay\fontsize{15}{17}\selectfont\color{popink}\MakeUppercase{#1}}\par
  \vspace{4pt}\noindent{\color{poporange}\rule{36pt}{3.6pt}}
  \end{minipage}\par\nopagebreak\vspace{8pt}
}
\newcounter{popsub}
\newcommand{\courssub}[1]{%
  \stepcounter{popsub}%
  \par\vspace{9pt}\noindent{\popxbold\fontsize{11.5}{13}\selectfont\color{popink}{\color{poporange}\thepopsec.\thepopsub}\hspace{6pt}#1}\par\nopagebreak\vspace{4pt}
}

% ============================================================
% Boîtes pédagogiques — même recette : fond blanc, bordure épaisse colorée,
% ombre dure encre, badge plein. Titre optionnel [#1].
% ============================================================
\tcbset{popbase/.style={breakable,enhanced,colback=white,boxrule=2.3pt,arc=9pt,
  shadow={3pt}{-3pt}{0pt}{popink},left=14pt,right=14pt,top=11pt,bottom=11pt,before skip=11pt,after skip=15pt,
  fontupper=\popsemi\fontsize{10.6}{14.5}\selectfont\color{popink},
  fontlower=\popsemi\fontsize{10.6}{14.5}\selectfont\color{popink}}}
% Coche dessinée (le glyphe ✓ n'existe pas dans les polices du document)
\newcommand{\popcheck}[1]{\tikz[baseline=-0.5ex]\draw[#1,line width=1.6pt,line cap=round,line join=round](-0.13,0.0)--(-0.04,-0.09)--(0.13,0.1);}
\providecommand{\coche}{\popcheck{popgreen}}
\providecommand{\resultat}[1]{\par\smallskip\noindent\fcolorbox{popgreen}{popgreenL}{\parbox{\dimexpr\linewidth-2\fboxsep-2\fboxrule\relax}{\textbf{Résultat.} #1}}\par\smallskip}
\newcommand{\popboxhead}[3]{\poppill{#1}{#2}{white}\ifx\relax#3\relax\else\hspace{6pt}{\popxbold\fontsize{10.6}{12}\selectfont\color{popink}#3}\fi\par\vspace{7pt}}

\newtcolorbox{aretenir}[1][]{popbase,colframe=popgreen,before upper={\popboxhead{À retenir}{popgreen}{#1}}}
\newtcolorbox{exemplebox}[1][]{popbase,colframe=poporange,before upper={\popboxhead{Exemple}{poporange}{#1}}}
\newtcolorbox{definition}[1][]{popbase,colframe=popblue,before upper={\popboxhead{Définition}{popblue}{#1}}}
\newtcolorbox{propriete}[1][]{popbase,colframe=poppurple,before upper={\popboxhead{Propriété}{poppurple}{#1}}}
\newtcolorbox{methode}[1][]{popbase,colframe=popgold,before upper={\popboxhead{Méthode}{popgold}{#1}}}
\newtcolorbox{attention}[1][]{popbase,colframe=poppink,before upper={\popboxhead{Attention}{poppink}{#1}}}
\newtcolorbox{experience}[1][]{popbase,colframe=popblue,colback=popblueL,before upper={\popboxhead{Expérience}{popblue}{#1}}}
% « Le savais-tu ? » : carte violette pleine (contexte, culture, Cameroun)
\newtcolorbox{savaistu}[1][]{popbase,colback=poppurple,colframe=popink,
  fontupper=\popsemi\fontsize{10.4}{14.2}\selectfont\color{white},
  before upper={\poppill{Le savais-tu\,?}{white}{poppurple}\ifx\relax#1\relax\else\hspace{6pt}{\popxbold\color{white}#1}\fi\par\vspace{7pt}}}
% Exercice résolu : énoncé en haut, \tcblower puis la solution sur fond crème
\newcounter{popexo}\newcounter{popexr}
\newtcolorbox{exoresolu}[1][]{popbase,colframe=poporange,colbacklower=popcream,
  segmentation style={popink,line width=1.2pt,dashed},
  before upper={\stepcounter{popexr}\popboxhead{Exercice résolu \thepopexr}{poporange}{#1}},
  before lower={\poppill{Solution}{popink}{popcream}\par\vspace{6pt}}}
% Exercice (non résolu en place) — la correction est dans la partie Corrigés
\newtcolorbox{exercice}[1][]{popbase,colframe=popink,
  before upper={\stepcounter{popexo}\popboxhead{Exercice \thepopexo}{popink}{#1}}}
% Corrigé
\newtcolorbox{corrige}[1][]{popbase,colframe=popgreen,colback=popgreenL,
  before upper={\popboxhead{Corrigé}{popgreen}{#1}}}
% Objectifs / prérequis / bilan
\newtcolorbox{objectifs}{popbase,colframe=popink,colback=poporangeL,
  before upper={\popboxhead{Objectifs de la leçon}{popink}{}}}
\newtcolorbox{prerequis}{popbase,colframe=popink,colback=popblueL,
  before upper={\popboxhead{Prérequis}{popblue}{}}}
\newtcolorbox{bilan}{popbase,colframe=popink,colback=white,boxrule=2.8pt,
  before upper={\popboxhead{Fiche bilan}{poporange}{L'essentiel à connaître pour l'examen}}}
% Figure encadrée avec légende
\newtcolorbox{popfigure}[1][]{enhanced,breakable=false,colback=white,colframe=popline,boxrule=1.4pt,arc=8pt,
  left=10pt,right=10pt,top=10pt,bottom=8pt,before skip=10pt,after skip=12pt,halign=center,
  lower separated=false,halign lower=center,
  fontlower=\popsemi\fontsize{9}{11.5}\selectfont\color{popmuted},
  #1}
\newcommand{\legende}[1]{\par\vspace{6pt}{\popsemi\fontsize{9}{11.5}\selectfont\color{popmuted}#1}}

% ============================================================
% Signature OnBuch+
% ============================================================
\newcommand{\poplogoinline}[1]{{\poplogo\fontsize{#1}{#1}\selectfont\color{popink}OnBuch\color{poporange}+}}
\newcommand{\signatureonbuch}{%
  \begin{tcolorbox}[enhanced,colback=popink,colframe=popink,boxrule=2.2pt,arc=10pt,
      drop shadow=poporange,left=14pt,right=14pt,top=10pt,bottom=10pt,before skip=0pt,after skip=0pt]
    \tikz[baseline=-0.7ex]\node[circle,fill=popgreen,draw=popcream,line width=1.3pt,minimum size=22pt,inner sep=0]{\popcheck{white}};%
    \hspace{9pt}\begin{minipage}[c]{0.62\linewidth}
      {\popxbold\fontsize{12}{13}\selectfont\color{popcream}Signé\ \textbullet\ {\poplogo OnBuch\color{poporange}+}}\par\vspace{2pt}
      {\raggedright\popsemi\fontsize{8}{10.5}\selectfont\color{popline}\MakeUppercase{Équipe pédagogique OnBuch+}\\\MakeUppercase{Conforme au programme officiel MINESEC}\par}
    \end{minipage}\hfill
    {\poplogo\fontsize{22}{22}\selectfont\color{popcream}OnBuch\color{poporange}+}
  \end{tcolorbox}%
}

% ============================================================
% Page de garde
% #1 titre · #2 matière · #3 niveau · #4 module · #5 n° leçon · #6 durée ·
% #7 description / objectifs (texte ou itemize)
% ============================================================
\newcommand{\pagegarde}[7]{%
  \thispagestyle{empty}
  \newgeometry{a4paper,margin=1.7cm,top=1.6cm,bottom=1.4cm}%
  \begin{tikzpicture}[remember picture,overlay]
    \fill[poporange!18] ([xshift=-3.2cm,yshift=-3.6cm]current page.north east) circle (4.6cm);
    \fill[poppurple!14] ([xshift=2.4cm,yshift=7.6cm]current page.south west) circle (3.8cm);
  \end{tikzpicture}%
  {\poplogo\fontsize{30}{30}\selectfont\color{popink}OnBuch\color{poporange}+}\hfill
  \raisebox{0.6ex}{\poppill{Cours signé \textbullet\ Premium}{popink}{popcream}}\par
  \vspace{1pt}\popsquiggle{poppurple}\par
  \vspace{8pt}
  \begin{tcolorbox}[enhanced,colback=poporange,colframe=popink,boxrule=2.8pt,arc=13pt,
      drop shadow=popink,left=18pt,right=18pt,top=12pt,bottom=13pt,after skip=10pt]
    \poppill{#2 \textbullet\ #3}{popink}{popcream}\hspace{5pt}\poppill{#4}{white}{popink}\par\vspace{8pt}
    {\popdisplay\fontsize{14}{14}\selectfont\color{popink}LEÇON #5}\par\vspace{4pt}
    {\raggedright\hyphenpenalty=10000\exhyphenpenalty=10000\popdisplay\fontsize{25}{27}\selectfont\color{popcream}\MakeUppercase{#1}\par}
    \vspace{8pt}
    {\popxbold\fontsize{9.5}{9.5}\selectfont\color{popink}\MakeUppercase{Durée conseillée : #6}}
  \end{tcolorbox}
  \begin{tcolorbox}[enhanced,colback=white,colframe=popink,boxrule=2.2pt,arc=10pt,
      drop shadow=popink,left=16pt,right=16pt,top=10pt,bottom=10pt,after skip=8pt]
    \poppill{Dans cette leçon}{poppurple}{white}\par\vspace{5pt}
    \popsemi\fontsize{9.3}{12.6}\selectfont\color{popink}#7
  \end{tcolorbox}
  \noindent\begin{minipage}[t]{0.487\linewidth}
    \begin{tcolorbox}[enhanced,colback=popgreen,colframe=popink,boxrule=2.2pt,arc=10pt,
        drop shadow=popink,left=11pt,right=11pt,top=10pt,bottom=10pt,height=3.1cm,valign=center]
      \tcbox[colback=white,colframe=popink,boxrule=1.3pt,arc=5pt,boxsep=0pt,nobeforeafter,
        left=3pt,right=3pt,top=3pt,bottom=3pt,valign=center]{\qrcode[height=1.9cm]{https://play.google.com/store/apps/details?id=com.luvvix.onbuch&hl=fr}}%
      \hspace{8pt}\begin{minipage}[c]{0.5\linewidth}
        {\popdisplay\fontsize{11}{12.5}\selectfont\color{white}\MakeUppercase{Télécharge OnBuch}}\par\vspace{4pt}
        {\raggedright\popsemi\fontsize{8.4}{11}\selectfont\color{white}Gratuit \textbullet\ Google Play\\Tuteur IA, annales, concours, résultats\par}
      \end{minipage}
    \end{tcolorbox}
  \end{minipage}\hfill
  \begin{minipage}[t]{0.487\linewidth}
    \begin{tcolorbox}[enhanced,colback=poppurple,colframe=popink,boxrule=2.2pt,arc=10pt,
        drop shadow=popink,left=11pt,right=11pt,top=10pt,bottom=10pt,height=3.1cm,valign=center]
      \tikz[baseline=-0.7ex]\node[circle,fill=white,draw=popink,line width=1.3pt,minimum size=1.6cm,inner sep=0]{\popdisplay\fontsize{24}{24}\selectfont\color{poppurple}+};%
      \hspace{8pt}\begin{minipage}[c]{0.6\linewidth}
        {\popdisplay\fontsize{11}{12.5}\selectfont\color{white}\MakeUppercase{Passe à OnBuch+}}\par\vspace{4pt}
        {\raggedright\popsemi\fontsize{8.4}{11}\selectfont\color{white}Tous les cours signés, Tuteur IA illimité, fiches PDF — onglet OnBuch+ de l'app\par}
      \end{minipage}
    \end{tcolorbox}
  \end{minipage}
  \vspace{8pt}
  \popcontenu
  \vfill
  \signatureonbuch
  \restoregeometry
  \newpage
}

% Rangée « Ce cours contient » (page de garde)
\newcommand{\poptile}[3]{% #1 couleur #2 gros texte #3 légende
  \begin{tcolorbox}[enhanced,colback=#1,colframe=popink,boxrule=2pt,arc=9pt,shadow={2.5pt}{-2.5pt}{0pt}{popink},
      left=6pt,right=6pt,top=9pt,bottom=9pt,halign=center,valign=center,height=1.75cm,width=\linewidth,nobeforeafter]
    {\popdisplay\fontsize{12.5}{13}\selectfont\color{white}\MakeUppercase{#2}}\par\vspace{3pt}
    {\centering\popsemi\fontsize{7.8}{9.5}\selectfont\color{white}#3\par}
  \end{tcolorbox}}
\newcommand{\popcontenu}{%
  \par\noindent{\popxboldls\fontsize{7.5}{7.5}\selectfont\color{popmuted}CE COURS SIGNÉ CONTIENT}\par\vspace{7pt}
  \noindent\begin{minipage}[t]{0.235\linewidth}\poptile{popblue}{Cours}{complet, illustré, pas à pas}\end{minipage}\hfill
  \begin{minipage}[t]{0.235\linewidth}\poptile{poporange}{Méthodes}{et exercices résolus}\end{minipage}\hfill
  \begin{minipage}[t]{0.235\linewidth}\poptile{poppink}{Exercices}{type examen + corrigés}\end{minipage}\hfill
  \begin{minipage}[t]{0.235\linewidth}\poptile{popgreen}{Fiche bilan}{l'essentiel à retenir}\end{minipage}\par}

% Sommaire
\newtcolorbox{sommaire}{enhanced,colback=white,colframe=popink,boxrule=2.2pt,arc=10pt,
  drop shadow=popink,left=18pt,right=18pt,top=13pt,bottom=13pt,before skip=6pt,after skip=20pt,
  before upper={\poppill{Sommaire}{poppurple}{white}\par\vspace{9pt}\popsemi\fontsize{10.6}{14.5}\selectfont\color{popink}}}

% Signature de fin de cours — poussée tout en bas de la dernière page
\newcommand{\findecours}{%
  \par\vfill\nopagebreak
  \begin{center}\popsquiggle{poporange}\end{center}\vspace{4pt}
  \signatureonbuch
}
\AtBeginDocument{\def\llbracket{\mathopen{[\mkern-3.5mu[}}\def\rrbracket{\mathclose{]\mkern-3.5mu]}}} % intervalles d'entiers (glyphe ⟦ absent de la police)
% Glyphes absents de Fira Math : redéfinis à partir de glyphes disponibles
\AtBeginDocument{%
  \def\setminus{\mathbin{\backslash}}\let\smallsetminus\setminus
  \def\vdots{\mathord{\vbox{\baselineskip3.2pt\lineskiplimit0pt\kern4pt\hbox{.}\hbox{.}\hbox{.}}}}%
  \def\ddots{\mathinner{\mkern1mu\raise6.5pt\hbox{.}\mkern2mu\raise3.6pt\hbox{.}\mkern2mu\raise0.7pt\hbox{.}\mkern1mu}}%
  \def\triangle{\mathord{\scalebox{0.9}{$\symup{\Delta}$}}}%
  \def\nabla{\mathord{\raisebox{\depth}{\scalebox{1}[-1]{$\symup{\Delta}$}}}}%
  \def\checkmark{\popcheck{popdark}}%
  \def\star{\mathbin{\ast}}\def\bigstar{\mathord{\ast}}%
  \def\diamond{\mathbin{\scalebox{0.6}{$\Diamond$}}}%
  \def\bigcup{\mathop{\mathchoice{\raisebox{-0.3em}{\scalebox{1.7}{$\cup$}}}{\scalebox{1.25}{$\cup$}}{\cup}{\cup}}\displaylimits}%
  \def\bigcap{\mathop{\mathchoice{\raisebox{-0.3em}{\scalebox{1.7}{$\cap$}}}{\scalebox{1.25}{$\cap$}}{\cap}{\cap}}\displaylimits}%
  \def\lesseqqgtr{\mathrel{\lessgtr}}\def\bigoplus{\mathop{\oplus}\displaylimits}%
}
\providecommand{\ssi}{si et seulement si} % abréviation fréquente
\AtBeginDocument{\providecommand{\wideparen}{\overparen}} % arc de cercle

%% FIGFIX-BEGIN v3 — correctifs globaux des figures (docs/onbuch-generation/tools/figfix)
\usepackage{adjustbox}\usepackage{etoolbox}
% toute figure TikZ/pgfplots plus large que la ligne est réduite à la largeur disponible
\BeforeBeginEnvironment{tikzpicture}{\begin{adjustbox}{max width=\linewidth}}
\AfterEndEnvironment{tikzpicture}{\end{adjustbox}}
% légendes pgfplots lisibles par-dessus les courbes
\pgfplotsset{every axis/.append style={legend style={fill=white,fill opacity=0.92,text opacity=1,draw=black!15,inner sep=3pt}}}
% étiquettes d'arêtes (syntaxe edge["..."]) avec fond blanc
\tikzset{every edge quotes/.append style={fill=white,inner sep=1.2pt}}
% bulles de cartes mentales : texte à la ligne dans la bulle, bulle un peu plus grande
\tikzset{every concept/.append style={text width=2.0cm,align=center,minimum size=2.3cm,inner sep=2pt}}
%% FIGFIX-END
\begin{document}

\pagegarde{Limitation de la fréquence de certaines maladies géniques et/ou chromosomiques au sein des familles}{SVT}{Tle D}{Module I : Le Monde Vivant — Prévision en génétique humaine}{N06}{8 h}{Cette leçon étudie les modes de transmission des maladies géniques et chromosomiques au sein des familles humaines à travers l'exploitation des arbres généalogiques. Elle permet d'évaluer les risques génétiques et de justifier l'intérêt des examens prénuptiaux et prénataux dans la prévention des maladies héréditaires.
\vspace{4pt}
\begin{multicols}{2}\small
\begin{itemize}
\item À la fin de cette leçon, je sais définir un arbre généalogique et utiliser les symboles conventionnels du pedigree.
\item À la fin de cette leçon, je sais construire un arbre généalogique à partir d'informations familiales.
\item À la fin de cette leçon, je sais distinguer une hérédité autosomique d'une hérédité gonosomique à partir d'un pedigree.
\item À la fin de cette leçon, je sais déterminer si une transmission est dominante, récessive ou codominante en exploitant un arbre généalogique.
\item À la fin de cette leçon, je sais citer des exemples de maladies autosomiques (albinisme, drépanocytose, achondroplasie) et gonosomiques (daltonisme, hémophilie, rachitisme vitamino-résistant).
\item À la fin de cette leçon, je sais expliquer la transmission des groupes sanguins des systèmes ABO et Rhésus.
\end{itemize}\end{multicols}}

\begin{sommaire}
\begin{multicols}{2}
\begin{enumerate}
\item L'arbre généalogique : outil d'étude de l'hérédité humaine
\item Hérédité autosomique récessive : albinisme et drépanocytose
\item Hérédité autosomique dominante et codominante : achondroplasie et groupes sanguins
\item Hérédité gonosomique : daltonisme, hémophilie, rachitisme vitamino-résistant
\item Évaluation d'un risque génétique
\item Applications : examens prénuptiaux et prénataux
\item Activité d'intégration
\item Méthodes et exercices résolus
\item Exercices
\item Corrigés des exercices
\item Fiche bilan
\end{enumerate}
\end{multicols}
\end{sommaire}

\begin{objectifs}
\begin{itemize}
\item À la fin de cette leçon, je sais définir un arbre généalogique et utiliser les symboles conventionnels du pedigree.
\item À la fin de cette leçon, je sais construire un arbre généalogique à partir d'informations familiales.
\item À la fin de cette leçon, je sais distinguer une hérédité autosomique d'une hérédité gonosomique à partir d'un pedigree.
\item À la fin de cette leçon, je sais déterminer si une transmission est dominante, récessive ou codominante en exploitant un arbre généalogique.
\item À la fin de cette leçon, je sais citer des exemples de maladies autosomiques (albinisme, drépanocytose, achondroplasie) et gonosomiques (daltonisme, hémophilie, rachitisme vitamino-résistant).
\item À la fin de cette leçon, je sais expliquer la transmission des groupes sanguins des systèmes ABO et Rhésus.
\item À la fin de cette leçon, je sais évaluer un risque génétique pour un couple à partir de son pedigree.
\item À la fin de cette leçon, je sais expliquer la nécessité des examens prénuptiaux et prénataux pour prévenir les maladies héréditaires.
\end{itemize}
\end{objectifs}

\begin{prerequis}
\begin{itemize}
\item Notions de gène, d'allèle, de génotype et de phénotype (Première).
\item Caryotype humain : autosomes et gonosomes X et Y (Première).
\item Méiose, fécondation et brassage génétique (Première).
\item Échiquiers de croisement et proportions génétiques (Première).
\item Notion de maladie héréditaire et exemples courants observés au Cameroun (drépanocytose, albinisme).
\end{itemize}
\end{prerequis}

\newpage

\coursec{L'arbre généalogique : outil d'étude de l'hérédité humaine}

Dans un quartier de Yaoundé, un couple dont les deux familles comptent des cas de drépanocytose hésite à se marier : leurs parents exigent un examen prénuptial avant les fiançailles. Ailleurs, une mère daltonienne et un père à la vision normale s'interrogent : leurs fils seront-ils daltoniens ? Et pourquoi l'albinisme apparaît-il parfois chez des enfants dont les deux parents sont normalement pigmentés ? Ces questions, fréquentes dans les familles camerounaises, trouvent leurs réponses dans l'analyse des \cle{arbres généalogiques} et les lois de transmission des caractères héréditaires.

\textbf{Problématique :} comment peut-on étudier scientifiquement la transmission des caractères héréditaires chez l'Homme, alors qu'il est impossible de réaliser des croisements expérimentaux comme chez les plantes ou la drosophile ?

\courssub{Pourquoi une méthode d'étude spécifique chez l'Homme ?}

Pour établir les lois de la génétique, Mendel a croisé des lignées pures de pois, et Morgan a croisé des drosophiles porteuses de caractères différents. La démarche expérimentale classique repose donc sur des \cle{croisements dirigés} : on choisit les parents, on contrôle les fécondations, on analyse une descendance nombreuse sur plusieurs générations.

Chez l'Homme, cette démarche est \textbf{impossible} pour deux raisons majeures :

\begin{itemize}
\item \textbf{Une raison éthique :} on ne peut pas imposer à des êtres humains des unions expérimentales dans le but d'étudier la transmission d'un caractère. La liberté individuelle et la dignité humaine l'interdisent absolument.
\item \textbf{Une raison biologique :} la durée d'une génération humaine est longue (environ 20 à 30 ans), le nombre de descendants par couple est faible (quelques enfants, contre des centaines de graines chez le pois ou des centaines de mouches chez la drosophile), et le chercheur ne peut pas observer plus de 3 ou 4 générations au cours de sa propre vie.
\end{itemize}

Il a donc fallu imaginer une méthode d'étude \textbf{spécifique à l'espèce humaine} : au lieu de provoquer les croisements, le généticien \emph{observe} les unions qui se sont produites naturellement au sein des familles, recense les individus qui présentent le caractère étudié (par exemple une maladie héréditaire), et reconstitue l'histoire familiale de ce caractère. C'est la \cle{méthode généalogique}, dont l'outil central est l'\cle{arbre généalogique}.

\begin{aretenir}[Savoir admis]
L'étude de l'hérédité humaine ne peut pas reposer sur des croisements expérimentaux (raisons éthique et longueur des générations). Elle s'appuie sur l'\textbf{observation des familles} et l'analyse des \textbf{arbres généalogiques} (ou pedigrees).
\end{aretenir}

\courssub{Définition de l'arbre généalogique (pedigree)}

\begin{definition}[Arbre généalogique ou pedigree]
Un \cle{arbre généalogique} (ou \cle{pedigree}) est la \textbf{représentation schématique, à l'aide de symboles conventionnels, de la transmission d'un caractère au sein d'une famille sur plusieurs générations}. Il indique, pour chaque individu, son sexe, sa position dans la famille et la présence ou l'absence du caractère étudié (par exemple : atteint ou non par une maladie héréditaire).
\end{definition}

L'intuition est simple : un pedigree est à la génétique humaine ce que le tableau de croisement est à la génétique expérimentale. En regardant \emph{qui} est atteint, \emph{à quelle génération}, et \emph{dans quel sexe}, on peut déduire le mode de transmission du caractère : autosomique ou gonosomique, dominant ou récessif. C'est exactement ce que nous ferons dans les sections suivantes avec l'albinisme, la drépanocytose, le daltonisme ou l'hémophilie.

\courssub{Les symboles conventionnels du pedigree}

Pour que tout généticien du monde puisse lire un pedigree, des symboles \textbf{conventionnels} ont été adoptés. Ils sont exigibles au Baccalauréat : tu dois savoir les utiliser et les interpréter sans hésitation.

\begin{itemize}
\item un \textbf{carré} représente un \textbf{homme} ; un \textbf{rond} (cercle) représente une \textbf{femme} ;
\item un symbole \textbf{plein} (noirci) représente un individu \textbf{atteint} (qui présente le caractère étudié, par exemple la maladie) ; un symbole \textbf{vide} représente un individu \textbf{sain} (qui ne présente pas le caractère) ;
\item un \textbf{trait horizontal} reliant un carré et un rond représente une \textbf{union} (couple) ;
\item un \textbf{trait vertical} descendant du trait d'union représente la \textbf{filiation} : il mène aux enfants du couple, reliés entre eux par un trait horizontal de fratrie ;
\item les \textbf{générations} sont numérotées en chiffres romains de haut en bas : I, II, III\ldots{} ; les \textbf{individus} sont numérotés en chiffres arabes de gauche à droite au sein de chaque génération : I-1, I-2, II-1, II-2\ldots{}
\end{itemize}

\begin{popfigure}
\begin{center}
\begin{tikzpicture}[scale=1, every node/.style={font=\small}]
% Homme sain
\draw[thick, popblue] (0,0) rectangle (0.5,0.5);
\node[right] at (0.7,0.25) {Homme sain};
% Femme saine
\draw[thick, poppink] (4.2,0) circle (0.28);
\node[right] at (4.6,0.25) {Femme saine};
% Homme atteint
\filldraw[thick, popblue, fill=popblue!70] (8.6,0) rectangle (9.1,0.5);
\node[right] at (9.3,0.25) {Homme atteint};
% Femme atteinte
\filldraw[thick, poppink, fill=poppink!70] (12.6,0) circle (0.28);
\node[right] at (13.0,0.25) {Femme atteinte};
% Union
\draw[thick, popblue] (0, -1.6) rectangle (0.5,-1.1);
\draw[thick] (0.5,-1.35) -- (1.7,-1.35);
\draw[thick, poppink] (1.98,-1.35) circle (0.28);
\node[right] at (2.5,-1.35) {Trait horizontal : union (couple)};
% Filiation
\draw[thick, popblue] (8.6,-2.4) rectangle (9.1,-1.9);
\draw[thick] (9.1,-2.15) -- (10.3,-2.15);
\draw[thick, poppink] (10.58,-2.15) circle (0.28);
\draw[thick] (9.7,-2.15) -- (9.7,-2.9);
\draw[thick] (8.9,-2.9) -- (10.5,-2.9);
\draw[thick] (8.9,-2.9) -- (8.9,-3.2);
\draw[thick] (10.5,-2.9) -- (10.5,-3.2);
\draw[thick, popblue] (8.65,-3.5) rectangle (9.15,-3.2);
\draw[thick, poppink] (10.5,-3.35) circle (0.15);
\node[right] at (11.0,-2.9) {Trait vertical : filiation (descendance)};
\end{tikzpicture}
\end{center}
\legende{Symboles conventionnels de l'arbre généalogique. Carré = homme ; rond = femme ; symbole plein = individu atteint ; symbole vide = individu sain ; trait horizontal = union ; trait vertical = filiation.}
\end{popfigure}

\begin{attention}[Erreur fréquente : union ou filiation ?]
Ne confonds jamais le \textbf{trait d'union} (horizontal, il relie deux conjoints, un carré et un rond) et le \textbf{trait de filiation} (vertical, il descend du trait d'union vers les enfants). Au Baccalauréat, un pedigree mal construit — par exemple des enfants reliés directement à un seul parent, ou un trait horizontal entre un père et son fils — est sanctionné. Autre piège : noircir un symbole alors que l'énoncé dit que l'individu est sain. Lis le texte descriptif deux fois avant de tracer.
\end{attention}

\courssub{Lecture d'un pedigree : un exemple guidé}

Lisons ensemble l'arbre généalogique ci-dessous, qui représente la transmission d'un caractère (par exemple une maladie héréditaire) sur trois générations d'une famille.

\begin{popfigure}
\begin{center}
\begin{tikzpicture}[scale=1.05, every node/.style={font=\small}]
% Génération I
\draw[thick, popblue] (3,3) rectangle (3.5,3.5);
\node[above] at (3.25,3.5) {I-1};
\draw[thick, poppink] (4.5,3.25) circle (0.28);
\node[above] at (4.5,3.55) {I-2};
\draw[thick] (3.5,3.25) -- (4.22,3.25);
% Descendance vers génération II
\draw[thick] (3.86,3.25) -- (3.86,2.4);
\draw[thick] (1.6,2.4) -- (6.1,2.4);
\draw[thick] (1.6,2.4) -- (1.6,2.0);
\draw[thick] (3.86,2.4) -- (3.86,2.0);
\draw[thick] (6.1,2.4) -- (6.1,2.0);
% Génération II : II-1 (femme saine), II-2 (homme atteint), II-3 (femme saine) + conjoint de II-3
\draw[thick, poppink] (1.6,1.72) circle (0.28);
\node[above] at (1.6,2.0) {II-1};
\filldraw[thick, popblue, fill=popblue!70] (3.61,1.45) rectangle (4.11,1.95);
\node[above] at (3.86,1.95) {II-2};
\draw[thick, poppink] (6.1,1.72) circle (0.28);
\node[above] at (6.1,2.0) {II-3};
% Conjoint de II-3
\draw[thick, popblue] (7.4,1.45) rectangle (7.9,1.95);
\node[above] at (7.65,1.95) {II-4};
\draw[thick] (6.38,1.72) -- (7.4,1.72);
% Descendance vers génération III
\draw[thick] (6.89,1.72) -- (6.89,0.9);
\draw[thick] (5.9,0.9) -- (7.9,0.9);
\draw[thick] (5.9,0.9) -- (5.9,0.5);
\draw[thick] (6.9,0.9) -- (6.9,0.5);
\draw[thick] (7.9,0.9) -- (7.9,0.5);
% Génération III
\filldraw[thick, poppink, fill=poppink!70] (5.9,0.22) circle (0.28);
\node[below] at (5.9,-0.08) {III-1};
\draw[thick, popblue] (6.65,-0.06) rectangle (7.15,0.44);
\node[below] at (6.9,-0.08) {III-2};
\filldraw[thick, popblue, fill=popblue!70] (7.65,-0.06) rectangle (8.15,0.44);
\node[below] at (7.9,-0.08) {III-3};
% Étiquettes générations
\node[left, font=\bfseries] at (0.4,3.25) {I};
\node[left, font=\bfseries] at (0.4,1.72) {II};
\node[left, font=\bfseries] at (0.4,0.2) {III};
\end{tikzpicture}
\end{center}
\legende{Exemple d'arbre généalogique sur trois générations. Le couple fondateur I-1 $\times$ I-2 (tous deux sains) a trois enfants : II-1 (fille saine), II-2 (garçon atteint), II-3 (fille saine). II-3, unie à II-4 (homme sain), a trois enfants : III-1 (fille atteinte), III-2 (garçon sain), III-3 (garçon atteint).}
\end{popfigure}

La lecture d'un pedigree suit toujours le même ordre logique :

\begin{enumerate}
\item \textbf{Identifier les générations} : ici trois générations (I, II, III), numérotées de haut en bas.
\item \textbf{Repérer les individus atteints} : II-2, III-1 et III-3 ont les symboles noircis ; tous les autres sont sains.
\item \textbf{Décrire les liens de parenté} : I-1 et I-2 sont les grands-parents ; II-1, II-2, II-3 sont frère et sœurs (une fratrie) ; III-1, III-2, III-3 sont les enfants de II-3 et II-4, donc les petits-enfants de I-1 et I-2.
\item \textbf{Dégager les faits générateurs d'hypothèses} : ici, des parents sains (I-1 et I-2) ont un enfant atteint (II-2) ; le caractère touche les deux sexes. Ces observations, exploitées dans les sections suivantes, suggèrent déjà une transmission \emph{autosomique récessive} : le caractère peut « sauter » une génération et être porté silencieusement par des individus sains.
\end{enumerate}

\begin{methode}[Construire un pedigree à partir d'un texte descriptif]
\begin{enumerate}
\item \textbf{Repérer les générations} dans le texte (grands-parents, parents, enfants) et placer la génération la plus ancienne (I) en haut de la feuille.
\item \textbf{Placer les couples} de la génération I : carré à gauche, rond à droite, reliés par un trait horizontal.
\item \textbf{Tracer le trait de filiation} vertical depuis le milieu du trait d'union, puis un trait horizontal de fratrie, et placer les enfants (génération II) de gauche à droite, de l'aîné au cadet si l'information est donnée.
\item \textbf{Ajouter les conjoints} éventuels des individus de la génération II, puis leurs enfants (génération III) selon le même procédé.
\item \textbf{Noircir} les symboles des individus atteints, laisser vides ceux des individus sains.
\item \textbf{Numéroter} les générations (I, II, III) et chaque individu (I-1, I-2, II-1\ldots), puis vérifier la cohérence avec le texte.
\end{enumerate}
\end{methode}

\begin{exoresolu}[Construction d'un pedigree à partir d'un texte]
\textbf{Énoncé.} Dans une famille de Bafoussam, les grands-parents paternels, tous deux normalement pigmentés, ont eu quatre enfants : trois filles normalement pigmentées et un garçon albinos. Ce garçon, uni à une femme normalement pigmentée, a eu deux enfants : une fille albinos et un garçon normalement pigmenté. Construis l'arbre généalogique de cette famille.
\tcblower
\textbf{Solution détaillée.}
\begin{enumerate}
\item \textbf{Génération I} : un couple, tous deux sains $\Rightarrow$ carré vide (I-1) et rond vide (I-2) reliés par un trait horizontal.
\item \textbf{Génération II} : quatre enfants. On place de gauche à droite les trois filles saines (ronds vides : II-1, II-2, II-3) puis le garçon albinos (carré plein : II-4). On ajoute la conjointe de II-4 : rond vide (II-5), reliée à II-4 par un trait horizontal.
\item \textbf{Génération III} : les enfants de II-4 et II-5 : une fille albinos (rond plein : III-1) et un garçon sain (carré vide : III-2), suspendus au trait de filiation issu du couple II-4 $\times$ II-5.
\item \textbf{Vérification} : 2 générations de couples, 9 individus, 2 atteints (II-4, III-1) — le texte et le schéma concordent.
\end{enumerate}
Observation clé pour la suite de la leçon : des parents sains (I-1, I-2) ont un enfant atteint (II-4) $\Rightarrow$ l'allèle responsable de l'albinisme est \emph{récessif} et porté à l'état hétérozygote par les grands-parents ; le caractère touche les deux sexes $\Rightarrow$ il est vraisemblablement \emph{autosomique}. Nous démontrerons cela rigoureusement dans la section 2.
\end{exoresolu}

\courssub{Limites de la méthode généalogique}

La méthode des pedigrees est puissante, mais elle a des \textbf{limites} qu'il faut connaître (savoir admis) :

\begin{itemize}
\item \textbf{Les familles humaines sont petites} : avec deux ou trois enfants par couple, les proportions théoriques (par exemple 1/4 d'enfants atteints pour un caractère récessif) ne se vérifient pas toujours dans une fratrie donnée. Un couple à risque peut n'avoir que des enfants sains par simple effet du hasard, ou au contraire plusieurs enfants atteints.
\item \textbf{Les informations sont parfois incomplètes ou incertaines} : dossiers médicaux anciens absents, individus décédés jeunes, diagnostics approximatifs, adoptions ou paternités non déclarées, générations lointaines mal connues.
\item \textbf{Un caractère peut avoir plusieurs causes} : deux familles présentant des pedigrees semblables peuvent être atteintes de maladies génétiquement différentes (hétérogénéité génétique), ce qui complique l'interprétation.
\end{itemize}

C'est pourquoi le généticien croise toujours l'analyse des pedigrees avec des examens biologiques modernes (analyse de l'hémoglobine pour la drépanocytose, tests de coagulation pour l'hémophilie, étude chromosomique ou de l'ADN), comme nous le verrons avec les examens prénuptiaux et prénataux.

\begin{savaistu}[L'hémophilie des familles royales d'Europe]
L'analyse des pedigrees a permis de retracer la transmission de l'\textbf{hémophilie} (maladie du sang qui ne coagule pas) dans les familles royales d'Europe. La reine \textbf{Victoria} d'Angleterre (1819--1901) était porteuse saine de l'allèle responsable, situé sur le chromosome X. Par les mariages de ses filles et petites-filles avec les cours d'Espagne, de Russie et d'Allemagne, la maladie s'est répandue : le tsarévitch Alexis de Russie, fils du tsar Nicolas II, en était atteint. Ce cas historique célèbre illustre la transmission \emph{gonosomique récessive} : les femmes transmettent l'allèle sans être malades, les hommes qui le reçoivent sont atteints. Nous l'étudierons en détail dans la section 4.
\end{savaistu}

\begin{aretenir}[L'essentiel de la section]
\begin{itemize}
\item L'étude de l'hérédité humaine ne peut pas utiliser de croisements expérimentaux (éthique, longueur des générations, faible descendance) : on utilise la \textbf{méthode généalogique}.
\item Un \textbf{pedigree} est la représentation schématique de la transmission d'un caractère sur plusieurs générations, à l'aide de symboles conventionnels : carré (homme), rond (femme), plein (atteint), vide (sain), trait horizontal (union), trait vertical (filiation).
\item La lecture d'un pedigree (générations, individus atteints, liens de parenté) permet de formuler des hypothèses sur le \textbf{mode de transmission} du caractère : autosomique ou gonosomique, dominant ou récessif.
\item La méthode a des limites : petites familles, informations incomplètes ; elle est complétée par les examens biologiques.
\end{itemize}
\end{aretenir}

\begin{exercice}[Pour t'entraîner]
Dans une famille de Douala, un couple sain a trois enfants : une fille saine, un garçon atteint de daltonisme et une fille saine. La fille aînée, mariée à un homme sain, a deux fils dont l'un est daltonien. (1) Construis le pedigree de cette famille en respectant les conventions. (2) Cite deux observations du pedigree qui te permettront, dans la section 4, de proposer un mode de transmission gonosomique récessif.
\end{exercice}

\begin{corrige}[Exercice 1]
\textbf{(1) Construction du pedigree :}
\begin{itemize}
\item \textbf{Génération I} : Couple sain $\rightarrow$ Carré vide (I-1) et Rond vide (I-2) unis par trait horizontal.
\item \textbf{Génération II} : Trois enfants (de gauche à droite) :
      \begin{itemize}
      \item Fille aînée saine $\rightarrow$ Rond vide (II-1).
      \item Garçon daltonien $\rightarrow$ Carré \textbf{plein} (II-2).
      \item Deuxième fille saine $\rightarrow$ Rond vide (II-3).
      \end{itemize}
      Conjoint de II-1 : Homme sain $\rightarrow$ Carré vide (II-4), uni à II-1 par trait horizontal.
\item \textbf{Génération III} : Enfants de II-1 $\times$ II-4 (deux fils) :
      \begin{itemize}
      \item Fils aîné daltonien $\rightarrow$ Carré \textbf{plein} (III-1).
      \item Fils cadet sain $\rightarrow$ Carré vide (III-2).
      \end{itemize}
      (Suspendus au trait de filiation descendant du couple II-1/II-4).
\end{itemize}
\textbf{Numérotation :} I-1, I-2 ; II-1, II-2, II-3, II-4 ; III-1, III-2.

\textbf{(2) Observations orientant vers une transmission gonosomique récessive (liée au X) :}
\begin{enumerate}
\item \textbf{Atteinte prédominante des hommes :} Dans cette famille, seuls des hommes sont atteints (II-2, III-1). Les femmes (I-2, II-1, II-3) sont phénotypiquement saines. Cela suggère que l'allèle muté est récessif et porté par le chromosome X : un homme (XY) n'ayant qu'un seul X exprime la maladie s'il porte l'allèle, alors qu'une femme (XX) doit être homozygote pour l'exprimer (ce qui est plus rare).
\item \textbf{Transmission par les femmes (absence de transmission père-fils) :} Le père (I-1) est sain et a un fils atteint (II-2). Le fils atteint (II-2) tient donc son allèle muté de sa mère (I-2), qui est une femme saine $\rightarrow$ elle est \textbf{porteuse saine} (hétérozygote). De même, la fille aînée (II-1), née de parents dont la mère est porteuse, a 1/2 chance de l'être ; elle transmet l'allèle à l'un de ses deux fils (III-1 atteint, III-2 sain). L'absence de transmission directe père $\rightarrow$ fils (I-1 sain $\rightarrow$ II-2 atteint ; II-4 sain $\rightarrow$ III-1 atteint) confirme que le gène n'est pas sur le chromosome Y ni sur un autosome (sinon le père devrait être atteint ou le caractère toucherait les deux sexes équitablement).
\end{enumerate}
Ces deux arguments (excès d'hommes atteints, transmission par des femmes saines, pas de transmission père-fils) sont la signature classique de l'hérédité \textbf{gonosomique récessive liée au chromosome X}.
\end{corrige}

\coursec{Hérédité autosomique récessive : albinisme et drépanocytose}

Dans la section précédente, tu as appris à construire et à lire un \cle{arbre généalogique} : cercles pour les femmes, carrés pour les hommes, symboles remplis pour les individus atteints, traits horizontaux pour les unions et verticaux pour les descendances. Cet outil n'a d'intérêt que s'il permet de \emph{raisonner}. C'est précisément ce que nous allons faire maintenant : à partir d'un pedigree réel, nous allons découvrir, pas à pas, les deux questions fondamentales que le généticien se pose toujours : l'allèle responsable de la maladie est-il \cle{dominant} ou \cle{récessif} ? Le gène est-il porté par un \cle{autosome} ou par un \cle{gonosome} (chromosome sexuel) ? L'albinisme et la drépanocytose, deux affections fréquentes au Cameroun, nous serviront de fil conducteur.

\courssub{Étude de cas : un pedigree qui pose question}

\begin{experience}[Observation d'une famille camerounaise]
Dans une famille suivie au Centre Hospitalier d'Essos (Yaoundé), on recense les faits suivants : le couple I (un homme et une femme, tous deux à la peau normalement pigmentée) a quatre enfants. Parmi eux, une fille et un garçon sont \textbf{albinos} : peau très claire, cheveux blancs, iris gris-bleu, forte sensibilité à la lumière. Les deux autres enfants ont une pigmentation normale. Aucun des deux parents n'est albinos, et aucun des grands-parents ne l'était.
\begin{center}
\begin{tikzpicture}[scale=0.95, every node/.style={font=\small}]
% Génération I
\node[draw, rectangle, minimum size=0.7cm, thick, popblue] (pere) at (0,0) {};
\node[draw, circle, minimum size=0.7cm, thick, popblue] (mere) at (2,0) {};
\draw[thick] (pere.east) -- (mere.west);
\node[below=2pt of pere] {I-1 (sain)};
\node[below=2pt of mere] {I-2 (saine)};
% Ligne de descendance
\coordinate (mid) at (1,0);
\draw[thick] (mid) -- ++(0,-0.6) coordinate (sib);
\draw[thick] (-1.5,-0.6) -- (3.5,-0.6);
% Génération II
\node[draw, rectangle, minimum size=0.7cm, thick, popblue] (f1) at (-1.5,-1.4) {};
\node[draw, circle, minimum size=0.7cm, thick, popblue, fill=poporange] (f2) at (0.2,-1.4) {};
\node[draw, rectangle, minimum size=0.7cm, thick, popblue, fill=poporange] (f3) at (1.9,-1.4) {};
\node[draw, circle, minimum size=0.7cm, thick, popblue] (f4) at (3.5,-1.4) {};
\draw[thick] (-1.5,-0.6) -- (f1.north);
\draw[thick] (0.2,-0.6) -- (f2.north);
\draw[thick] (1.9,-0.6) -- (f3.north);
\draw[thick] (3.5,-0.6) -- (f4.north);
\node[below=2pt of f1] {II-1 (sain)};
\node[below=2pt of f2] {II-2 (albinos)};
\node[below=2pt of f3] {II-3 (albinos)};
\node[below=2pt of f4] {II-4 (saine)};
\end{tikzpicture}
\end{center}
\textbf{Fait à interpréter :} deux parents phénotypiquement sains donnent naissance à des enfants atteints, et la maladie frappe aussi bien un garçon qu'une fille.
\end{experience}

Face à ce pedigree, la démarche d'investigation scientifique consiste à formuler des hypothèses puis à les tester.

\textbf{Hypothèse 1 : l'allèle responsable est dominant.} Si tel était le cas, tout individu portant cet allèle serait malade. Or les parents I-1 et I-2 sont sains : ils ne porteraient donc pas l'allèle, et ne pourraient pas le transmettre. L'hypothèse est \emph{contredite} par l'existence des enfants II-2 et II-3.

\textbf{Hypothèse 2 : l'allèle responsable est récessif.} L'allèle morbide, noté $a$, ne s'exprimerait qu'à l'état homozygote $a/a$. Les parents, sains mais parents d'enfants atteints, posséderaient chacun un allèle $a$ « caché » par l'allèle normal $A$ : ils seraient de génotype $A/a$, c'est-à-dire \cle{hétérozygotes}. Chacun peut transmettre $a$ ; l'enfant qui reçoit $a$ du père \emph{et} $a$ de la mère est $a/a$, donc albinos. L'hypothèse rend parfaitement compte de toutes les observations.

\textbf{Localisation du gène.} La maladie touche indifféremment le garçon (II-3) et la fille (II-2). Si le gène était porté par le chromosome sexuel $X$, on observerait une répartition inégale entre les sexes (nous le verrons en section 4). La transmission identique chez les deux sexes indique que le gène est porté par un \cle{autosome}, c'est-à-dire un chromosome non sexuel.

\begin{definition}[Maladie autosomique récessive]
Une maladie génétique est dite \cle{autosomique récessive} lorsque le gène responsable est porté par un autosome et que l'allèle morbide ne s'exprime qu'à l'état \textbf{homozygote} ($a/a$). Les individus hétérozygotes ($A/a$) sont phénotypiquement sains : on les appelle \cle{porteurs sains} (ou conducteurs), car ils peuvent transmettre l'allèle morbide à leur descendance sans en manifester les symptômes.
\end{definition}

\begin{propriete}[Signature d'une transmission autosomique récessive sur un pedigree]
Deux indices, lisibles directement sur l'arbre généalogique, permettent de conclure :
\begin{itemize}
\item des \textbf{enfants atteints naissent de parents sains} : le caractère peut donc « sauter » une ou plusieurs générations $\Rightarrow$ l'allèle est \textbf{récessif} ;
\item la maladie touche \textbf{aussi bien les garçons que les filles} $\Rightarrow$ le gène est porté par un \textbf{autosome}.
\end{itemize}
Cette propriété, dégagée de l'exploitation de nombreux pedigrees, est un outil de diagnostic génétique exigible au Baccalauréat.
\end{propriete}

\begin{methode}[Conclure à partir d'un pedigree]
\begin{enumerate}
\item Repérer s'il existe des enfants atteints nés de parents sains. Si oui : allèle \textbf{récessif} et les deux parents sont obligatoirement \textbf{hétérozygotes}.
\item Comparer la répartition des atteints entre garçons et filles. Répartition équilibrée : gène \textbf{autosomique}.
\item Attribuer les génotypes : atteint $= a/a$ ; parents sains d'un enfant atteint $= A/a$.
\item Vérifier la cohérence en écrivant le croisement et son échiquier.
\end{enumerate}
\end{methode}

\courssub{L'albinisme : une enzyme en panne, un allèle récessif}

L'\cle{albinisme} est une maladie génétique caractérisée par l'\textbf{absence de mélanine}, le pigment brun-noir qui colore la peau, les poils, les cheveux et l'iris, et qui protège la peau des rayons ultraviolets du soleil. La mélanine est fabriquée par des cellules spécialisées, les mélanocytes, grâce à une enzyme appelée \emph{tyrosinase}. L'allèle normal $A$ commande la synthèse d'une tyrosinase fonctionnelle ; l'allèle muté $a$ commande une enzyme inactive. Chez un hétérozygote $A/a$, l'allèle normal suffit à produire assez d'enzyme : la pigmentation est normale, ce qui explique pourquoi $A$ est dominant et $a$ récessif. Seul l'individu $a/a$, dépourvu de toute tyrosinase active, est albinos.

\begin{center}
\begin{tabularx}{\linewidth}{|c|X|X|}
\hline
\rowcolor{popblueL} \textbf{Génotype} & \textbf{Phénotype} & \textbf{Explication biochimique} \\
\hline
$A/A$ & sain (pigmentation normale) & deux allèles fonctionnels, tyrosinase active \\
\hline
$A/a$ & sain (porteur sain) & un allèle fonctionnel suffit à produire la mélanine \\
\hline
$a/a$ & albinos & aucune tyrosinase active, pas de mélanine \\
\hline
\end{tabularx}
\end{center}

Prévoyons maintenant la descendance du couple I-1 $\times$ I-2, tous deux hétérozygotes $A/a$. Chaque parent produit deux types de gamètes en proportions égales : $A$ ($1/2$) et $a$ ($1/2$). L'échiquier de croisement donne :

\begin{popfigure}
\begin{center}
\begin{tikzpicture}[scale=1.05]
% Grille de l'échiquier
\draw[thick, popdark] (0,0) grid (3,3);
% Gamètes
\node at (1.5,3.45) {\textbf{Gamètes du père} $A/a$};
\node[rotate=90] at (-0.75,1.5) {\textbf{Gamètes de la mère} $A/a$};
\node at (1.5,2.5) {\large $A$};
\node at (2.5,2.5) {\large $a$};
\node at (0.5,1.5) {\large $A$};
\node at (0.5,0.5) {\large $a$};
% Cases
\node[fill=popgreenL, minimum width=0.95cm, minimum height=0.95cm] at (1.5,1.5) {$A/A$};
\node[fill=popgreenL, minimum width=0.95cm, minimum height=0.95cm] at (2.5,1.5) {$A/a$};
\node[fill=popgreenL, minimum width=0.95cm, minimum height=0.95cm] at (1.5,0.5) {$A/a$};
\node[fill=poporangeL, minimum width=0.95cm, minimum height=0.95cm] at (2.5,0.5) {$a/a$};
% Proportions
\node[below] at (1.5,-0.25) {\small $1/4$ sain \quad $2/4$ porteurs sains \quad $1/4$ albinos};
\end{tikzpicture}
\end{center}
\legende{Échiquier de croisement $A/a \times A/a$. Cases vertes : descendants sains ($A/A$ ou $A/a$) ; case orange : descendant albinos ($a/a$). Proportions génotypiques : $1/4\ A/A$, $2/4\ A/a$, $1/4\ a/a$ ; proportions phénotypiques : $3/4$ sains, $1/4$ albinos.}
\end{popfigure}

À chaque naissance, ce couple a donc une probabilité de $1/4$, soit \SI{25}{\%}, d'avoir un enfant albinos, quel que soit le sexe de l'enfant. C'est exactement ce que montre le pedigree étudié : sur quatre enfants, deux sont atteints — les proportions théoriques sont des \emph{probabilités par naissance}, pas des quotas obligatoirement respectés dans une petite fratrie.

\begin{attention}[Pièges fréquents]
\begin{itemize}
\item Un porteur sain $A/a$ n'est \textbf{pas} malade : ne jamais écrire qu'il est « un peu albinos ». L'allèle récessif ne s'exprime pas du tout chez lui.
\item Le risque de $1/4$ est le même \textbf{à chaque grossesse} : un couple qui a déjà eu un enfant albinos n'est pas « protégé » pour les naissances suivantes.
\item Ne pas confondre « récessif » et « rare » : un allèle récessif peut être fréquent dans une population, porté silencieusement par de nombreux hétérozygotes.
\end{itemize}
\end{attention}

\courssub{La drépanocytose : une hémoglobine anormale}

La \cle{drépanocytose} (ou anémie falciforme) est une maladie du sang très fréquente en zone tropicale, notamment au Cameroun. Elle résulte d'une mutation du gène qui commande la synthèse de l'hémoglobine, la protéine des globules rouges chargée du transport du dioxygène. L'allèle normal $A$ (ou $Hb^A$) commande l'hémoglobine normale $HbA$ ; l'allèle muté $S$ (ou $Hb^S$) commande une hémoglobine anormale $HbS$ qui, en milieu pauvre en dioxygène, se solidifie en longs filaments et déforme le globule rouge en \textbf{faucille} (d'où « falciforme »). Ces globules déformés se bloquent dans les petits vaisseaux (crises douloureuses) et sont détruits prématurément (anémie chronique).

\begin{experience}[Électrophorèse de l'hémoglobine : le fait expérimental clé]
On dépose l'hémoglobine extraite du sang de trois individus ($A/A$, $A/S$, $S/S$) sur un gel, puis on soumet le gel à un champ électrique. Les protéines migrent d'autant plus vite qu'elles sont plus chargées ; $HbA$ et $HbS$ n'ayant pas la même charge électrique, elles se séparent en bandes distinctes. On observe :
\begin{itemize}
\item individu $A/A$ : une \textbf{seule bande}, au niveau $HbA$ ;
\item individu $S/S$ : une \textbf{seule bande}, au niveau $HbS$ (migré différemment) ;
\item individu $A/S$ : \textbf{deux bandes}, l'une au niveau $HbA$, l'autre au niveau $HbS$.
\end{itemize}
\textbf{Interprétation :} l'hétérozygote fabrique \emph{à la fois} l'hémoglobine normale et l'hémoglobine anormale. Aucun des deux allèles ne masque l'autre : ils s'expriment tous les deux. C'est le fait expérimental qui fonde la notion de \cle{codominance}.
\end{experience}

\begin{popfigure}
\begin{center}
\begin{tikzpicture}[scale=1.0]
% Puits et bandes pour 3 profils
\foreach \x/\nom/\type in {0/{A/A}/AA, 3.2/{A/S}/AS, 6.4/{S/S}/SS}{
  \draw[thick, popdark, fill=popcream] (\x,0) rectangle (\x+1.6,4.2);
  \node[below] at (\x+0.8,-0.1) {\nom};
  % puits
  \draw[fill=popmuted!40] (\x+0.5,3.7) rectangle (\x+1.1,4.0);
}
% Bande HbA (haut, y=2.9)
\draw[fill=popblue] (0.5,2.8) rectangle (1.1,3.0);
\draw[fill=popblue] (3.7,2.8) rectangle (4.3,3.0);
% Bande HbS (bas, y=1.9)
\draw[fill=poporange] (3.7,1.8) rectangle (4.3,2.0);
\draw[fill=poporange] (6.9,1.8) rectangle (7.5,2.0);
% Annotations
\node[right, popblue] at (8.2,2.9) {\small bande $HbA$};
\node[right, poporange] at (8.2,1.9) {\small bande $HbS$};
\draw[->, thick, popdark] (-0.6,4.1) -- (-0.6,0.2) node[midway, left, rotate=90]{\small sens de migration};
\node[above] at (0.8,4.2) {\small dépôt};
\end{tikzpicture}
\end{center}
\legende{Profils d'électrophorèse de l'hémoglobine. Le sujet $A/A$ ne présente que la bande $HbA$ ; le sujet $S/S$ (drépanocytaire) ne présente que la bande $HbS$ ; le sujet $A/S$ (porteur sain) présente les \textbf{deux} bandes : les deux allèles s'expriment simultanément, ce qui traduit la codominance.}
\end{popfigure}

\begin{definition}[Codominance]
Deux allèles sont dits \cle{codominants} lorsqu'ils s'expriment \textbf{tous les deux} chez l'individu hétérozygote, qui présente alors un phénotype intermédiaire ou composite. Pour la drépanocytose, l'hétérozygote $A/S$ synthétise $HbA$ \emph{et} $HbS$ : une partie de ses globules rouges peut se déformer en faucille lors d'efforts intenses ou en altitude. On parle de \textbf{drépanocytose mineure} (ou trait drépanocytaire) ; le sujet $A/S$ est cliniquement quasi sain, c'est un porteur sain.
\end{definition}

\begin{exemplebox}[Couple $A/S \times A/S$ : le calcul du risque]
Deux porteurs sains se marient. Leurs gamètes sont $A$ ($1/2$) et $S$ ($1/2$). L'échiquier de croisement (identique en structure à celui de l'albinisme) donne :
\begin{itemize}
\item $1/4$ d'enfants $A/A$ : non porteurs, hémoglobine normale ;
\item $2/4$ d'enfants $A/S$ : porteurs sains (drépanocytose mineure) ;
\item $1/4$ d'enfants $S/S$ : drépanocytaires majeurs, atteints d'anémie falciforme grave.
\end{itemize}
Le risque d'avoir un enfant malade est donc de \SI{25}{\%} à chaque naissance.
\end{exemplebox}

\begin{savaistu}[Drépanocytose et paludisme au Cameroun]
Au Cameroun, environ \SI{20}{\%} à \SI{25}{\%} de la population est porteuse de l'allèle $S$. Pourquoi une fréquence si élevée pour un allèle grave ? Parce que l'hétérozygote $A/S$ possède un avantage sélectif en zone d'endémie palustre : ses globules rouges, partiellement falciformes, offrent un milieu défavorable au développement du \emph{Plasmodium}, parasite du paludisme. Les porteurs sains résistent mieux au paludisme que les sujets $A/A$ : l'allèle $S$ a ainsi été maintenu par la sélection naturelle dans les régions tropicales. C'est un exemple classique d'\textbf{avantage de l'hétérozygote}, souvent cité au Baccalauréat. Conséquence pratique : deux jeunes gens choisis au hasard au Cameroun ont une probabilité non négligeable d'être tous deux $A/S$, d'où l'importance capitale du \textbf{dépistage prénuptial} (test d'électrophorèse de l'hémoglobine avant le mariage).
\end{savaistu}

\begin{exoresolu}[Exploitation d'un pedigree de drépanocytose]
Énoncé. Dans une famille de Bafoussam, M. et Mme N., tous deux en bonne santé apparente, ont trois enfants : une fille atteinte d'anémie falciforme grave et deux garçons sains. L'électrophorèse montre que M. et Mme N. possèdent chacun les deux types d'hémoglobine. 1) Déterminer le mode de transmission de la maladie. 2) Donner les génotypes de chaque membre de la famille. 3) Calculer le risque que le prochain enfant du couple soit malade.
\tcblower
Solution détaillée.
1) Des parents sains ont un enfant atteint : l'allèle morbide $S$ ne s'exprime pas chez eux, il est donc \textbf{récessif} au sens clinique (la maladie majeure n'apparaît qu'à l'état $S/S$). La maladie touche une fille dans une fratrie comportant des garçons sains, et le père comme la mère transmettent : la transmission est identique dans les deux sexes, le gène est \textbf{autosomique}. (Au niveau moléculaire, l'électrophorèse révèle une codominance entre $A$ et $S$ : les deux notions ne se contredisent pas, elles décrivent deux niveaux d'observation.)
2) L'électrophorèse montre deux bandes chez chaque parent : M. N. et Mme N. sont $A/S$ (porteurs sains). La fille malade est $S/S$. Chaque garçon sain est soit $A/A$, soit $A/S$ ; seule une électrophorèse permettrait de trancher (probabilités : $1/3\ A/A$, $2/3\ A/S$ parmi les enfants sains).
3) Croisement $A/S \times A/S$ : l'échiquier donne $1/4\ A/A$, $2/4\ A/S$, $1/4\ S/S$. Le risque d'un enfant malade ($S/S$) à la prochaine naissance est de $1/4$, soit \SI{25}{\%}, indépendamment des naissances précédentes.
\end{exoresolu}

\begin{aretenir}[L'essentiel de la section]
\begin{itemize}
\item Une maladie \textbf{autosomique récessive} ne s'exprime qu'à l'état homozygote $a/a$ ; les hétérozygotes $A/a$ sont des \textbf{porteurs sains}.
\item Signature sur un pedigree : enfants atteints nés de parents sains (allèle récessif) et atteinte des deux sexes (gène autosomique).
\item Albinisme : absence de mélanine ; croisement $A/a \times A/a \rightarrow 3/4$ sains, $1/4$ albinos.
\item Drépanocytose : hémoglobine $HbS$ anormale ; l'électrophorèse révèle la \textbf{codominance} des allèles $A$ et $S$ (deux bandes chez $A/S$) ; couple $A/S \times A/S \rightarrow$ \SI{25}{\%} d'enfants $S/S$ malades.
\item La fréquence élevée de l'allèle $S$ au Cameroun (\SIrange{20}{25}{\%} de porteurs) justifie le dépistage prénuptial par électrophorèse.
\end{itemize}
\end{aretenir}

\coursec{Hérédité autosomique dominante et codominante : achondroplasie et groupes sanguins}

Dans la section précédente, nous avons vu que l'albinisme et la drépanocytose se transmettent selon un mode \cle{autosomique récessif} : l'allèle muté ne s'exprime qu'à l'état homozygote, la maladie peut \og sauter \fg{} des générations, et deux parents sains mais conducteurs peuvent donner naissance à un enfant malade. Mais toutes les maladies génétiques ne se comportent pas ainsi. Certaines apparaissent à \emph{chaque génération}, sans discontinuité, comme si l'allèle responsable s'imposait dès qu'il est présent. C'est le cas de l'\cle{achondroplasie}, forme la plus fréquente de nanisme. Par ailleurs, il existe des situations où \emph{deux} allèles s'expriment simultanément, sans qu'aucun ne domine l'autre : c'est la \cle{codominance}, illustrée magistralement par les groupes sanguins du système ABO. C'est tout l'objet de cette section.

\courssub{L'achondroplasie : une maladie autosomique dominante}

\textbf{Observation et document d'étude : un pedigree d'achondroplasie}

L'achondroplasie est une maladie du développement osseux : la croissance en longueur des os des membres est fortement ralentie, ce qui donne des adultes de petite taille (environ \SI{1,30}{m}) aux membres courts mais au tronc et à la tête de dimensions quasi normales. Elle touche environ une naissance sur \num{25000} dans le monde, y compris au Cameroun. Étudions le pedigree d'une famille atteinte.

\begin{popfigure}
\begin{tikzpicture}[
  male/.style={draw, thick, popblue, minimum size=0.6cm, inner sep=0pt},
  female/.style={draw, thick, poppink, circle, minimum size=0.6cm, inner sep=0pt},
  maleA/.style={draw, thick, fill=popblue!70, minimum size=0.6cm, inner sep=0pt},
  femaleA/.style={draw, thick, fill=poppink!70, circle, minimum size=0.6cm, inner sep=0pt},
  union/.style={thick},
  gen/.style={font=\small\bfseries, popmuted},
  node distance=0.8cm and 1.2cm
]
% Génération I
\node[gen] at (-2.0,0.3) {I};
\node[maleA] (I1) at (0,0) {};
\node[female] (I2) at (1.6,0) {};
\draw[union] (I1) -- (I2);
\coordinate (u1) at (0.8,0);
% Génération II
\node[gen] at (-2.0,-2.2) {II};
\node[maleA] (II1) at (-1.0,-2.2) {};
\node[female] (II2) at (0.6,-2.2) {};
\node[male] (II3) at (2.2,-2.2) {};
\draw[union] (u1) -- (0.8,-0.7) -- (-1.0,-0.7) -- (II1);
\draw[union] (0.8,-0.7) -- (2.2,-0.7) -- (II3);
\draw[union] (II1) -- (II2);
\coordinate (u2) at (-0.2,-2.2);
% Génération III
\node[gen] at (-2.0,-4.4) {III};
\node[femaleA] (III1) at (-1.6,-4.4) {};
\node[male] (III2) at (-0.2,-4.4) {};
\node[femaleA] (III3) at (1.2,-4.4) {};
\draw[union] (u2) -- (-0.2,-2.9) -- (-1.6,-2.9) -- (III1);
\draw[union] (-0.2,-2.9) -- (III2);
\draw[union] (-0.2,-2.9) -- (1.2,-2.9) -- (III3);
% Numéros
\node[below=0.1cm of I1, font=\scriptsize] {1};
\node[below=0.1cm of I2, font=\scriptsize] {2};
\node[below=0.1cm of II1, font=\scriptsize] {1};
\node[below=0.1cm of II2, font=\scriptsize] {2};
\node[below=0.1cm of II3, font=\scriptsize] {3};
\node[below=0.1cm of III1, font=\scriptsize] {1};
\node[below=0.1cm of III2, font=\scriptsize] {2};
\node[below=0.1cm of III3, font=\scriptsize] {3};
% Légende des symboles
\node[maleA, label=right:{\scriptsize homme atteint}] at (3.5,0.3) {};
\node[femaleA, label=right:{\scriptsize femme atteinte}] at (3.5,-0.7) {};
\node[male, label=right:{\scriptsize homme sain}] at (3.5,-1.7) {};
\node[female, label=right:{\scriptsize femme saine}] at (3.5,-2.7) {};
\end{tikzpicture}
\legende{Pedigree d'une famille atteinte d'achondroplasie sur trois générations. Carré = homme ; cercle = femme ; symbole rempli = individu atteint ; trait horizontal = union ; trait vertical = descendance. On observe que la maladie est présente à chaque génération et que chaque enfant atteint (II1, III1, III3) a au moins un parent atteint.}
\end{popfigure}

\textbf{Analyse du document.} Trois faits frappent immédiatement :
\begin{itemize}
  \item la maladie est présente à \emph{chaque génération} : I1, puis II1, puis III1 et III3 ;
  \item tout individu atteint a \emph{au moins un parent atteint} : II1 a pour père I1 ; III1 et III3 ont pour père II1 ;
  \item les deux sexes sont touchés de façon comparable, et la transmission se fait aussi bien par le père que par la mère : le gène n'est donc pas porté par un chromosome sexuel, il est \cle{autosomique}.
\end{itemize}

Comparons avec ce que nous avons vu pour l'albinisme : là, des parents sains pouvaient avoir un enfant albinos, et la maladie sautait des générations. Ici, c'est exactement l'inverse. Ce contraste est la signature d'une transmission \emph{dominante}.

\begin{definition}[Maladie autosomique dominante]
Une maladie héréditaire est dite \cle{autosomique dominante} lorsque l'allèle muté, porté par un autosome, s'exprime dès la présence d'\emph{un seul} exemplaire : l'individu \cle{hétérozygote} $A/a$ est atteint, tout comme l'homozygote $A/A$. Il suffit donc d'hériter de l'allèle muté de l'un des deux parents pour être malade.
\end{definition}

\textbf{Interprétation génétique du pedigree.} Notons $A$ l'allèle muté responsable de l'achondroplasie et $a$ l'allèle normal. Les individus sains sont nécessairement $a/a$. L'individu I1, atteint, a eu des enfants sains (II3) avec une femme saine $a/a$ : il a donc transmis l'allèle $a$ à II3, ce qui prouve qu'il est \emph{hétérozygote} $A/a$. Le croisement I1 $\times$ I2 s'écrit :

\[ A/a \times a/a \longrightarrow \tfrac{1}{2}\, A/a \text{ (atteints)} \;+\; \tfrac{1}{2}\, a/a \text{ (sains)} \]

Chaque enfant d'un parent atteint hétérozygote a donc une probabilité de $1/2$ (soit 50 \%) d'être atteint, quel que soit son sexe. C'est exactement ce que montre le pedigree (environ un enfant atteint sur deux à chaque union).

\begin{aretenir}[Un savoir admis : l'homozygote $A/A$ n'est pas viable]
Dans l'achondroplasie, l'allèle muté dominant s'exprime dès l'état \cle{hétérozygote} : les malades sont donc presque toujours $A/a$. Les individus \cle{homozygotes dominants} $A/A$ présentent des malformations osseuses si graves qu'ils sont généralement \emph{non viables} (mort \emph{in utero} ou peu après la naissance). C'est pourquoi, en pratique, tous les achondroplases rencontrés sont hétérozygotes.
\end{aretenir}

\begin{propriete}[Conséquence fondamentale de la dominance]
Dans une transmission dominante, \textbf{deux parents sains ne peuvent pas avoir d'enfant atteint} : étant $a/a \times a/a$, ils ne transmettent que l'allèle normal $a$. (Seule une \emph{mutation nécessaire}, c'est-à-dire une mutation nouvelle survenue dans les gamètes, peut contredire cette règle — cas hors programme.) À l'inverse, tout enfant atteint d'une maladie dominante a au moins un parent atteint.
\end{propriete}

\begin{methode}[Distinguer transmission dominante et récessive sur un pedigree]
Face à un arbre généalogique, procède en trois étapes :
\begin{enumerate}
  \item \textbf{Cherche des parents sains ayant un enfant atteint.} Si un tel couple existe, l'allèle muté était caché chez les parents : la transmission est \textbf{récessive}.
  \item \textbf{Vérifie si tout enfant atteint a un parent atteint.} Si oui, et si la maladie apparaît à chaque génération, la transmission est \textbf{dominante}.
  \item \textbf{Précise la localisation du gène.} Si pères et mères transmettent indifféremment et si les deux sexes sont touchés de façon comparable, le gène est \textbf{autosomique}. (La question du gonosomique sera traitée dans la section suivante.)
\end{enumerate}
Résumé mnémotechnique : \og saut de génération $\Rightarrow$ récessif ; présence à chaque génération $\Rightarrow$ dominant \fg{}.
\end{methode}

\begin{attention}[Piège de rédaction]
Ne dis jamais qu'une maladie dominante \og ne saute jamais de génération \fg{} de façon absolue : un parent atteint peut n'avoir, par hasard, que des enfants sains. La formulation rigoureuse est : \emph{tout individu atteint possède au moins un parent atteint}, et \emph{deux parents sains n'ont jamais d'enfant atteint}. C'est ce dernier critère, certain, qu'il faut invoquer au Baccalauréat.
\end{attention}

\courssub{La codominance : le système de groupes sanguins ABO}

\textbf{Observation.} Lors d'une transfusion sanguine — par exemple au Centre Hospitalier d'Essos à Yaoundé —, le personnel soignant doit impérativement connaître le groupe sanguin du donneur et du receveur : A, B, AB ou O. Ce caractère, déterminé par des molécules (antigènes) présentes à la surface des globules rouges, est héréditaire. Mais son étude révèle une situation nouvelle : il n'existe pas deux, mais \emph{trois} allèles pour un même gène.

\begin{definition}[Allèles multiples et codominance]
Le gène du système \cle{ABO}, porté par la paire d'autosomes n° 9, existe sous \emph{trois} formes alléliques : $I^{A}$, $I^{B}$ et $i$.
\begin{itemize}
  \item $I^{A}$ et $I^{B}$ sont \cle{codominants} entre eux : chez un individu $I^{A}/I^{B}$, \emph{les deux} allèles s'expriment simultanément, et les deux antigènes A et B sont présents sur les globules rouges (groupe AB).
  \item $I^{A}$ et $I^{B}$ sont tous deux \emph{dominants} sur l'allèle $i$, qui est \emph{récessif} et ne code aucun antigène.
\end{itemize}
\end{definition}

Chaque individu ne portant que \emph{deux} allèles (un sur chaque chromosome homologue), il existe six génotypes possibles, pour seulement quatre phénotypes :

\begin{center}
\begin{tabularx}{\linewidth}{|c|c|X|}
\hline
\rowcolor{popblueL} \textbf{Phénotype (groupe)} & \textbf{Génotype(s)} & \textbf{Antigènes sur les hématies} \\
\hline
A & $I^{A}/I^{A}$ ou $I^{A}/i$ & antigène A \\
\hline
B & $I^{B}/I^{B}$ ou $I^{B}/i$ & antigène B \\
\hline
AB & $I^{A}/I^{B}$ & antigènes A \emph{et} B (codominance) \\
\hline
O & $i/i$ & aucun antigène \\
\hline
\end{tabularx}
\end{center}

\begin{attention}[Erreur fréquente à bannir]
Ne dis jamais \og A est dominant sur B \fg{} ! Dans le système ABO, $I^{A}$ et $I^{B}$ sont \textbf{codominants entre eux} (ils s'expriment tous les deux chez $I^{A}/I^{B}$), mais chacun est \textbf{dominant sur $i$}. La relation de dominance concerne $I^{A}$ (ou $I^{B}$) \emph{face à $i$}, jamais $I^{A}$ face à $I^{B}$.
\end{attention}

\courssub{Le système Rhésus : un cas simple de dominance}

Le second système de groupes sanguins, découvert chez le singe \emph{Macacus rhesus} (d'où son nom), est plus simple : il n'existe que deux allèles.
\begin{itemize}
  \item L'allèle $Rh^{+}$ commande la présence de l'antigène Rhésus (facteur D) sur les hématies ; il est \emph{dominant}.
  \item L'allèle $Rh^{-}$ ne code aucun antigène ; il est \emph{récessif}.
\end{itemize}

\begin{aretenir}[Génotypes du système Rhésus]
\begin{itemize}
  \item Individu de phénotype \textbf{Rhésus positif} : génotype $Rh^{+}/Rh^{+}$ \emph{ou} $Rh^{+}/Rh^{-}$ (l'allèle dominant suffit) ;
  \item Individu de phénotype \textbf{Rhésus négatif} : génotype obligatoirement $Rh^{-}/Rh^{-}$.
\end{itemize}
Un individu Rhésus positif peut donc être homozygote ou hétérozygote : seul l'examen de sa descendance (ou de ses parents) permet parfois de trancher. La notation complète d'un groupe sanguin combine les deux systèmes, indépendants : par exemple A Rh$^{+}$, O Rh$^{-}$, AB Rh$^{-}$.
\end{aretenir}

\courssub{Application : prévision des groupes sanguins des enfants d'un couple}

Considérons un couple dont le père est de groupe A et la mère de groupe B. Quels groupes sanguins leurs enfants peuvent-ils avoir ? Tout dépend des génotypes. Si le père est $I^{A}/i$ et la mère $I^{B}/i$ (tous deux hétérozygotes), l'échiquier de croisement donne :

\begin{popfigure}
\begin{tikzpicture}
\matrix (pun) [
  matrix of nodes,
  nodes={draw, thick, minimum size=1.1cm, anchor=center, font=\small, inner sep=2pt},
  column sep=-\pgflinewidth,
  row sep=-\pgflinewidth,
  row 1/.style={nodes={fill=popblueL, font=\small\bfseries, text=popdark}},
  column 1/.style={nodes={fill=poppinkL, font=\small\bfseries, text=popdark}},
  row 2/.style={nodes={fill=popcream}},
  row 3/.style={nodes={fill=popcream}}
] {
& $I^{A}$ & $i$ \\
$I^{B}$ & $I^{A}/I^{B}$\newline\textbf{AB} & $I^{B}/i$\newline\textbf{B} \\
$i$ & $I^{A}/i$\newline\textbf{A} & $i/i$\newline\textbf{O} \\
};
% Étiquettes extérieures
\node[above=0.2cm of pun-1-2, font=\small\bfseries, popblue] {Gamètes du père ($I^{A}/i$)};
\node[left=0.3cm of pun-2-1, font=\small\bfseries, poppink, rotate=90] {Gamètes de la mère ($I^{B}/i$)};
\end{tikzpicture}
\legende{Échiquier de croisement d'un couple $I^{A}/i \times I^{B}/i$ (père de groupe A, mère de groupe B, tous deux hétérozygotes). Chaque case a une probabilité de $1/4$. Les quatre groupes sanguins A, B, AB et O sont possibles dans la descendance.}
\end{popfigure}

\begin{exemplebox}[Résultat chiffré du croisement $I^{A}/i \times I^{B}/i$]
Les quatre cases de l'échiquier étant équiprobables, les enfants de ce couple ont :
\begin{itemize}
  \item 25 \% de chances d'être de groupe \textbf{AB} ($I^{A}/I^{B}$) ;
  \item 25 \% de chances d'être de groupe \textbf{A} ($I^{A}/i$) ;
  \item 25 \% de chances d'être de groupe \textbf{B} ($I^{B}/i$) ;
  \item 25 \% de chances d'être de groupe \textbf{O} ($i/i$).
\end{itemize}
Ainsi, un couple A $\times$ B peut avoir des enfants de \emph{tous} les groupes — y compris O, alors qu'aucun des deux parents n'est O ! Ce résultat, surprenant pour le sens commun, s'explique parfaitement par la récessivité de $i$ et la ségrégation des allèles à la méiose.
\end{exemplebox}

\begin{exoresolu}[Prévision génétique dans une famille]
\textbf{Énoncé.} Monsieur N., de groupe A Rh$^{+}$, et Madame N., de groupe B Rh$^{+}$, ont un premier enfant de groupe O Rh$^{-}$. (a) Déterminer les génotypes complets des deux parents. (b) Quelle est la probabilité que leur deuxième enfant soit de groupe AB Rh$^{+}$ ?
\tcblower
\textbf{Solution détaillée.}
(a) \emph{Génotypes des parents.} L'enfant O est $i/i$ : il a reçu un allèle $i$ de chaque parent. Le père, de groupe A, est donc $I^{A}/i$ ; la mère, de groupe B, est donc $I^{B}/i$. De même, l'enfant Rh$^{-}$ est $Rh^{-}/Rh^{-}$ : chaque parent, bien que Rh$^{+}$, porte un allèle $Rh^{-}$ ; ils sont donc tous deux $Rh^{+}/Rh^{-}$.
Génotypes complets : Père $[I^{A}/i ; Rh^{+}/Rh^{-}]$, Mère $[I^{B}/i ; Rh^{+}/Rh^{-}]$.

(b) \emph{Probabilité recherchée.} Les deux systèmes ABO et Rhésus sont portés par des chromosomes différents (chr 9 et chr 1) : ils se transmettent indépendamment.
\begin{itemize}
  \item Probabilité du groupe AB : d'après l'échiquier $I^{A}/i \times I^{B}/i$, $P(AB) = \dfrac{1}{4}$.
  \item Probabilité du phénotype Rh$^{+}$ : le croisement $Rh^{+}/Rh^{-} \times Rh^{+}/Rh^{-}$ donne $\dfrac{1}{4}\,Rh^{+}/Rh^{+} + \dfrac{2}{4}\,Rh^{+}/Rh^{-} + \dfrac{1}{4}\,Rh^{-}/Rh^{-}$, soit $P(Rh^{+}) = \dfrac{3}{4}$.
\end{itemize}
Par indépendance des deux caractères :
\[ P(AB\;Rh^{+}) = \frac{1}{4} \times \frac{3}{4} = \frac{3}{16} = 0{,}1875 \approx 18{,}75\ \%. \]
Le deuxième enfant a donc environ 18,75 \% de chances d'être AB Rh$^{+}$.
\end{exoresolu}

\begin{savaistu}[Groupes sanguins et transfusion au Cameroun]
Les individus de groupe O Rh$^{-}$ sont dits \og donneurs universels \fg{} (leurs hématies ne portent aucun antigène A, B ou Rhésus), et les AB Rh$^{+}$ \og receveurs universels \fg{}. Dans les banques de sang des hôpitaux camerounais, la pénurie de sang O Rh$^{-}$ est un enjeu de santé publique : connaître son groupe et donner son sang sauve des vies, notamment lors des hémorragies post-partum et des anémies sévères dues au paludisme.
\end{savaistu}

\begin{exercice}[Pour t'entraîner]
Une femme de groupe A, dont le père était de groupe O, épouse un homme de groupe AB. (1) Donner les génotypes des deux époux. (2) Établir l'échiquier de croisement et déterminer les groupes sanguins possibles de leurs enfants, avec leurs probabilités. (3) Ce couple peut-il avoir un enfant de groupe O ? Justifier.
\end{exercice}

\begin{corrige}[Exercice : couple A $\times$ AB]
(1) La femme, de groupe A, a un père $i/i$ dont elle a obligatoirement reçu un allèle $i$ : elle est $I^{A}/i$. L'homme AB est nécessairement $I^{A}/I^{B}$.

(2) Échiquier $I^{A}/i \times I^{A}/I^{B}$ :

\begin{center}
\begin{tabular}{|c|c|c|}
\hline
\rowcolor{popblueL} \textbf{Gamètes} & $I^{A}$ (mère) & $i$ (mère) \\
\hline
$I^{A}$ (père) & $I^{A}/I^{A}$ : A & $I^{A}/i$ : A \\
\hline
$I^{B}$ (père) & $I^{A}/I^{B}$ : AB & $I^{B}/i$ : B \\
\hline
\end{tabular}
\end{center}

Probabilités : 50 \% A, 25 \% AB, 25 \% B.

(3) Non, ce couple ne peut pas avoir d'enfant O : un enfant $i/i$ devrait recevoir un allèle $i$ de chaque parent, or le père $I^{A}/I^{B}$ ne possède aucun allèle $i$ à transmettre.
\end{corrige}

\begin{aretenir}[L'essentiel de la section]
\begin{itemize}
  \item Une maladie \textbf{autosomique dominante} (achondroplasie) s'exprime dès l'état hétérozygote ; elle apparaît à chaque génération, et tout enfant atteint a un parent atteint. Deux parents sains n'ont jamais d'enfant atteint.
  \item Le système \textbf{ABO} met en jeu trois allèles : $I^{A}$ et $I^{B}$ codominants entre eux, tous deux dominants sur $i$ récessif.
  \item Le système \textbf{Rhésus} : $Rh^{+}$ dominant, $Rh^{-}$ récessif ; seul $Rh^{-}/Rh^{-}$ est de phénotype négatif.
  \item L'échiquier de croisement permet de prévoir les groupes sanguins possibles d'une descendance : un couple $I^{A}/i \times I^{B}/i$ peut engendrer les quatre groupes à parts égales (25 \% chacun).
\end{itemize}
\end{aretenir}

\coursec{Hérédité gonosomique : daltonisme, hémophilie, rachitisme vitamino-résistant}

Dans la section précédente, nous avons étudié des caractères dont les gènes sont portés par les \cle{autosomes} : l'achondroplasie, transmise de façon dominante, et les groupes sanguins des systèmes ABO et Rhésus, qui illustrent la codominance. Dans tous ces cas, le sexe de l'individu ne joue aucun rôle : garçons et filles sont atteints avec la même fréquence. Or l'observation de certaines familles montre un fait troublant : \textbf{certaines maladies ne frappent presque que les garçons}, alors que les mères, parfaitement saines, transmettent la maladie. Comment expliquer cette inégalité devant la maladie ? C'est toute la question de l'\cle{hérédité gonosomique}, c'est-à-dire liée aux chromosomes sexuels.

\courssub{Les gonosomes X et Y : un support inégal entre les sexes}

Rappelons un fait établi par l'étude des \cle{caryotypes} humains : les cellules d'une femme possèdent 22 paires d'autosomes et deux chromosomes X identiques (formule chromosomique $2n = 44\,\text{AA} + XX$), tandis que celles d'un homme possèdent 22 paires d'autosomes, un chromosome X et un chromosome Y beaucoup plus petit ($2n = 44\,\text{AA} + XY$). Le chromosome X est grand et porte de nombreux gènes ; le chromosome Y, très réduit, ne porte que peu de gènes, essentiellement impliqués dans la détermination du sexe masculin.

Conséquence capitale : pour un gène porté par le chromosome X, la femme possède \textbf{deux exemplaires} (deux allèles), comme pour n'importe quel gène autosomique, alors que l'homme n'en possède \textbf{qu'un seul}. On dit que l'homme est \cle{hémizygote} pour les gènes portés par X : il ne peut être ni homozygote ni hétérozygote, et tout allèle présent sur son unique X s'exprime obligatoirement, qu'il soit dominant ou récessif.

\begin{definition}[Hérédité gonosomique]
Une hérédité est dite \cle{gonosomique} (ou liée au sexe) lorsque le gène responsable du caractère étudié est porté par un chromosome sexuel (gonosome). Dans les cas étudiés au programme, le gène est porté par le \textbf{chromosome X} ; le chromosome Y ne porte pas d'allèle correspondant. On note les génotypes en indiquant l'allèle en exposant du X : par exemple $X^{D}$, $X^{d}$, et $Y$ sans exposant.
\end{definition}

\begin{attention}[Ce que le caryotype ne dit pas]
Le caryotype montre les chromosomes sexuels, mais il ne permet pas, à lui seul, d'affirmer qu'un gène donné est porté par X : il faut pour cela des informations complémentaires (pedigrees, croisements statistiques). Au Baccalauréat, la localisation sur X est un \textbf{savoir admis} pour le daltonisme, l'hémophilie et le rachitisme vitamino-résistant : tu dois l'utiliser pour interpréter les arbres généalogiques, non la redémontrer à partir des seuls caryotypes.
\end{attention}

\courssub{Étude de cas : le daltonisme, une maladie de garçons transmise par des mères saines}

\textbf{Observation.} Le \cle{daltonisme} est une anomalie de la vision des couleurs (confusion typique du rouge et du vert), due à un défaut des pigments des cônes rétiniens. Elle n'empêche pas de vivre normalement, mais elle interdit certains métiers (pilote, électricien). Considérons le pedigree suivant, relevé dans une famille : les grands-parents I-1 et I-2 sont sains ; parmi leurs enfants, une fille (II-2), saine, épouse un homme sain (II-4) ; parmi leurs quatre enfants, \textbf{deux garçons sont daltoniens}, les deux filles sont saines.

\begin{popfigure}
\begin{center}
\begin{tikzpicture}[
  male/.style={draw=popblue, thick, minimum size=0.55cm, inner sep=0pt},
  female/.style={draw=poppink, thick, circle, minimum size=0.55cm, inner sep=0pt},
  maleaff/.style={draw=popblue, thick, fill=popblue!70, minimum size=0.55cm, inner sep=0pt},
  femalecond/.style={draw=poppink, thick, circle, fill=poppink!25, minimum size=0.55cm, inner sep=0pt},
  link/.style={thick, popdark},
  every label/.style={font=\small\itshape, popmuted}
]
% Génération I
\node[male, label=above:{I-1 : $X^{D}Y$}] (I1) at (0,0) {};
\node[female, label=above:{I-2 : $X^{D}X^{d}$}] (I2) at (1.4,0) {};
\draw[link] (I1) -- (I2);
\draw[link] (0.7,0) -- (0.7,-0.5);
\draw[link] (-0.9,-0.5) -- (3.1,-0.5);
% Génération II : enfants de I
\node[male, label=below:{II-1 : $X^{D}Y$}] (II1) at (-0.9,-1.3) {};
\node[femalecond, label=below:{II-2 : $X^{D}X^{d}$}] (II2) at (0.7,-1.3) {};
\node[male, label=below:{II-3 : $X^{D}Y$}] (II3) at (3.1,-1.3) {};
\draw[link] (0.7,-0.5) -- (II2);
\draw[link] (-0.9,-0.5) -- (II1);
\draw[link] (3.1,-0.5) -- (II3);
% Conjoint de II-2
\node[male, label=below:{II-4 : $X^{D}Y$}] (II4) at (2.1,-1.3) {};
\draw[link] (II2) -- (II4);
% Génération III (enfants de II-2 x II-4)
\draw[link] (1.4,-1.3) -- (1.4,-1.9);
\draw[link] (0.2,-1.9) -- (2.6,-1.9);
\node[maleaff, label=below:{III-1 : $X^{d}Y$}] (III1) at (0.2,-2.7) {};
\node[female, label=below:{III-2 : $X^{D}X^{?}$}] (III2) at (1.0,-2.7) {};
\node[maleaff, label=below:{III-3 : $X^{d}Y$}] (III3) at (1.8,-2.7) {};
\node[female, label=below:{III-4 : $X^{D}X^{?}$}] (III4) at (2.6,-2.7) {};
\draw[link] (0.2,-1.9) -- (III1);
\draw[link] (1.0,-1.9) -- (III2);
\draw[link] (1.8,-1.9) -- (III3);
\draw[link] (2.6,-1.9) -- (III4);
\end{tikzpicture}
\end{center}
\legende{Pedigree d'une famille atteinte de daltonisme. Carré = homme, cercle = femme ; symbole rempli = individu daltonien ; cercle à fond clair = femme conductrice saine. $X^{D}X^{?}$ : fille saine dont le second allèle est indéterminé ($X^{D}X^{D}$ ou $X^{D}X^{d}$). Seuls des garçons sont atteints, et la maladie est transmise par des mères saines.}
\end{popfigure}

\textbf{Fait à interpréter.} Deux observations sautent aux yeux : (1) seuls les garçons sont atteints ; (2) les parents des garçons malades sont tous deux sains. Appliquons la démarche d'investigation.

\begin{itemize}
  \item \emph{Étape 1 : dominance ou récessivité.} Des parents sains (II-2 et II-4) ont des enfants malades : chacun porte donc l'allèle morbide sans l'exprimer. L'allèle responsable, noté $d$, est \textbf{récessif} ; l'allèle normal $D$ est dominant.
  \item \emph{Étape 2 : localisation du gène.} Si le gène était autosomique, les garçons et les filles seraient atteints avec la même fréquence. Or seuls les garçons sont malades. De plus, le père II-4, sain, a des fils malades : si le gène est sur X, le père transmet son Y à ses fils et ne peut donc pas leur donner l'allèle morbide — c'est la mère qui transmet. Tout concorde : le gène est \textbf{porté par le chromosome X}.
  \item \emph{Conclusion.} La mère II-2 est de génotype $X^{D}X^{d}$ : elle est saine car l'allèle normal $D$ domine, mais elle peut transmettre $X^{d}$ à ses enfants. On l'appelle \cle{femme conductrice} (ou porteuse saine). Un fils qui reçoit $X^{d}$ de sa mère et $Y$ de son père est $X^{d}Y$ : n'ayant qu'un seul X, il est obligatoirement daltonien.
\end{itemize}

\begin{definition}[Femme conductrice]
Une \cle{femme conductrice} (ou porteuse saine) est une femme hétérozygote $X^{D}X^{d}$ : elle possède un allèle morbide récessif sur l'un de ses chromosomes X, masqué par l'allèle normal dominant porté par l'autre X. Elle est phénotypiquement saine mais peut transmettre l'allèle morbide à sa descendance.
\end{definition}

\textbf{Les cinq génotypes possibles} pour le daltonisme se résument ainsi :

\begin{center}
\begin{tabularx}{\linewidth}{|l|X|X|}
\hline
\rowcolor{popblueL} \textbf{Génotype} & \textbf{Sexe} & \textbf{Phénotype} \\
\hline
$X^{D}X^{D}$ & Femme & Saine (non conductrice) \\
\hline
$X^{D}X^{d}$ & Femme & Saine conductrice \\
\hline
$X^{d}X^{d}$ & Femme & Daltonienne (très rare) \\
\hline
$X^{D}Y$ & Homme & Sain \\
\hline
$X^{d}Y$ & Homme & Daltonien \\
\hline
\end{tabularx}
\end{center}

Une femme daltonienne $X^{d}X^{d}$ ne peut naître que d'un père daltonien \emph{et} d'une mère au moins conductrice : cette combinaison rare explique pourquoi environ 8 \% des hommes sont daltoniens contre moins de 0,5 \% des femmes.

\begin{propriete}[Pourquoi les hommes sont-ils plus souvent atteints ?]
Dans une transmission \textbf{récessive liée à X}, les hommes sont beaucoup plus souvent atteints que les femmes, car ils ne possèdent qu'\textbf{un seul chromosome X} : un seul allèle récessif suffit à rendre un homme malade, alors qu'une femme doit en recevoir deux (un de chaque parent). Il n'existe pas d'homme « conducteur sain » : un homme est soit sain ($X^{D}Y$), soit malade ($X^{d}Y$).
\end{propriete}

\courssub{L'échiquier de croisement : prévoir la descendance d'une mère conductrice}

Croisons une mère conductrice $X^{D}X^{d}$ avec un père sain $X^{D}Y$. La mère produit deux types d'ovules ($X^{D}$ et $X^{d}$) en proportions égales ; le père produit deux types de spermatozoïdes ($X^{D}$ et $Y$). L'échiquier de croisement rassemble les quatre fécondations possibles, équiprobables :

\begin{popfigure}
\begin{center}
\begin{tikzpicture}[scale=1.0]
% Grille 3x3
\draw[thick, popdark] (0,0) grid (3,3);
% En-têtes des colonnes : gamètes du père
\node[font=\small\bfseries, popdark] at (1.5,3.4) {Gamètes du père};
\node at (1.5,2.5) {$X^{D}$};
\node at (2.5,2.5) {$Y$};
% En-têtes des lignes : gamètes de la mère
\node[font=\small\bfseries, popdark, rotate=90] at (-0.6,1.0) {Gamètes de la mère};
\node at (0.5,1.5) {$X^{D}$};
\node at (0.5,0.5) {$X^{d}$};
% Cases de fécondation
\node[fill=popgreenL, rounded corners=2pt, minimum width=0.85cm, minimum height=0.5cm] (c1) at (1.5,1.75) {$X^{D}X^{D}$};
\node[font=\small, popdark] at (1.5,1.2) {fille saine};
\node[fill=popgreenL, rounded corners=2pt, minimum width=0.85cm, minimum height=0.5cm] (c2) at (2.5,1.75) {$X^{D}Y$};
\node[font=\small, popdark] at (2.5,1.2) {garçon sain};
\node[fill=popblueL, rounded corners=2pt, minimum width=0.85cm, minimum height=0.5cm] (c3) at (1.5,0.75) {$X^{D}X^{d}$};
\node[font=\small, popdark] at (1.5,0.2) {fille conductrice};
\node[fill=poporangeL, rounded corners=2pt, minimum width=0.85cm, minimum height=0.5cm] (c4) at (2.5,0.75) {$X^{d}Y$};
\node[font=\small, popdark] at (2.5,0.2) {garçon \textbf{daltonien}};
\end{tikzpicture}
\end{center}
\legende{Échiquier de croisement $X^{D}X^{d} \times X^{D}Y$. Vert : individus sains non conducteurs ; bleu : fille conductrice saine ; orange : garçon daltonien. Chaque case a une probabilité de 1/4.}
\end{popfigure}

\begin{exemplebox}[Résultats chiffrés du croisement $X^{D}X^{d} \times X^{D}Y$]
Pour chaque naissance : 1/4 de filles saines non conductrices, 1/4 de filles saines conductrices, 1/4 de garçons sains, 1/4 de garçons daltoniens. Autrement dit :
\begin{itemize}
  \item \textbf{0 \% de filles daltoniennes} (toutes les filles sont saines, mais 50 \% d'entre elles sont conductrices) ;
  \item parmi les garçons, \textbf{50 \% de daltoniens} ;
  \item sur l'ensemble des enfants, 25 \% de daltoniens (tous des garçons).
\end{itemize}
\end{exemplebox}

\begin{attention}[Le piège classique : la transmission père $\to$ fils]
Un homme transmet \textbf{toujours} son chromosome Y à ses fils et son chromosome X à ses filles. Par conséquent :
\begin{itemize}
  \item un père daltonien $X^{d}Y$ \textbf{ne peut pas} transmettre le daltonisme à ses garçons (il leur donne Y) ;
  \item en revanche, il transmet $X^{d}$ à \textbf{toutes ses filles}, qui seront conductrices (ou daltoniennes si la mère est conductrice).
\end{itemize}
Dans un pedigree, si un père atteint a un fils atteint, le fils n'a pas reçu l'allèle de son père mais de sa mère. Ne rédige jamais « le père a transmis la maladie à son fils » pour un caractère lié à X.
\end{attention}

\begin{methode}[Reconnaître une hérédité récessive liée à X dans un pedigree]
Face à un arbre généalogique, applique la grille d'analyse suivante :
\begin{enumerate}
  \item \textbf{Qui est atteint ?} Si les malades sont presque exclusivement des hommes, pense à une hérédité gonosomique récessive.
  \item \textbf{Des parents sains ont-ils des enfants malades ?} Si oui, l'allèle est récessif.
  \item \textbf{Par qui la maladie est-elle transmise ?} Si ce sont des femmes saines (conductrices) qui ont des fils malades, le gène est porté par X.
  \item \textbf{Vérification :} un père atteint ne transmet jamais l'allèle à ses fils ; toutes ses filles reçoivent son X morbide et sont conductrices.
  \item \textbf{Conclus} en rédigeant : « l'allèle morbide est récessif et porté par le chromosome X ».
\end{enumerate}
\end{methode}

\courssub{L'hémophilie : autre maladie récessive liée à X}

L'\cle{hémophilie} est un trouble grave de la \textbf{coagulation sanguine} : un déficit en facteur de coagulation (facteur VIII pour l'hémophilie A, facteur IX pour l'hémophilie B) empêche la formation normale du caillot. Une simple coupure, une extraction dentaire ou un hématome peuvent provoquer une hémorragie prolongée, parfois mortelle. Le gène du facteur VIII est porté par le chromosome X : la transmission est donc \textbf{exactement superposable} à celle du daltonisme. En notant $X^{h}$ l'allèle morbide et $X^{H}$ l'allèle normal :

\begin{itemize}
  \item femmes : $X^{H}X^{H}$ saine, $X^{H}X^{h}$ conductrice saine, $X^{h}X^{h}$ hémophile (très rare) ;
  \item hommes : $X^{H}Y$ sain, $X^{h}Y$ hémophile.
\end{itemize}

Une mère conductrice $X^{H}X^{h}$ et un père sain donnent donc, comme pour le daltonisme, 50 \% de garçons hémophiles et des filles toutes saines dont la moitié sont conductrices. Au Cameroun, la prise en charge des hémophiles (concentrés de facteurs de coagulation, surveillance dans les hôpitaux de référence comme le CHU de Yaoundé) reste un enjeu de santé publique, et le dépistage des femmes conductrices dans les familles à risque est une application directe de cette leçon.

\begin{savaistu}[La reine Victoria, « grand-mère de l'Europe »]
La reine \textbf{Victoria d'Angleterre} (1819--1901) était conductrice de l'hémophilie. Elle a transmis l'allèle $X^{h}$ à plusieurs de ses enfants : son fils Léopold était hémophile, et ses filles, conductrices, ont épousé des souverains européens. La maladie s'est ainsi répandue dans les cours royales d'Espagne, de Prusse et de Russie : le tsarévitch \textbf{Alexis}, fils du tsar Nicolas II, était hémophile, ce qui fragilisa la dynastie des Romanov. C'est l'exemple historique le plus célèbre de transmission récessive liée à X — et un pedigree que l'on retrouve souvent au Baccalauréat !
\end{savaistu}

\courssub{Le rachitisme vitamino-résistant : une hérédité dominante liée à X}

Toutes les hérédités gonosomiques ne sont pas récessives. Le \cle{rachitisme vitamino-résistant} est une maladie des os due à une mauvaise réabsorption rénale du phosphate : les os, mal minéralisés, se déforment (jambes arquées), et — fait caractéristique — la maladie \textbf{ne guérit pas avec la vitamine D}, d'où son nom. L'allèle responsable, noté $X^{R}$, est \textbf{dominant} sur l'allèle normal $X^{r}$ et porté par le chromosome X.

\textbf{Fait de pedigree à interpréter.} Dans une famille, un père atteint et une mère saine ont cinq enfants : \textbf{toutes les filles sont atteintes, tous les fils sont sains}. Interprétons.

\begin{popfigure}
\begin{center}
\begin{tikzpicture}[
  male/.style={draw=popblue, thick, minimum size=0.55cm, inner sep=0pt},
  female/.style={draw=poppink, thick, circle, minimum size=0.55cm, inner sep=0pt},
  maleaff/.style={draw=popblue, thick, fill=popblue!70, minimum size=0.55cm, inner sep=0pt},
  femaleaff/.style={draw=poppink, thick, circle, fill=poppink!60, minimum size=0.55cm, inner sep=0pt},
  link/.style={thick, popdark},
  every label/.style={font=\small\itshape, popmuted}
]
% Génération I
\node[maleaff, label=above:{I-1 : $X^{R}Y$ (atteint)}] (I1) at (0,0) {};
\node[female, label=above:{I-2 : $X^{r}X^{r}$ (saine)}] (I2) at (1.6,0) {};
\draw[link] (I1) -- (I2);
\draw[link] (0.8,0) -- (0.8,-0.6);
\draw[link] (-1.6,-0.6) -- (3.2,-0.6);
% Génération II
\node[femaleaff, label=below:{II-1 : $X^{R}X^{r}$}] (II1) at (-1.6,-1.4) {};
\node[femaleaff, label=below:{II-2 : $X^{R}X^{r}$}] (II2) at (-0.4,-1.4) {};
\node[male, label=below:{II-3 : $X^{r}Y$}] (II3) at (0.8,-1.4) {};
\node[femaleaff, label=below:{II-4 : $X^{R}X^{r}$}] (II4) at (2.0,-1.4) {};
\node[male, label=below:{II-5 : $X^{r}Y$}] (II5) at (3.2,-1.4) {};
\draw[link] (-1.6,-0.6) -- (II1);
\draw[link] (-0.4,-0.6) -- (II2);
\draw[link] (0.8,-0.6) -- (II3);
\draw[link] (2.0,-0.6) -- (II4);
\draw[link] (3.2,-0.6) -- (II5);
\end{tikzpicture}
\end{center}
\legende{Pedigree du rachitisme vitamino-résistant (dominance liée à X). Le père atteint $X^{R}Y$ transmet son X porteur de l'allèle morbide à toutes ses filles (atteintes) et son Y à tous ses fils (sains, car leur X vient de la mère saine).}
\end{popfigure}

\textbf{Interprétation.} Le père atteint possède un seul X, porteur de $X^{R}$. Or un père donne son X à \textbf{toutes} ses filles : elles reçoivent donc toutes $X^{R}$ et, l'allèle étant dominant, elles sont toutes atteintes ($X^{R}X^{r}$). Il donne son Y à tous ses fils, qui reçoivent leur X de la mère saine ($X^{r}X^{r}$) : ils sont tous $X^{r}Y$, donc sains. Ce schéma de transmission — \emph{père atteint, toutes les filles atteintes, aucun fils atteint} — est la \textbf{signature} d'une hérédité dominante liée à X.

\begin{attention}[Dominante liée à X : deux conséquences à ne pas inverser]
\begin{itemize}
  \item Un \textbf{père} atteint transmet à toutes ses filles et à aucun de ses fils.
  \item Une \textbf{mère} atteinte hétérozygote $X^{R}X^{r}$ transmet en moyenne à la moitié de ses enfants, \textbf{garçons et filles confondus}.
  \item Contrairement au cas récessif, il n'y a \textbf{pas de conducteur sain} : tout porteur de l'allèle $X^{R}$ est malade, homme ou femme.
\end{itemize}
\end{attention}

\courssub{Bilan comparatif et exercice de synthèse}

\begin{center}
\begin{tabularx}{\linewidth}{|l|X|X|}
\hline
\rowcolor{popblueL} \textbf{Critère} & \textbf{Gonosomique récessive} (daltonisme, hémophilie) & \textbf{Gonosomique dominante} (rachitisme vitamino-résistant) \\
\hline
Sexe le plus atteint & Hommes (un seul X suffit) & Femmes environ deux fois plus souvent (deux X) \\
\hline
Conducteur sain & Oui : la femme $X^{D}X^{d}$ & Non : tout porteur est malade \\
\hline
Père atteint & Aucun fils atteint ; toutes les filles conductrices saines & Aucun fils atteint ; toutes les filles \textbf{atteintes} \\
\hline
Saut de génération & Fréquent (allèle caché chez les conductrices) & Rare (tout porteur s'exprime) \\
\hline
\end{tabularx}
\end{center}

\begin{aretenir}[L'essentiel de la section]
\begin{itemize}
  \item Une hérédité gonosomique est liée aux chromosomes sexuels ; au programme, le gène est porté par \textbf{X}.
  \item \textbf{Récessive liée à X} (daltonisme, hémophilie) : surtout des hommes atteints, transmission par les femmes conductrices saines ; croisement $X^{D}X^{d} \times X^{D}Y$ $\rightarrow$ 50 \% de garçons malades, 0 \% de filles malades.
  \item \textbf{Dominante liée à X} (rachitisme vitamino-résistant) : un père atteint a toutes ses filles atteintes et aucun fils atteint.
  \item Un homme ne transmet \textbf{jamais} son X à ses fils.
\end{itemize}
\end{aretenir}

\begin{exoresolu}[Un couple inquiet avant un mariage à Douala]
\textbf{Énoncé.} Awa, dont le père est daltonien, envisage de se marier avec Joel, dont la vision des couleurs est normale. Awa est elle-même saine. (1) Détermine le génotype d'Awa en le justifiant. (2) Construis l'échiquier de croisement du couple et évalue le risque, pour ce couple, d'avoir un enfant daltonien. (3) Précise le sexe des enfants éventuellement atteints.
\tcblower
\textbf{Solution.} (1) Le père d'Awa est daltonien : son génotype est $X^{d}Y$. Or un père transmet obligatoirement son chromosome X à toutes ses filles. Awa a donc reçu $X^{d}$ de son père. Étant saine, elle possède un allèle normal dominant sur son second X, reçu de sa mère. Awa est donc $X^{D}X^{d}$ : \textbf{conductrice saine}.

(2) Joel est sain : $X^{D}Y$. Le croisement $X^{D}X^{d} \times X^{D}Y$ donne : $X^{D}X^{D}$ (fille saine, 1/4), $X^{D}X^{d}$ (fille conductrice saine, 1/4), $X^{D}Y$ (garçon sain, 1/4), $X^{d}Y$ (garçon daltonien, 1/4).

(3) Le risque d'avoir un enfant daltonien est de \textbf{1/4 (25 \%) par naissance}, et seuls des \textbf{garçons} peuvent être atteints : le risque est de 50 \% pour chaque garçon et de 0 \% pour chaque fille. Ce type de raisonnement est exactement celui qui fonde les examens prénuptiaux, que nous étudierons dans la dernière section.
\end{exoresolu}

\begin{exercice}[À toi de jouer]
Dans une famille, un homme atteint de rachitisme vitamino-résistant épouse une femme saine. Ils ont trois filles et deux fils. (1) Donne les génotypes des parents. (2) Prédis le phénotype de chaque enfant. (3) L'un des fils, sain, épouse une femme saine : leurs enfants peuvent-ils être atteints ? Justifie.
\end{exercice}

\begin{corrige}[Exercice : rachitisme vitamino-résistant]
(1) Père atteint : $X^{R}Y$ ; mère saine : $X^{r}X^{r}$. (2) Toutes les filles reçoivent $X^{R}$ du père et $X^{r}$ de la mère : elles sont $X^{R}X^{r}$, toutes atteintes. Les fils reçoivent Y du père et $X^{r}$ de la mère : ils sont $X^{r}Y$, tous sains. (3) Le fils sain ($X^{r}Y$) ne possède pas l'allèle morbide ; avec une femme saine ($X^{r}X^{r}$), tous les enfants seront $X^{r}X^{r}$ ou $X^{r}Y$ : \textbf{aucun enfant ne peut être atteint}, car aucun parent ne transmet $X^{R}$.
\end{corrige}

\coursec{Évaluation d'un risque génétique}

Dans la section précédente, nous avons appris à lire un pedigree et à identifier le mode de transmission d'une maladie héréditaire : autosomique ou gonosomique, récessive ou dominante. Mais identifier le mode de transmission ne suffit pas toujours. Imaginons un couple de jeunes mariés de Yaoundé dont les deux partenaires sont de phénotype sain, mais dont les familles comptent des cas de drépanocytose. Ces futurs parents se posent une question très concrète : \emph{quel est le risque que notre enfant soit atteint ?} Répondre à cette question, c'est évaluer un \cle{risque génétique}. C'est l'une des applications les plus importantes de la génétique humaine, et une compétence exigible au Baccalauréat.

\courssub{Qu'est-ce qu'un risque génétique ?}

\begin{definition}[Risque génétique]
Le \cle{risque génétique} est la \textbf{probabilité}, exprimée en fraction ou en pourcentage, qu'un enfant d'un couple donné soit atteint d'une maladie héréditaire déterminée. Il se calcule à partir des génotypes des parents, déduits de l'étude de l'arbre généalogique (pedigree) et du mode de transmission de la maladie.
\end{definition}

L'idée intuitive est simple : la fécondation est un tirage au sort. Chaque parent produit des gamètes qui ne portent qu'\textbf{un seul} des deux allèles de chaque gène, et la rencontre des gamètes est due au hasard. Connaître les génotypes parentaux permet donc de prévoir, non pas le destin d'un enfant précis, mais la \textbf{proportion attendue} d'enfants atteints parmi toutes les naissances possibles de ce couple.

\begin{attention}[Un risque n'est pas une prédiction individuelle]
Un risque de \SI{25}{\%} ne signifie \textbf{pas} « un enfant atteint sur quatre dans la famille ». Chaque grossesse est un \textbf{événement indépendant} : le hasard n'a pas de mémoire. Un couple à risque de \SI{25}{\%} peut avoir quatre enfants sains, ou deux enfants atteints de suite. De même, avoir déjà eu un enfant atteint ne « protège » pas les naissances suivantes : le risque reste de \SI{25}{\%} à chaque grossesse. C'est l'erreur la plus fréquente au Baccalauréat.
\end{attention}

\courssub{La démarche d'évaluation d'un risque génétique}

\begin{methode}[Les trois étapes du calcul d'un risque génétique]
\begin{enumerate}
\item \textbf{Déterminer les génotypes des parents} à partir du pedigree et du mode de transmission (un individu atteint d'une maladie récessive est obligatoirement homozygote ; un parent sain ayant un enfant atteint d'une maladie récessive est obligatoirement hétérozygote conducteur).
\item \textbf{Construire l'échiquier de croisement} en plaçant les gamètes de chaque parent en ligne et en colonne, puis en remplissant les cases (génotypes possibles des enfants).
\item \textbf{Compter les cases} correspondant au phénotype atteint et exprimer le résultat en fraction ou en pourcentage. Si la question porte sur un sexe particulier (maladie gonosomique), préciser si le risque est calculé parmi tous les enfants ou parmi les enfants d'un sexe donné.
\end{enumerate}
\end{methode}

\courssub{Cas 1 : maladie autosomique récessive — deux parents conducteurs sains}

\textbf{Observation.} Un couple, tous deux de phénotype normal, a déjà un enfant albinos. Comment est-ce possible, et quel est le risque pour les naissances suivantes ?

\textbf{Raisonnement.} L'albinisme est une maladie autosomique récessive : l'allèle normal $A$ (pigmentation normale) domine l'allèle muté $a$ (albinisme). L'enfant albinos est de génotype $a/a$ : il a reçu un allèle $a$ de chaque parent. Les deux parents, bien que sains, sont donc obligatoirement \textbf{hétérozygotes conducteurs} $A/a$ (on dit aussi « porteurs sains »).

\textbf{Échiquier de croisement} $A/a \times A/a$ :

\begin{center}
\begin{tabular}{|c|c|c|}
\hline
\rowcolor{popblueL} Gamètes & $A$ (\SI{1}{/2}) & $a$ (\SI{1}{/2}) \\
\hline
$A$ (\SI{1}{/2}) & $A/A$ sain (\SI{1}{/4}) & $A/a$ sain conducteur (\SI{1}{/4}) \\
\hline
$a$ (\SI{1}{/2}) & $A/a$ sain conducteur (\SI{1}{/4}) & $a/a$ \textbf{albinos} (\SI{1}{/4}) \\
\hline
\end{tabular}
\end{center}

\begin{propriete}[Risque pour un couple de porteurs sains]
Pour toute maladie autosomique récessive, le croisement de deux hétérozygotes $A/a \times A/a$ donne à chaque naissance :
\begin{itemize}
\item \SI{1}{/4} (\SI{25}{\%}) d'enfant homozygote sain $A/A$ ;
\item \SI{2}{/4} (\SI{50}{\%}) d'enfants hétérozygotes sains conducteurs $A/a$ ;
\item \SI{1}{/4} (\SI{25}{\%}) d'enfant atteint $a/a$.
\end{itemize}
Le risque génétique d'avoir un enfant atteint est donc de \textbf{\SI{25}{\%} à chaque naissance}, quel que soit le sexe de l'enfant (caractère autosomique).
\end{propriete}

\courssub{Cas 2 : la drépanocytose, un enjeu majeur de santé publique au Cameroun}

\textbf{Situation.} La drépanocytose (anémie falciforme) est particulièrement fréquente en Afrique subsaharienne : au Cameroun, on estime qu'environ \SI{20}{\%} à \SI{25}{\%} de la population est porteuse du trait drépanocytaire (génotype $A/S$, dit « AS »), c'est-à-dire hétérozygote et cliniquement saine. L'allèle $S$ code une hémoglobine anormale qui fait prendre aux globules rouges une forme de faucille en cas de faible oxygénation, provoquant crises douloureuses et anémie grave chez les homozygotes $S/S$.

\textbf{Évaluation du risque.} Si deux porteurs sains $A/S$ se marient — situation fréquente faute de dépistage systématique —, le raisonnement est identique au cas précédent : $A/S \times A/S$ donne \SI{1}{/4} $A/A$, \SI{2}{/4} $A/S$, \SI{1}{/4} $S/S$. Le risque d'avoir un enfant drépanocytaire ($S/S$) est de \textbf{\SI{25}{\%} à chaque grossesse}.

\begin{popfigure}
\begin{center}
\begin{tikzpicture}[
grow=right, level distance=3.4cm,
level 1/.style={sibling distance=2.6cm},
level 2/.style={sibling distance=1.3cm},
every node/.style={font=\small},
edge from parent/.style={draw=popmuted, thick}]
\node[draw=popdark, fill=popcream, rounded corners, thick, inner sep=4pt] {Fécondation $A/S \times A/S$}
  child { node[draw=poporange, fill=poporangeL, rounded corners, inner sep=3pt] {Ovule $S$ (\SI{1}{/2})}
    child { node[draw=popink, fill=poppinkL, rounded corners, inner sep=3pt] {$S/S$ atteint : \SI{1}{/4}}
      edge from parent node[above, font=\scriptsize] {$S$ (\SI{1}{/2})} }
    child { node[draw=popgreen, fill=popgreenL, rounded corners, inner sep=3pt] {$A/S$ sain : \SI{1}{/4}}
      edge from parent node[above, font=\scriptsize] {$A$ (\SI{1}{/2})} }
    edge from parent node[above, font=\scriptsize] {\SI{1}{/2}} }
  child { node[draw=popblue, fill=popblueL, rounded corners, inner sep=3pt] {Ovule $A$ (\SI{1}{/2})}
    child { node[draw=popgreen, fill=popgreenL, rounded corners, inner sep=3pt] {$A/S$ sain : \SI{1}{/4}}
      edge from parent node[above, font=\scriptsize] {$S$ (\SI{1}{/2})} }
    child { node[draw=popgreen, fill=popgreenL, rounded corners, inner sep=3pt] {$A/A$ sain : \SI{1}{/4}}
      edge from parent node[above, font=\scriptsize] {$A$ (\SI{1}{/2})} }
    edge from parent node[above, font=\scriptsize] {\SI{1}{/2}} };
\end{tikzpicture}
\end{center}
\legende{Arbre de probabilités du croisement $A/S \times A/S$ (drépanocytose). Chaque branche porte la probabilité du gamète (\SI{1}{/2}) ; la probabilité d'un génotype est le produit des probabilités le long du chemin. Bilan : \SI{25}{\%} $S/S$ (enfant atteint), \SI{50}{\%} $A/S$ (porteurs sains), \SI{25}{\%} $A/A$ (sains non porteurs).}
\end{popfigure}

\begin{savaistu}[Pourquoi l'allèle $S$ est-il si fréquent en Afrique ?]
Les hétérozygotes $A/S$ résistent mieux au paludisme grave que les homozygotes $A/A$ : le parasite \emph{Plasmodium falciparum} se développe mal dans leurs globules rouges. Dans les régions où le paludisme sévit, comme au Cameroun, la sélection naturelle a donc maintenu l'allèle $S$ à une fréquence élevée malgré la gravité de la drépanocytose homozygote. C'est un exemple célèbre d'\textbf{avantage sélectif de l'hétérozygote}.
\end{savaistu}

\courssub{Cas 3 : maladie gonosomique récessive — femme conductrice et homme sain}

\textbf{Situation.} Une femme dont le père était hémophile épouse un homme sain. Quel est le risque pour leurs enfants ?

\textbf{Raisonnement.} L'hémophilie est récessive et liée au chromosome $X$. La femme, saine mais fille d'un père hémophile ($X^{h}Y$), a obligatoirement reçu le $X^{h}$ paternel : elle est conductrice $X^{H}X^{h}$. L'homme sain est $X^{H}Y$.

\begin{exemplebox}[Échiquier de croisement $X^{H}X^{h} \times X^{H}Y$]
\begin{center}
\begin{tabular}{|c|c|c|}
\hline
\rowcolor{popblueL} Gamètes & $X^{H}$ (\SI{1}{/2}) & $Y$ (\SI{1}{/2}) \\
\hline
$X^{H}$ (\SI{1}{/2}) & $X^{H}X^{H}$ fille saine (\SI{1}{/4}) & $X^{H}Y$ garçon sain (\SI{1}{/4}) \\
\hline
$X^{h}$ (\SI{1}{/2}) & $X^{H}X^{h}$ fille conductrice (\SI{1}{/4}) & $X^{h}Y$ \textbf{garçon hémophile} (\SI{1}{/4}) \\
\hline
\end{tabular}
\end{center}
Bilan : \SI{1}{/4} $X^{H}X^{H}$, \SI{1}{/4} $X^{H}X^{h}$, \SI{1}{/4} $X^{H}Y$, \SI{1}{/4} $X^{h}Y$.
\end{exemplebox}

\begin{propriete}[Risque pour une femme conductrice et un homme sain]
Pour une maladie récessive liée au chromosome $X$ (hémophilie, daltonisme), le croisement $X^{H}X^{h} \times X^{H}Y$ donne :
\begin{itemize}
\item \textbf{aucune fille atteinte} (toutes reçoivent le $X^{H}$ paternel), mais \SI{1}{/2} fille sur \SI{2}{/2} est conductrice ;
\item \textbf{1 garçon sur 2 atteint}, soit un risque de \textbf{\SI{50}{\%} parmi les garçons} ;
\item un risque de \textbf{\SI{25}{\%} sur l'ensemble des naissances} (tous sexes confondus) d'avoir un enfant atteint.
\end{itemize}
\end{propriete}

\begin{attention}[Bien préciser la population de référence]
« \SI{50}{\%} de garçons hémophiles » et « \SI{25}{\%} d'enfants atteints » sont \textbf{les deux mêmes réalités} exprimées différemment : la première probabilité est calculée parmi les garçons, la seconde parmi toutes les naissances. Au Baccalauréat, lis attentivement la question : « risque d'avoir un garçon hémophile » (\SI{1}{/4} des naissances) n'est pas « risque qu'un garçon soit hémophile » (\SI{1}{/2} des garçons).
\end{attention}

\courssub{Cas 4 : maladie autosomique dominante — un parent hétérozygote atteint}

\textbf{Situation.} Un homme atteint d'achondroplasie (nanisme avec membres courts), dont le père était également atteint mais dont la mère était de taille normale, épouse une femme de taille normale. Quel est le risque pour leurs enfants ?

\textbf{Raisonnement.} L'achondroplasie est autosomique \textbf{dominante} : un seul allèle muté suffit pour être atteint. Notons $N$ l'allèle normal et $A$ l'allèle muté. L'homme atteint a une mère saine $N/N$ : il a donc reçu $N$ de sa mère et $A$ de son père ; il est hétérozygote $A/N$ (les homozygotes $A/A$ sont d'ailleurs généralement non viables). La femme saine est $N/N$.

\textbf{Échiquier de croisement} $A/N \times N/N$ :

\begin{center}
\begin{tabular}{|c|c|}
\hline
\rowcolor{popblueL} Gamètes & $N$ (mère, \SI{100}{\%}) \\
\hline
$A$ (père, \SI{1}{/2}) & $A/N$ \textbf{atteint} (\SI{1}{/2}) \\
\hline
$N$ (père, \SI{1}{/2}) & $N/N$ sain (\SI{1}{/2}) \\
\hline
\end{tabular}
\end{center}

\begin{propriete}[Risque pour une maladie autosomique dominante]
Un hétérozygote atteint croisé avec un homozygote sain ($A/N \times N/N$) transmet la maladie à \textbf{\SI{50}{\%} de ses enfants, à chaque naissance}, indifféremment aux deux sexes. À la différence des maladies récessives, il n'existe pas de « porteur sain » : tout individu porteur de l'allèle muté est atteint.
\end{propriete}

\courssub{Tableau récapitulatif des risques génétiques}

\begin{center}
\begin{tabularx}{\linewidth}{|l|X|X|X|}
\hline
\rowcolor{popblueL} \textbf{Mode de transmission} & \textbf{Génotypes parentaux} & \textbf{Risque d'enfant atteint} & \textbf{Exemple} \\
\hline
Autosomique récessive & $A/a \times A/a$ (deux porteurs sains) & \SI{25}{\%} à chaque naissance, tout sexe & Albinisme, drépanocytose ($A/S \times A/S$) \\
\hline
Autosomique récessive & $A/A \times A/a$ (sain non porteur $\times$ conducteur) & \SI{0}{\%} (\SI{50}{\%} de conducteurs sains) & Albinisme \\
\hline
Autosomique dominante & $A/N \times N/N$ (atteint hétérozygote $\times$ sain) & \SI{50}{\%} à chaque naissance, tout sexe & Achondroplasie \\
\hline
Gonosomique récessive (liée à $X$) & $X^{H}X^{h} \times X^{H}Y$ (conductrice $\times$ homme sain) & \SI{50}{\%} des garçons, soit \SI{25}{\%} des naissances & Hémophilie, daltonisme \\
\hline
Gonosomique récessive (liée à $X$) & $X^{H}X^{H} \times X^{h}Y$ (femme saine $\times$ homme atteint) & \SI{0}{\%} d'enfants atteints (filles toutes conductrices) & Hémophilie \\
\hline
Gonosomique dominante (liée à $X$) & $X^{R}X^{r} \times X^{r}Y$ (femme atteinte hétérozygote $\times$ homme sain) & \SI{50}{\%} des enfants, tout sexe & Rachitisme vitamino-résistant \\
\hline
\end{tabularx}
\end{center}

\courssub{Le conseil génétique : de la science à la décision}

\begin{definition}[Conseil génétique]
Le \cle{conseil génétique} est une démarche médicale par laquelle un médecin (généticien) informe un couple, \textbf{avant la conception} ou en début de grossesse, des risques de transmission d'une maladie héréditaire, à partir de l'étude des arbres généalogiques et, si nécessaire, d'examens biologiques (électrophorèse de l'hémoglobine, caryotype, tests sanguins). Il s'agit d'une \textbf{information} et non d'une injonction : la décision finale appartient toujours au couple.
\end{definition}

Le conseil génétique est particulièrement recommandé lorsqu'il existe des cas de maladie héréditaire dans la famille, lors de mariages entre apparentés (consanguinité, qui augmente la probabilité que les deux conjoints portent le même allèle récessif), ou dans les populations où certains allèles sont fréquents, comme l'allèle $S$ de la drépanocytose au Cameroun. Il prépare naturellement les examens prénuptiaux et prénataux que nous étudierons dans la section suivante.

\begin{exoresolu}[Risque d'hémophilie dans une famille]
\textbf{Énoncé.} Dans une famille, une femme (III-2) de phénotype sain a un frère hémophile. Leur mère (II-3) est saine, ainsi que leur père (II-2). III-2 épouse un homme sain. 1) Déterminer le génotype de la mère II-3. 2) Quel est le risque que III-2 soit conductrice ? 3) Quel est le risque que le couple III-2 ait un garçon hémophile ?
\tcblower
\textbf{Solution détaillée.}
\begin{enumerate}
\item Le frère hémophile est $X^{h}Y$ : il a reçu $X^{h}$ de sa mère (le père donne $Y$ à ses fils). La mère II-3, saine, est donc obligatoirement conductrice : $X^{H}X^{h}$.
\item III-2 est une fille saine issue du croisement $X^{H}X^{h} \times X^{H}Y$. Parmi les filles, deux génotypes équiprobables : $X^{H}X^{H}$ et $X^{H}X^{h}$. Le risque qu'elle soit conductrice est donc de \textbf{\SI{1}{/2} (\SI{50}{\%})}.
\item Deux probabilités se combinent : probabilité que III-2 soit conductrice (\SI{1}{/2}) $\times$ probabilité, si elle l'est, d'avoir un garçon hémophile (\SI{1}{/4} des naissances, car $X^{h}Y$ représente une case sur quatre de l'échiquier). Risque total : $\dfrac{1}{2} \times \dfrac{1}{4} = \dfrac{1}{8}$, soit \textbf{\SI{12,5}{\%} des naissances} (ou, de façon équivalente, \SI{1}{/2} $\times$ \SI{1}{2} $=$ \SI{1}{4} des garçons).
\end{enumerate}
\end{exoresolu}

\begin{exercice}[À toi de jouer]
Deux époux, tous deux de phénotype sain, consultent avant leur mariage à Douala. L'électrophorèse de l'hémoglobine révèle qu'ils sont tous deux de génotype $A/S$. 1) Construis l'échiquier de croisement. 2) Calcule le risque, à chaque grossesse, d'avoir un enfant drépanocytaire, puis un enfant porteur sain. 3) Le couple a déjà un enfant $S/S$. Quel est le risque pour la grossesse suivante ? Justifie ta réponse.
\end{exercice}

\begin{corrige}[Exercice — À toi de jouer]
1) L'échiquier $A/S \times A/S$ comporte quatre cases : $A/A$, $A/S$, $A/S$, $S/S$.
2) Risque d'enfant drépanocytaire $S/S$ : \SI{1}{/4}, soit \textbf{\SI{25}{\%}}. Risque d'enfant porteur sain $A/S$ : \SI{2}{/4}, soit \textbf{\SI{50}{\%}}.
3) Le risque reste de \textbf{\SI{25}{\%}} : chaque grossesse est un événement indépendant, les naissances précédentes ne modifient pas les probabilités.
\end{corrige}

\begin{aretenir}[L'essentiel de la section]
\begin{itemize}
\item Le risque génétique est la probabilité qu'un enfant d'un couple donné soit atteint d'une maladie héréditaire ; il se calcule en trois étapes : génotypes parentaux, échiquier de croisement, comptage des cases.
\item Deux porteurs sains d'une maladie autosomique récessive : \SI{25}{\%} d'enfants atteints (albinisme, drépanocytose).
\item Femme conductrice $\times$ homme sain (maladie liée à $X$) : \SI{50}{\%} des garçons atteints, soit \SI{25}{\%} des naissances.
\item Maladie autosomique dominante (parent hétérozygote) : \SI{50}{\%} d'enfants atteints.
\item Les probabilités s'appliquent à \textbf{chaque naissance indépendamment} des précédentes.
\item Le conseil génétique informe le couple de ces risques avant la conception ; la décision lui appartient.
\end{itemize}
\end{aretenir}

\coursec{Applications : examens prénuptiaux et prénataux}

Après avoir appris à \emph{évaluer un risque génétique} à partir d'arbres généalogiques (section précédente), nous disposons maintenant d'un outil puissant : la prévision. Mais prévoir ne suffit pas ; il faut pouvoir \emph{agir} en amont pour limiter l'apparition des maladies héréditaires au sein des familles et de la population. C'est tout l'objet de la génétique médicale appliquée : proposer des examens de dépistage \textbf{avant} la conception (examen prénuptial) et \textbf{pendant} la grossesse (examen prénatal). Au Cameroun, comme dans de nombreux pays, ces examens s'inscrivent dans une politique de santé publique visant à réduire la charge des maladies génétiques, notamment la drépanocytose, première maladie génétique du pays.

\courssub{L'examen prénuptial : informer avant l'union}

\begin{definition}[Examen prénuptial]
L'\cle{examen prénuptial} (ou bilan prénuptial) est l'ensemble des examens médicaux et biologiques réalisés \textbf{avant le mariage} (ou avant tout projet de grossesse) chez les deux futurs époux. Son but est de dépister d'éventuelles anomalies génétiques, infectieuses ou immunologiques susceptibles d'être transmises à la descendance ou de compromettre la grossesse.
\end{definition}

\begin{savaistu}[Le certificat prénuptial au Cameroun]
Au Cameroun, le Code de la Famille et les textes d'application (notamment l'arrêté n°004/MSP/SG/DOS du 19 janvier 2002) rendent \textbf{obligatoire la présentation d'un certificat prénuptial} pour tout mariage civil. Ce certificat, délivré par un médecin habilité après examen des deux futurs époux, atteste notamment : 
\begin{itemize}
    \item du groupe sanguin et du \textbf{Rhésus} de chacun ;
    \item du \textbf{statut drépanocytaire} (électrophorèse de l'hémoglobine) ;
    \item de la sérologie VIH, syphilis, hépatites B et C.
\end{itemize}
Ce certificat a une validité de \num{3} mois. Il ne \emph{conditionne pas} la célébration du mariage (le mariage ne peut être refusé pour raison médicale), mais il \emph{informe} le couple pour une procréation responsable.
\end{savaistu}

\paragraph{Le dépistage de la drépanocytose par électrophorèse de l'hémoglobine.}
La drépanocytose (anémie falciforme) est une maladie autosomique récessive liée à l'allèle $S$ du gène $\beta$-globine ($HBB$). Le génotype normal est $AA$ (hémoglobine A adulte), le génotype hétérozygote est $AS$ (porteur sain, trait drépanocytaire), le génotype homozygote mutant est $SS$ (malade drépanocytaire majeur). L'électrophorèse sur gel (ou en capillaire) sépare les hémoglobines selon leur charge électrique et leur taille à pH alcalin (pH \num{8,6}).

\begin{methode}[Interprétation d'un électrophérogramme d'hémoglobine pour le conseil prénuptial]
\begin{enumerate}
    \item \textbf{Observer les bandes :} À pH alcalin, l'hémoglobine A (HbA) et l'hémoglobine F (HbF) migrent rapidement vers l'anode (charge négative forte), l'hémoglobine S (HbS) migre moins vite.
    \item \textbf{Identifier le profil :}
    \begin{itemize}
        \item Une seule bande en position HbA/HbF $\rightarrow$ Profil \textbf{AA} (non porteur).
        \item Deux bandes : HbA (majoritaire, $\approx$ \SI{60}{\%}) et HbS (minoritaire, $\approx$ \SI{40}{\%}) $\rightarrow$ Profil \textbf{AS} (porteur sain / trait drépanocytaire).
        \item Une seule bande en position HbS (éventuellement HbF résiduelle en position HbA) $\rightarrow$ Profil \textbf{SS} (malade drépanocytaire majeur).
    \end{itemize}
    \item \textbf{Conseiller le couple :} Croiser les profils des deux partenaires pour évaluer le risque à chaque grossesse (voir tableau ci-dessous).
\end{enumerate}
\end{methode}

\begin{popfigure}
\begin{tikzpicture}[scale=0.9, every node/.style={font=\small}]
    % Axes
    \draw[->] (0,0) -- (12,0) node[right] {Distance de migration (cm) $\rightarrow$};
    \draw[->] (0,0) -- (0,5) node[above] {Concentration / Absorbance};
    
    % Point de départ (cathode)
    \draw[dashed] (1,0) -- (1,4.5);
    \node[below] at (1,0) {Origine (cathode)};
    
    % Anode direction
    \node[above right] at (11,4.5) {Anode (+)};
    
    % Profils : AA (HbA + HbF résiduelle)
    \draw[popblue, thick, domain=1:11, samples=100] plot (\x, {3*exp(-(\x-3.5)^2/0.8)});
    \node[popblue, anchor=south] at (3.5, 3.2) {\textbf{AA} : HbA (+ HbF)};
    
    % Profils : AS
    \draw[poporange, thick, domain=1:11, samples=100] plot (\x, {2.2*exp(-(\x-3.5)^2/0.8) + 1.5*exp(-(\x-7)^2/1.0)});
    \node[poporange, anchor=south] at (3.5, 2.4) {\textbf{AS} : HbA};
    \node[poporange, anchor=south] at (7, 1.7) {\textbf{AS} : HbS};
    
    % Profils : SS
    \draw[popred, thick, domain=1:11, samples=100] plot (\x, {2.5*exp(-(\x-7)^2/1.0)});
    \node[popred, anchor=south] at (7, 2.7) {\textbf{SS} : HbS (+ HbF)};
    
    % Repères de migration
    \draw[dashed, thin] (3.5,0) -- (3.5,3.5);
    \draw[dashed, thin] (7,0) -- (7,3);
    \node[below, font=\tiny] at (3.5,0) {Zone HbA / HbF};
    \node[below, font=\tiny] at (7,0) {Zone HbS};
\end{tikzpicture}
\legende{Profils d'électrophorèse de l'hémoglobine (pH alcalin \num{8,6}) pour le dépistage prénuptial. HbA et HbF migrent vite vers l'anode ; HbS migre plus lentement. AA : bande unique zone rapide. AS : deux bandes. SS : bande unique zone lente (éventuel pic HbF en zone rapide).}
\end{popfigure}

\begin{center}
\begin{tabularx}{\linewidth}{|l|X|X|X|}\hline
\rowcolor{popblueL}
\textbf{Croisement parental} & \textbf{Génotypes parentaux} & \textbf{Risque à chaque grossesse} & \textbf{Conseil génétique} \\ \hline
Non porteur $\times$ Non porteur & $AA \times AA$ & \SI{0}{\%} malade, \SI{0}{\%} porteur & Risque nul. \\ \hline
Non porteur $\times$ Porteur & $AA \times AS$ & \SI{0}{\%} malade, \SI{50}{\%} porteurs (AS) & Enfants sains ; 1 sur 2 porteur sain. \\ \hline
Porteur $\times$ Porteur & $AS \times AS$ & \textbf{\SI{25}{\%} malades (SS)}, \SI{50}{\%} porteurs (AS), \SI{25}{\%} sains (AA) & \textbf{Risque élevé}. Information sur la maladie, option de diagnostic prénatal. \\ \hline
Porteur $\times$ Malade & $AS \times SS$ & \SI{50}{\%} malades (SS), \SI{50}{\%} porteurs (AS) & Risque très élevé. \\ \hline
Malade $\times$ Malade & $SS \times SS$ & \SI{100}{\%} malades (SS) & Certitude de transmission. \\ \hline
\end{tabularx}
\end{center}

\begin{attention}[Erreur fréquente : l'examen prénuptial n'interdit pas le mariage]
L'examen prénuptial (et le certificat qui en découle) a une visée \textbf{d'information et de conseil}, non de \textbf{contrainte juridique}. Au Cameroun, l'article 194 du Code Civil dispose que « le mariage ne peut être refusé pour cause d'inaptitude physique ». Le médecin \textbf{ne doit pas} dire « vous ne devez pas vous marier », mais « \emph{si vous avez des enfants, voici le risque génétique encouru à chaque grossesse, et voici les options qui s'offrent à vous (diagnostic prénatal, adoption, don de gamètes, etc.)} ». La décision finale revient au couple (autonomie reproductive).
\end{attention}

\paragraph{L'incompatibilité Rhésus (système Rhésus D).}
L'examen prénuptial dépiste aussi le phénotype Rhésus (Rh1/D). Le risque survient si la mère est \textbf{Rhésus négatif (Rh- / $dd$)} et le père \textbf{Rhésus positif (Rh+ / $DD$ ou $Dd$)}. Le fœtus peut être Rh+. Lors du premier accouchement (ou avortement, grossesse extra-utérine, amniocentèse), du sang fœtal Rh+ peut passer dans la circulation maternelle $\rightarrow$ \textbf{immunisation (sensibilisation)} de la mère : elle produit des anti-D (IgG) qui traversent le placenta lors des grossesses ultérieures $\rightarrow$ \textbf{maladie hémolytique du nouveau-né (ictère grave, anémie, hydrops fœtalis)}.
\begin{propriete}[Prévention de l'alloinmunisation anti-D]
L'examen prénuptial permet d'identifier les couples à risque (Femme Rh- $\times$ Homme Rh+). La prévention repose sur l'injection d'\textbf{immunoglobulines anti-D} (Rhophylac\textsuperscript{\textregistered} ou équivalent) à la mère Rh- :
\begin{itemize}
    \item À \num{28}-\num{30} SA (prophylaxie anténatale de routine) ;
    \item Dans les \num{72}h après tout événement à risque (accouchement, IVG, amniocentèse, traumatisme abdominal) si le nouveau-né est Rh+.
\end{itemize}
Cela détruit les globules rouges fœtaux Rh+ entrés dans la circulation maternelle \emph{avant} qu'elle ne produise ses propres anticorps.
\end{propriete}

\courssub{L'examen prénatal : diagnostiquer pendant la grossesse}

\begin{definition}[Examen prénatal (Diagnostic prénatal - DP)]
L'\cle{examen prénatal} (ou diagnostic prénatal) est l'ensemble des techniques médicales permettant de \textbf{détecter in utero} (pendant la grossesse) des anomalies chromosomiques, génétiques, structurelles ou infectieuses chez le fœtus. Il ne concerne que les grossesses à risque identifié (antécédents, âge maternel, examen prénuptial, marqueurs sériques, échographie).
\end{definition}

\begin{aretenir}[Deux niveaux d'examen prénatal]
\begin{enumerate}
    \item \textbf{Le dépistage prénatal (non invasif) :} Évalue une \emph{probabilité} (risque). Ex : dosage des marqueurs sériques maternels (PAPP-A, $\beta$-hCG libre) au 1er trimestre + clarté nucale à l'échographie $\rightarrow$ score de risque pour trisomie 21. Si risque élevé ($> 1/250$), on propose un diagnostic.
    \item \textbf{Le diagnostic prénatal (invasif) :} Donne un \emph{résultat certain} (caryotype, ADN). Ex : \textbf{amniocentèse}, choriocentèse (prélèvement de trophoblaste). Risque de fausse couche $\approx$ \SI{0,5}{\%} à \SI{1}{\%}.
\end{enumerate}
\end{aretenir}

\paragraph{L'échographie : imagerie morphologique.}
Réalisée aux 3 trimestres (\num{11}-\num{13} SA, \num{22}-\num{24} SA, \num{32}-\num{34} SA). Elle détecte les \textbf{malformations structurelles} : clarté nucale augmentée (signe d'appel trisomie 21), anomalies cardiaques, spina bifida, agénésies, dysmorphies faciales, restriction de croissance intra-utérine (RCIU). Elle guide aussi les gestes invasifs (repérage du placenta, du cordon, poche de liquide amniotique).

\paragraph{L'amniocentèse : le standard pour le caryotype fœtal.}
C'est le prélèvement de \textbf{liquide amniotique} (riche en cellules fœtales exfoliées : cellules amniotiques, cellules cutanées, urinaires) par ponction transabdominale sous contrôle échographique, généralement entre \textbf{\num{15} et \num{18} SA} (amniocentèse tardive) ou plus tôt (amniocentèse précoce $\approx$ \num{14} SA, plus risquée).

\begin{popfigure}
\begin{tikzpicture}[scale=0.85, every node/.style={font=\small}]
    % Utérus
    \draw[popdark, thick, fill=poppinkL!30] (0,0) .. controls (0,2) and (6,3) .. (6,0) .. controls (6,-1) and (0,-1) .. cycle;
    % Paroi abdominale
    \draw[popdark!70!black, thick, fill=popdark!10] (-0.5,-0.2) -- (6.5,-0.2) -- (6.5,-1.5) -- (-0.5,-1.5) -- cycle;
    \node[popdark!70!black] at (3,-0.85) {Paroi abdominale maternelle};
    
    % Liquide amniotique
    \draw[popblueL, fill=popblueL!20] (1,0.2) .. controls (1,1.5) and (5,2) .. (5,0.2) .. controls (5,-0.2) and (1,-0.2) .. cycle;
    \node[popblue] at (3,1.2) {Liquide amniotique};
    
    % Fœtus
    \draw[popdark, fill=popcream] (3.5,0.8) ellipse (1.2 and 0.7);
    \node at (3.5,0.8) {Fœtus};
    
    % Placenta
    \draw[popred, fill=popred!20] (1.2, -0.1) rectangle (2.5, 1.5);
    \node[popred] at (1.85, 0.7) {Placenta};
    
    % Aiguille
    \draw[poporange, very thick, ->] (3, -1.5) -- (3, 0.5);
    \draw[poporange, thick] (2.8, -1.2) -- (3.2, -1.2); % Poignée seringue
    \node[poporange, left] at (2.5, -1.0) {Aiguille (20-22G)};
    \node[poporange, left] at (2.5, -1.4) {sous contrôle écho};
    
    % Prélèvement
    \draw[popblue, thick, ->] (3, 0.5) -- (3, 1.0);
    \node[popblue, right] at (3.3, 0.75) {Prélèvement \SI{15}{\mL}--\SI{20}{\mL}};
    
    % Flèche vers analyse
    \draw[popgreen, thick, ->] (6.5, 1.5) -- (9, 1.5) node[right, popgreen] {Culture cellulaire (\num{10}--\num{15} j)};
    \draw[popgreen, thick, ->] (9, 1.5) -- (11, 1.5) node[right, popgreen] {Caryotype (bande G)};
    
    % Résultat caryotype simplifié
    \begin{scope}[shift={(12,0)}, scale=0.4]
        \draw[popdark] (0,0) rectangle (4,6);
        \foreach \y in {0.5, 1.5, 2.5, 3.5, 4.5, 5.5} {
            \draw[popblue, thick] (0.2,\y) -- (1.8,\y);
            \draw[popblue, thick] (2.2,\y) -- (3.8,\y);
        }
        \node at (2, -0.5) {46,XY ou 47,XY,+21};
    \end{scope}
\end{tikzpicture}
\legende{Principe de l'amniocentèse. Sous contrôle échographique continu, une aiguille traverse la paroi abdominale et utérine pour prélever \SI{15}{\mL} à \SI{20}{\mL} de liquide amniotique contenant des cellules fœtales. Ces cellules sont mises en culture (\num{10} à \num{15} jours) pour obtenir des métaphases, permettant l'établissement du caryotype fœtal (détection trisomie 21, anomalies structurales, sexe).}
\end{popfigure}

\begin{experience}[Établissement du caryotype fœtal par amniocentèse (démarche d'investigation)]
\textbf{Question :} Comment le caryotype obtenu après amniocentèse permet-il de confirmer ou d'infirmer une trisomie 21 suspectée au dépistage ?

\textbf{Protocole (simplifié) :}
\begin{enumerate}
    \item \textbf{Prélèvement :} \SI{15}{\mL}--\SI{20}{\mL} de liquide amniotique (\num{15}--\num{18} SA).
    \item \textbf{Centrifugation :} Pélletage des cellules fœtales (amniocytes).
    \item \textbf{Mise en culture :} Milieu riche (F10 + sérum fœtal), \SI{37}{\celsius}, \SI{5}{\%} CO$_2$, $\approx$ \num{10}--\num{15} jours pour obtenir des mitoses.
    \item \textbf{Blocage en métaphase :} Colchicine (inhibe fuseau mitotique) $\rightarrow$ chromosomes condensés.
    \item \textbf{Lysis hypotonique :} Gonflement cellules $\rightarrow$ écartement chromosomes.
    \item \textbf{Fixation :} Méthanol/Acétique (\num{3}/\num{1}).
    \item \textbf{Étalement sur lame + Coloration bande G (Trypsine-Giemsa) :} Révèle bandes claires (ADN riche AT, réplication tardive) et foncées (ADN riche GC, réplication précoce).
    \item \textbf{Analyse microscopique / Automatique :} Photographie de \num{15}--\num{20} métaphases $\rightarrow$ caryotypage (logiciel ou manuel).
\end{enumerate}

\textbf{Observation / Résultat attendu :}
\begin{itemize}
    \item \textbf{Caryotype normal masculin :} \num{46} chromosomes, paires 1 à 22 homologues + XY. Notation : \textbf{46,XY}.
    \item \textbf{Trisomie 21 libre :} \num{47} chromosomes, \textbf{trois} chromosomes 21 (chromosomes acrocentriques, groupe G). Notation : \textbf{47,XY,+21} (ou 47,XX,+21).
    \item \textbf{Translocation robertsonienne 14;21 :} \num{46} chromosomes mais un chromosome 21 fusionné au 14. Notation : \textbf{46,XY,der(14;21)(q10;q10),+21}. Risque de récidive différent (transmissible par un parent porteur sain).
\end{itemize}

\textbf{Interprétation :} La présence d'un chromosome 21 en surnombre (trisomie libre) ou en excès de matériel génétique 21 (translocation) confirme le diagnostic de syndrome de Down. L'absence d'anomalie numérale ou structurale majeure sur les 22 paires d'autosomes et les gonosomes rassure (mais n'exclut pas les microdélétions ou mutations ponctuelles non visibles au caryotype standard $\rightarrow$ puce à ADN / séquençage si indication).
\end{experience}

\paragraph{Autres techniques de diagnostic prénatal invasif.}
\begin{itemize}
    \item \textbf{Choriocentèse (Prélèvement de villosités choriales - PVC) :} \num{11}--\num{13} SA. Prélèvement de trophoblaste (placenta). Avantage : résultat plus précoce (culture courte \num{24}--\num{48}h ou direct). Risque fausse couche légèrement supérieur ($\approx$ \SI{1}{\%}--\SI{2}{\%}). Risque de mosaïcisme placentaire (faux positif/négatif) $\rightarrow$ confirmation par amniocentèse parfois nécessaire.
    \item \textbf{Cordocentèse (Puncture funiculaire) :} \num{18}--\num{20} SA. Prélèvement sang fœtal au cordon. Indication : caryotype rapide (\num{72}h), hématologie fœtale (anémie, alloimmunisation), infectiologie. Risque plus élevé.
    \item \textbf{ADN fœtal circulant (cfDNA / NIPT - Test Prénatal Non Invasif) :} Dosage de l'ADN fœtal (originaire du trophoblaste) dans le sang maternel (dès \num{10} SA). Dépistage \textbf{très performant} pour T21 (Sensibilité $> \SI{99}{\%}$, Spécificité $> \SI{99,9}{\%}$), T18, T13, aneuploïdies gonosomiques. \textbf{N'est pas un diagnostic} (faux positifs possibles $\rightarrow$ confirmation par amniocentèse obligatoire avant toute décision).
\end{itemize}

\courssub{Anomalies chromosomiques : focus sur la trisomie 21}

\begin{propriete}[Origine chromosomique de la trisomie 21 (Rappel de Première - Savoir admis)]
La trisomie 21 (syndrome de Down) résulte dans $\approx$ \SI{95}{\%} des cas d'une \textbf{non-disjonction chromosomique} lors de la méiose maternelle (majoritairement Méiose I, plus rarement Méiose II) ou paternelle ($\approx$ \SI{5}{\%}).
\begin{itemize}
    \item \textbf{Non-disjonction en Méiose I :} Les chromosomes 21 homologues ne se séparent pas $\rightarrow$ gamète avec 2 chr21 (\num{24} chr) + gamète sans chr21 (\num{22} chr). Fécondation par gamète normal (\num{23} chr) $\rightarrow$ zygote \num{47} chr (trisomie 21) ou \num{45} chr (monosomie 21 $\rightarrow$ létale précoce).
    \item \textbf{Non-disjonction en Méiose II :} Les chromatides sœurs d'un chromosome 21 ne se séparent pas.
\end{itemize}
Le risque de non-disjonction \textbf{augmente significativement avec l'âge maternel} (ex : 1/1500 à 20 ans, 1/350 à 35 ans, 1/100 à 40 ans, 1/30 à 45 ans). C'est le principal indicateur pour proposer un diagnostic prénatal invasif (amniocentèse) ou un NIPT.
\end{propriete}

\begin{exoresolu}[Conseil génétique prénuptial et décision prénatale]
\textbf{Énoncé :} Un couple consulte pour un examen prénuptial au Centre de Dépistage de la Drépanocytose de Yaoundé. L'électrophorèse de l'hémoglobine montre : Homme : profil AS ; Femme : profil AS. Le groupe sanguin Rhésus : Homme Rh+ ($Dd$), Femme Rh- ($dd$). Ils souhaitent avoir des enfants.
\begin{enumerate}
    \item Préciser le risque drépanocytaire pour chaque grossesse.
    \item Préciser le risque Rhésus et la conduite à tenir.
    \item La femme tombe enceinte. Au 1er trimestre, le dépistage combiné (marqueurs sériques + clarté nucale) donne un risque de trisomie 21 de 1/150. Que propose-t-on ? Quelles sont les options du couple ?
\end{enumerate}

\textbf{Solution détaillée :}
\begin{enumerate}
    \item \textbf{Risque drépanocytaire :} Croisement $AS \times AS$.
    \begin{itemize}
        \item Gamètes : A ou S (1/2 chacun) pour chaque parent.
        \item Zygotes possibles : $AA$ (1/4), $AS$ (1/2), $SS$ (1/4).
        \item \textbf{Risque de \SI{25}{\%} (1/4) d'enfant malade (SS) à chaque grossesse.} Risque de \SI{50}{\%} de porteur sain (AS), \SI{25}{\%} de non porteur (AA).
    \end{itemize}
    \item \textbf{Risque Rhésus :} Mère Rh- ($dd$) $\times$ Père Rh+ ($Dd$ ou $DD$). Si père $Dd$ : \SI{50}{\%} fœtus Rh+ ($Dd$), \SI{50}{\%} Rh- ($dd$). Si père $DD$ : \SI{100}{\%} fœtus Rh+.
    \begin{itemize}
        \item \textbf{Risque d'alloinmunisation anti-D} si fœtus Rh+.
        \item \textbf{Conduite à tenir :} Injection d'Ig anti-D à \num{28}-\num{30} SA (prophylaxie anténatale) + dans les \num{72}h post-partum si nouveau-né Rh+. Recherche d'anticorps irréguliers (RAI) mensuelle pendant la grossesse.
    \end{itemize}
    \item \textbf{Risque Trisomie 21 :} 1/150 $>$ seuil de 1/250 $\rightarrow$ \textbf{Risque élevé}.
    \begin{itemize}
        \item \textbf{Proposition :} Diagnostic prénatal invasif : \textbf{Amniocentèse} (\num{15}--\num{18} SA) pour caryotype (résultat certain) \textbf{OU} Choriocentèse (\num{11}--\num{13} SA) pour résultat plus rapide.
        \item \textbf{Alternative (si refus invasif) :} NIPT (ADN fœtal circulant) pour affiner le risque (mais pas diagnostique).
        \item \textbf{Options du couple après résultat :}
        \begin{itemize}
            \item Caryotype normal (46,XX ou 46,XY) $\rightarrow$ Poursuite grossesse rassurée.
            \item Trisomie 21 confirmée (47,XX,+21 ou 47,XY,+21) $\rightarrow$ Information complète sur le syndrome (handicap intellectuel variable, cardiopathies, espérance de vie $\approx$ \num{60} ans aujourd'hui). \textbf{Choix parental :} Poursuite de la grossesse avec préparation pédiatrique spécialisée \textbf{OU} Demande d'Interruption Médicale de Grossesse (IMG) si la loi le permet (au Cameroun, l'IMG est autorisée si la grossesse met en danger la vie de la mère ; le handicap fœtal seul ne constitue pas un motif légal d'IMG au sens strict du Code Pénal, mais la pratique hospitalière évolue). Le couple est libre de sa décision, accompagné par l'équipe pluridisciplinaire.
        \end{itemize}
    \end{itemize}
\end{enumerate}
\end{exoresolu}

\courssub{Limites, éthique et bilan}

\begin{attention}[Limites techniques et éthiques du diagnostic prénatal]
\begin{itemize}
    \item \textbf{Risque fœtal :} Amniocentèse / Choriocentèse : risque de fausse couche $\approx$ \SI{0,5}{\%}--\SI{1}{\%}. À peser face au risque de la maladie recherchée.
    \item \textbf{Faux positifs / négatifs :} NIPT : faux positifs possibles (mosaïcisme placentaire, tumeur maternelle, jumelle disparue). \textbf{Confirmation par caryotype invasif obligatoire avant toute IMG.}
    \item \textbf{Portée du caryotype standard :} Ne détecte que les anomalies \textbf{numériques} (anéuploïdies) et \textbf{structurales majeures} ($> 5$--$10$ Mpb). Ne voit pas les mutations ponctuelles (ex : drépanocytose, myopathies) $\rightarrow$ nécessite biologie moléculaire ciblée (PCR, séquençage) sur l'ADN amniotique si indication (couple à risque connu).
    \item \textbf{Secret médical et autonomie :} Les résultats relèvent du \textbf{secret médical} partagé entre les deux membres du couple (ou la femme seule si elle le souhaite). Le médecin \emph{informe}, ne \emph{décide pas}. L'Interruption Médicale de Grossesse (IMG) est une décision médicale collégiale validée par un comité pluridisciplinaire, sur demande des parents, dans le cadre légal strict.
    \item \textbf{Eugénisme vs Médecine préventive :} La ligne est fine. Le but est la \textbf{prévention de la souffrance} (maladies graves, incurables, d'apparition précoce) et l'\textbf{autonomie reproductive}, non la « sélection » de traits non pathologiques.
\end{itemize}
\end{attention}

\begin{aretenir}[Bilan : La génétique au service de la santé publique]
La connaissance des modes de transmission (autosomique récessif/dominant, gonosomique, chromosomique) permet :
\begin{enumerate}
    \item \textbf{En amont (Prénuptial) :} D'identifier les couples à risque (ex: $AS \times AS$, Femme Rh- $\times$ Homme Rh+) $\rightarrow$ \textbf{Information} et \textbf{Conseil} pour une procréation éclairée (diagnostic prénatal, FIV-DPI, adoption, don de gamètes).
    \item \textbf{Pendant la grossesse (Prénatal) :} De détecter les anomalies chromosomiques (trisomie 21 via caryotype amniotique) ou génétiques (maladies mongéniques via ADN amniotique) $\rightarrow$ \textbf{Prise en charge adaptée} (naissance en centre spécialisé, chirurgie néonatale précoce) ou \textbf{IMG} si pathologie grave et incurable (cadre légal).
    \item \textbf{A l'échelle populationnelle :} Au Cameroun, le dépistage néonatal de la drépanocytose (test de Guthrie / électrophorèse sur sang de cordon ou talon à J3) couplé à la prise en charge précoce (pénicilline prophylactique, vaccins, hydroxyurée, éducation thérapeutique) a \textbf{réduit la mortalité infanto-juvénile} liée à la drépanocytose. La généralisation du certificat prénuptial participe à la \textbf{limitation de la fréquence} de l'allèle $S$ dans les générations futures par choix reproductifs informés.
\end{enumerate}
\end{aretenir}

\begin{exercice}[Application : Conseil génétique complet]
Un couple, M. et Mme N., se présente en consultation de génétique au CHU de Douala avant un projet de grossesse.
\begin{itemize}
    \item M. N. est drépanocytaire majeur (SS). Il est daltonien (récessif lié à l'X). Groupe sanguin A Rh+.
    \item Mme N. est phénotypiquement normale pour la drépanocytose et la vision des couleurs. Groupe sanguin O Rh-.
    \item Ils ont déjà un garçon de 3 ans, drépanocytaire (SS), non daltonien, groupe A Rh+.
\end{itemize}
\begin{enumerate}
    \item Déterminer le génotype probable de Mme N. pour le gène $\beta$-globine et pour le gène de la vision des couleurs. Justifier.
    \item Calculer le risque pour une future grossesse d'avoir : (a) un enfant drépanocytaire ; (b) un garçon daltonien ; (c) une fille daltonienne ; (d) un enfant de groupe O Rh-.
    \item Quel examen prénatal proposeriez-vous spécifiquement pour la drépanocytose ? Sur quel prélèvement ? À quel terme ?
    \item Quel suivi Rhésus mettre en place pour Mme N. ?
\end{enumerate}
\end{exercice}

\begin{corrige}[Exercice : Conseil génétique complet]
\begin{enumerate}
    \item \textbf{Génotype Mme N. :}
    \begin{itemize}
        \item \textbf{Drépanocytose :} Enfant SS $\rightarrow$ a reçu $S$ du père (SS) et $S$ de la mère. Donc Mme N. est obligatoirement \textbf{AS} (porteuse).
        \item \textbf{Daltonisme (récessif X) :} Père daltonien ($X^dY$). Fils non daltonien $\rightarrow$ a reçu $Y$ du père et $X^D$ de la mère. La mère a donné $X^D$. Comme elle est phénotypiquement normale, elle est soit $X^D X^D$, soit $X^D X^d$. \textbf{Indéterminé} (mais probabilité porteuse selon fréquence allélique).
        \item \textbf{Groupes sanguins :} Père A Rh+ (génotype possible $I^A i$ ou $I^A I^A$ ; $DD$ ou $Dd$). Mère O Rh- ($ii$, $dd$). Enfant A Rh+ $\rightarrow$ a reçu $I^A$ du père et $i$ de la mère $\rightarrow$ Père $I^A i$. A reçu $D$ du père $\rightarrow$ Père $Dd$.
    \end{itemize}
    \item \textbf{Risques par grossesse :}
    \begin{itemize}
        \item (a) Drépanocytose : $SS \times AS \rightarrow$ 1/2 $SS$ (malade), 1/2 $AS$ (porteur). \textbf{Risque = \SI{50}{\%}}.
        \item (b) Garçon daltonien : Mère $X^D X^?$ (inconnu). Si mère $X^D X^D$ : \SI{0}{\%}. Si mère $X^D X^d$ : 1/2 fils $X^d Y$ (daltonien). Risque global = $P(\text{mère porteuse}) \times 1/2$. Sans info fréquence, on note le risque conditionnel.
        \item (c) Fille daltonienne : Père $X^d Y$ donne $X^d$ à toutes ses filles. Mère donne $X^D$ ou $X^d$. Fille daltonienne = $X^d X^d$. Nécessite mère porteuse ($X^D X^d$) ET transmission $X^d$ (1/2). Risque = $P(\text{mère porteuse}) \times 1/2$.
        \item (d) Groupe O Rh- : Père $I^A i, Dd$. Mère $ii, dd$. Enfant O = $ii$ (1/2). Enfant Rh- = $dd$ (1/2). Indépendants $\rightarrow$ 1/4 = \textbf{\SI{25}{\%}}.
    \end{itemize}
    \item \textbf{Diagnostic prénatal drépanocytose :} Recherche de la mutation $\beta^S$ (Glu6Val) par \textbf{PCR / RFLP / Séquençage} sur \textbf{ADN extrait des amniocytes} (liquide amniotique) prélevé par \textbf{amniocentèse à \num{15}--\num{18} SA}. (Ou sur trophoblaste par choriocentèse \num{11}--\num{13} SA). Résultat rapide (quelques jours).
    \item \textbf{Suivi Rhésus :} Mme N. Rh- ($dd$), conjoint $Dd$. Risque fœtus Rh+ = \SI{50}{\%}.
    \begin{itemize}
        \item RAI (Recherche Anticorps Irréguliers) mensuelle.
        \item Injection Ig anti-D (\SI{300}{\micro\gram}) à \num{28}--\num{30} SA.
        \item Injection Ig anti-D dans \num{72}h post-partum si enfant Rh+ (confirmé sur sang cordon).
        \item Injection Ig anti-D si amniocentèse / saignement / traumatisme.
    \end{itemize}
\end{enumerate}
\end{corrige}

\newpage

\coursec{Activité d'intégration}

\courssub{Objectifs de l'activité}

À l'issue de cette activité d'intégration, tu seras capable de :
\begin{itemize}
\item construire un \cle{arbre généalogique} (pedigree) à partir d'un récit familial, en respectant les symboles conventionnels ;
\item déterminer le \cle{mode de transmission} (autosomique ou gonosomique, récessif ou dominant) d'une affection à partir d'un pedigree et de documents expérimentaux (électrophorèse, caryotype) ;
\item évaluer un \cle{risque génétique} pour la descendance d'un couple en croisant les informations de deux familles ;
\item rédiger un avis de \cle{conseil génétique} argumenté, en précisant les \cle{examens prénuptiaux et prénataux} recommandés.
\end{itemize}

\courssub{Situation}

\textbf{Conseil génétique : le cas de la famille Ngo Bell (Douala).}

M.~Etienne Ngo Bell, 28 ans, enseignant à Douala, souhaite épouser Mlle~Sandrine Mbappe, 26 ans, infirmière au Centre Hospitalier d'Essos. Avant le mariage, et sur les conseils du médecin du centre de santé, les deux familles acceptent de reconstituer leurs histoires médicales et de réaliser quelques examens.

\textbf{Document 1 — Récit de la famille Ngo Bell (famille du fiancé).}
Les grands-parents paternels d'Etienne, M.~et Mme~Ngo Bell (génération I), sont tous les deux sains et à peau normalement pigmentée. Ils ont eu quatre enfants : deux garçons sains (dont le père d'Etienne), une fille saine, et une fille \textbf{albinos} (dépigmentation de la peau, des cheveux et des iris). Le père d'Etienne, sain, a épousé une femme saine sans aucun cas d'albinisme dans sa famille ; de cette union sont nés Etienne et sa sœur, tous deux sains. Aucun cas de drépanocytose ni de daltonisme n'est connu dans cette famille.

\textbf{Document 2 — Récit de la famille Mbappe (famille de la fiancée).}
Le frère cadet de Sandrine est \textbf{drépanocytaire} (forme majeure SS, anémie sévère, crises douloureuses fréquentes, suivies à l'hôpital) ; ses parents et Sandrine sont cliniquement sains. Par ailleurs, le père de Sandrine est \textbf{daltonien} (il confond le rouge et le vert), alors que sa mère a une vision normale des couleurs ; on apprend aussi que le grand-père maternel de Sandrine était daltonien. Sandrine et son frère drépanocytaire ont une vision des couleurs normale.

\textbf{Document 3 — Électrophorèse de l'hémoglobine des deux fiancés.}
L'électrophorèse réalisée au laboratoire donne deux bandes (HbA et HbS) pour Etienne comme pour Sandrine : les deux fiancés sont de génotype \textbf{AS} (porteurs sains du trait drépanocytaire).

\textbf{Document 4 — Caryotypes.}
Les caryotypes des deux fiancés sont normaux : 46,XY pour Etienne et 46,XX pour Sandrine, sans anomalie de structure ni de nombre. Aucune anomalie chromosomique n'est donc à craindre dans ce couple ; l'étude portera uniquement sur les maladies \textbf{géniques}.

On rappelle : l'\cle{albinisme} est dû à l'allèle récessif $a$ (allèle normal $A$) ; la \cle{drépanocytose} met en jeu les allèles codominants $A$ (HbA) et $S$ (HbS), les hétérozygotes AS étant sains ; le \cle{daltonisme} est dû à l'allèle récessif $d$ porté par le chromosome $X$ (allèle normal $D$).

\courssub{Consignes}

\begin{enumerate}
\item Construis l'arbre généalogique de la famille Ngo Bell (trois générations) en utilisant les symboles conventionnels du pedigree.
\item À partir de ce pedigree, détermine le mode de transmission de l'albinisme (autosomique ou gonosomique ? dominant ou récessif ?) en justifiant par deux arguments tirés de l'arbre.
\item Construis l'arbre généalogique de la famille Mbappe pour la drépanocytose et pour le daltonisme.
\item Détermine le mode de transmission du daltonisme et déduis-en le génotype de Sandrine pour ce caractère (justifie à partir du phénotype de son père).
\item Calcule le risque, pour un enfant du couple Etienne $\times$ Sandrine :
\begin{enumerate}
\item d'être drépanocytaire SS ;
\item d'être daltonien, en distinguant garçons et filles ;
\item d'être albinos, en précisant les hypothèses nécessaires sur le génotype de Sandrine.
\end{enumerate}
\item Rédige, en 15 à 20 lignes, un avis de conseil génétique argumenté à l'intention du couple : rappel des risques chiffrés, examens prénuptiaux déjà réalisés ou à compléter, examens prénataux envisageables pendant la grossesse, et attitude à adopter.
\end{enumerate}

\courssub{Ressources à mobiliser}

\begin{itemize}
\item \textbf{Symboles du pedigree} : cercle = femme, carré = homme ; symbole rempli = individu atteint ; trait horizontal = union ; trait vertical = descendance ; générations numérotées I, II, III.
\item \textbf{Règles de diagnostic du mode de transmission} : des parents sains ayant un enfant atteint $\Rightarrow$ allèle \textbf{récessif} ; un père atteint transmettant à toutes ses filles (ou un caractère touchant surtout les garçons) $\Rightarrow$ transmission \textbf{gonosomique liée à $X$} ; un caractère présent à chaque génération $\Rightarrow$ transmission \textbf{dominante}.
\item \textbf{Échiquiers de croisement} et probabilités : chez les hétérozygotes $A/a \times A/a$, la descendance est $1/4$ $AA$, $1/2$ $A/a$, $1/4$ $a/a$ ; pour un caractère lié à $X$, le père transmet $X$ à ses filles et $Y$ à ses fils.
\item \textbf{Probabilité d'être porteur sain} : un individu sain issu de deux parents hétérozygotes a une probabilité de $2/3$ d'être hétérozygote.
\item \textbf{Examens} : électrophorèse de l'hémoglobine et test de compatibilité (prénuptiaux) ; amniocentèse, prélèvement de villosités choriales, échographie (prénataux).
\end{itemize}

\courssub{Démarche et corrigé commenté}

\begin{exoresolu}[Conseil génétique : le cas de la famille Ngo Bell]
Énoncé : voir la situation ci-dessus (documents 1 à 4). Il s'agit de construire les pedigrees, de déterminer les modes de transmission, d'évaluer les risques pour la descendance du couple et de formuler un avis de conseil génétique.
\tcblower
\textbf{Étape 1 — Pedigree de la famille Ngo Bell (albinisme).}

\begin{popfigure}
\begin{center}
\begin{tikzpicture}[
    scale=0.9,
    every node/.style={font=\small, inner sep=1.5pt},
    male/.style={draw, rectangle, minimum size=0.6cm},
    female/.style={draw, circle, minimum size=0.6cm},
    affected/.style={fill=popink},
    genlabel/.style={font=\bfseries\small, anchor=east, xshift=-1cm}
]
% Generation I
\node[genlabel] (GI) at (0,0) {I};
\node[male] (I1) at (1,0) {};
\node[female] (I2) at (2.5,0) {};
\draw (I1) -- (I2);
% Children line Gen I
\draw (1.75,0) -- (1.75,-0.5);
\draw (0.5,-0.5) -- (3.5,-0.5);
% Generation II
\node[genlabel] (GII) at (0,-1.2) {II};
\node[male] (II1) at (0.5,-1.2) {};
\node[male, label=below:{\footnotesize Père d'Etienne}] (II2) at (1.75,-1.2) {};
\node[female] (II3) at (3.0,-1.2) {};
\node[female, affected, label=below:{\footnotesize Tante albinos}] (II4) at (3.5,-1.2) {};
\draw (0.5,-0.5) -- (II1);
\draw (1.75,-0.5) -- (II2);
\draw (3.0,-0.5) -- (II3);
\draw (3.5,-0.5) -- (II4);
% Spouse of II2
\node[female] (II2w) at (1.75,-2.4) {};
\draw (II2) -- (II2w);
% Children line Gen II
\draw (1.75,-2.4) -- (1.75,-2.9);
\draw (1.0,-2.9) -- (2.5,-2.9);
% Generation III
\node[genlabel] (GIII) at (0,-3.6) {III};
\node[male, label=below:{\footnotesize Etienne}] (III1) at (1.0,-3.6) {};
\node[female] (III2) at (2.5,-3.6) {};
\draw (1.0,-2.9) -- (III1);
\draw (2.5,-2.9) -- (III2);
% Legend symbols
\node[male, xshift=5cm, yshift=1cm] (legM) at (0,0) {};
\node[female, right=0.5cm of legM] (legF) {};
\node[female, affected, right=0.5cm of legF] (legA) {};
\node[right=0.2cm of legM] {Homme sain};
\node[right=0.2cm of legF] {Femme saine};
\node[right=0.2cm of legA] {Individu albinos};
\end{tikzpicture}
\end{center}
\legende{Pedigree de la famille Ngo Bell (albinisme). Carré = homme ; cercle = femme ; symbole plein = individu atteint (albinos) ; trait horizontal = union ; trait vertical = descendance. Générations I (grands-parents), II (père, oncles, tantes), III (Etienne et sa sœur).}
\end{popfigure}

\textbf{Étape 2 — Mode de transmission de l'albinisme.}
Deux arguments tirés du pedigree permettent de conclure à une transmission \textbf{autosomique récessive} :
\begin{enumerate}
\item \textbf{Argument de récessivité :} Les grands-parents (génération I) sont phénotypiquement sains mais ont engendré une fille albinos (génération II). L'allèle responsable est donc masqué à l'état hétérozygote chez les parents : c'est un allèle \textbf{récessif} (noté $a$). Les grands-parents sont obligatoirement hétérozygotes $A/a$.
\item \textbf{Argument d'autosomie :} L'individu atteint de la génération II est une \textbf{fille}. Si l'allèle était porté par le chromosome $X$ (transmission gonosomique récessive), une fille atteinte ($X^a X^a$) aurait obligatoirement un père atteint ($X^a Y$), or le grand-père est sain. La transmission est donc \textbf{autosomique}.
\end{enumerate}

\textbf{Étape 3 — Pedigrees de la famille Mbappe.}

\begin{popfigure}
\begin{center}
\begin{tikzpicture}[
    scale=0.85,
    every node/.style={font=\small, inner sep=1.5pt},
    male/.style={draw, rectangle, minimum size=0.55cm},
    female/.style={draw, circle, minimum size=0.55cm},
    affectedS/.style={fill=poporange}, % Drépanocytose
    affectedD/.style={fill=popblue},   % Daltonisme
    carrierD/.style={draw, circle, pattern=north east lines, pattern color=popblue}, % Conductrice
    genlabel/.style={font=\bfseries\small, anchor=east, xshift=-1.2cm}
]
% Gen I (Maternal grandparents)
\node[genlabel] (GI) at (0,0) {I};
\node[male, affectedD, label=below:{\footnotesize \shortstack{G-père maternel\\daltonien}}] (I1) at (1.5,0) {};
\node[female, label=below:{\footnotesize G-mère maternelle}] (I2) at (3,0) {};
\draw (I1) -- (I2);
\draw (2.25,0) -- (2.25,-0.5);
\draw (1.5,-0.5) -- (3,-0.5);
% Gen II (Parents of Sandrine)
\node[genlabel] (GII) at (0,-1.2) {II};
\node[male, affectedD, label=below:{\footnotesize \shortstack{Père\\daltonien}}] (II1) at (1.5,-1.2) {};
\node[female, carrierD, label=below:{\footnotesize \shortstack{Mère\\conductrice}}] (II2) at (3,-1.2) {};
\draw (1.5,-0.5) -- (II1);
\draw (3,-0.5) -- (II2);
\draw (II1) -- (II2);
% Children line
\draw (2.25,-1.2) -- (2.25,-1.7);
\draw (1.5,-1.7) -- (3,-1.7);
% Gen III (Sandrine & Brother)
\node[genlabel] (GIII) at (0,-2.4) {III};
\node[female, label=below:{\footnotesize \shortstack{Sandrine\\AS, vision normale}}] (III1) at (1.5,-2.4) {};
\node[male, affectedS, label=below:{\footnotesize \shortstack{Frère\\SS, vision normale}}] (III2) at (3,-2.4) {};
\draw (1.5,-1.7) -- (III1);
\draw (3,-1.7) -- (III2);
% Legends
\node[male, affectedD, xshift=5.5cm, yshift=1.5cm] (l1) at (0,0) {};
\node[right=0.2cm of l1] {Homme daltonien};
\node[female, carrierD, right=0.5cm of l1] (l2) {};
\node[right=0.2cm of l2] {Femme conductrice daltonisme};
\node[male, affectedS, below=0.5cm of l1] (l3) {};
\node[right=0.2cm of l3] {Homme drépanocytaire SS};
\node[female, below=0.5cm of l2] (l4) {};
\node[right=0.2cm of l4] {Femme saine (AS ou AA)};
\end{tikzpicture}
\end{center}
\legende{Pedigree de la famille Mbappe. Haut : transmission du daltonisme (lié à X, récessif). Bas : transmission de la drépanocytose (autosomique, codominance). Symboles : carré = homme, cercle = femme ; plein orange = drépanocytaire SS ; plein bleu = daltonien ; hachuré bleu = conductrice daltonisme ; trait horizontal = union.}
\end{popfigure}

\textbf{Étape 4 — Mode de transmission du daltonisme et génotype de Sandrine.}
\begin{itemize}
\item \textbf{Transmission gonosomique liée à $X$ récessive :} Le grand-père maternel (I-1) est daltonien ; sa fille (II-2, mère de Sandrine) est saine mais conductrice (obligée, car elle a reçu le $X^d$ de son père). Le père de Sandrine (II-1) est daltonien ($X^d Y$). L'atteinte ne touche que des hommes (grand-père, père) et se transmet par les femmes conductrices : c'est la signature d'un allèle \textbf{récessif lié au chromosome $X$} (allèle $d$, allèle normal $D$).
\item \textbf{Génotype de Sandrine :} Son père est $X^d Y$. Il transmet son unique chromosome $X$ (porteur de $d$) à \textbf{toutes} ses filles. Sandrine a reçu $X^d$ de son père. Comme elle a une vision normale, elle a obligatoirement reçu un allèle $D$ de sa mère (conductrice $X^D X^d$). Sandrine est donc \textbf{conductrice hétérozygote : $X^D X^d$} (certitude, probabilité 1).
\end{itemize}

\textbf{Étape 5 — Calcul des risques pour les enfants du couple Etienne ($X^D Y$, $A/a$?, $A/S$) $\times$ Sandrine ($X^D X^d$, $A/A$?, $A/S$).}

\emph{a) Risque de drépanocytose majeure (SS).}
L'électrophorèse (Doc 3) montre que les deux fiancés sont \textbf{AS} (hétérozygotes codominants $A/S$).
Croisement : $A/S \times A/S$ $\rightarrow$ $1/4\ A/A$ (sain), $1/2\ A/S$ (porteur sain), $1/4\ S/S$ (drépanocytaire).
\[
\boxed{P(\text{enfant SS}) = \dfrac{1}{4} = 25\% \text{ à chaque grossesse, indépendamment du sexe.}}
\]

\emph{b) Risque de daltonisme.}
Etienne (vision normale) : $X^D Y$. Sandrine (conductrice) : $X^D X^d$.
Croisement : $X^D X^d \times X^D Y$.
\begin{itemize}
\item \textbf{Filles :} Reçoivent $X^D$ du père. Génotypes : $1/2\ X^D X^D$ (saines), $1/2\ X^D X^d$ (conductrices saines). \textbf{Aucune fille daltonienne} ($0\%$).
\item \textbf{Garçons :} Reçoivent $Y$ du père et $X$ de la mère. Génotypes : $1/2\ X^D Y$ (sains), $1/2\ X^d Y$ (daltoniens). \textbf{1 garçon sur 2 daltonien} ($50\%$).
\end{itemize}
Risque global par naissance (sexe non connu) : $P(\text{garçon}) \times P(\text{daltonien}| \text{garçon}) = 1/2 \times 1/2 = 1/4 = 25\%$, mais \textbf{exclusivement chez les garçons}.

\emph{c) Risque d'albinisme.}
\begin{itemize}
\item \textbf{Etienne :} Son père (II-2) est issu du croisement $A/a \times A/a$ (grands-parents obligatoirement hétérozygotes). Son père est sain $\Rightarrow$ Probabilité d'être $A/a = 2/3$. Sa mère (II-2w) sans antécédent familial $\Rightarrow$ Supposée $A/A$. Etienne a donc une probabilité de $2/3 \times 1/2 = \mathbf{1/3}$ d'être porteur $A/a$.
\item \textbf{Sandrine :} Aucun cas d'albinisme signalé dans sa famille (Doc 2). On suppose qu'elle est \textbf{$A/A$} (hypothèse de non-portage en l'absence d'information).
\end{itemize}
Croisement possible : Si Etienne est $A/a$ (prob. 1/3) et Sandrine $A/A$ $\rightarrow$ $0$ albinos.
Si Sandrine était porteuse (probabilité faible, fréquence allélique $q$ dans la population), le risque serait $1/3 \times 1/2 \times 1/4 = 1/24$.
\[
\boxed{\text{Risque d'albinisme \textbf{quasi nul} (hypothèse Sandrine $A/A$). Un test génétique sur Sandrine leverait l'incertitude résiduelle.}}
\]

\textbf{Étape 6 — Avis de conseil génétique (proposition de rédaction).}

\begin{aretenir}[Avis de conseil génétique pour le couple Ngo Bell -- Mbappe]
M. Etienne Ngo Bell et Mlle Sandrine Mbappe, votre demande de conseil génétique prénuptial a porté sur trois affections héréditaires identifiées dans vos familles.

\textbf{1. Risques identifiés.}
\begin{itemize}
\item \textbf{Drépanocytose (forme majeure SS) :} L'électrophorèse d'hémoglobine révèle que vous êtes \textbf{tous deux porteurs sains (AS)}. Le risque d'avoir un enfant drépanocytaire SS est de \textbf{25 \% à chaque grossesse}, quel que soit le sexe. C'est le risque majeur, la drépanocytose étant une affection grave (anémie hémolytique chronique, crises vaso-occlusives, infections répétées) nécessitant une prise en charge médicale spécialisée précoce (suivi en centre drépanocytose, vaccination antipneumococcique, prophylaxie antipaludique, hydratation).
\item \textbf{Daltonisme (récessif lié à $X$) :} Sandrine est conductrice certaine ($X^D X^d$), Etienne est sain ($X^D Y$). Le risque est de \textbf{50 \% pour chaque garçon} (daltonien), 0 \% pour les filles (toutes saines, 50 \% conductrices). Affection bénigne, sans conséquence vitale, mais à connaître pour l'orientation professionnelle.
\item \textbf{Albinisme (autosomique récessif) :} Etienne a 1 chance sur 3 d'être porteur ; Sandrine, sans antécédent, est supposée non porteuse. Le risque est \textbf{quasi nul}. Un dépistage du portage chez Sandrine par biologie moléculaire peut être proposé si le couple le souhaite.
\end{itemize}

\textbf{2. Examens réalisés et recommandations.}
Les examens prénuptiaux (électrophorèse d'hémoglobine, groupage ABO/Rhésus, caryotypes) ont permis cette évaluation. Nous recommandons de compléter par une \textbf{consultation spécialisée en génétique médicale} avant toute grossesse.

\textbf{3. Diagnostic prénatal (DPN).}
Au cours de chaque grossesse, un \textbf{diagnostic prénatal} peut être proposé pour la drépanocytose :
\begin{itemize}
\item \textbf{Prélèvement de villosités choriales (PVC)} vers 11--13 SA (semaines d'aménorrhée) : analyse d'ADN fœtal (recherche mutation $\beta^S$).
\item \textbf{Amniocentèse} vers 15--17 SA : culture de cellules amniotiques + analyse ADN.
\end{itemize}
Ces actes invasifs comportent un risque de fausse couche ($\approx 0,5$ à $1\%$). Ils permettent au couple, en toute liberté, de préparer la naissance d'un enfant SS (organisation du suivi néonatal immédiat) ou, conformément à la loi camerounaise sur la santé de la reproduction, de demander une interruption médicale de grossesse (IMG) si le fœtus est atteint.

\textbf{4. Attitude à adopter.}
Le conseil génétique n'interdit pas l'union : il \textbf{informe} pour un choix éclairé. En cas de grossesse, un suivi obstétrical rapproché en maternité de niveau 2 ou 3 (ex: CHU de Douala, Hôpital Laquintinie) est indispensable. À la naissance, le \textbf{dépistage néonatal systématique de la drépanocytose} (test de Guthrie / électrophorèse sur sang de cordon) permettra une prise en charge précoce (pénicilline prophylactique dès 3 mois, vaccins, éducation des parents), améliorant drastiquement la survie et la qualité de vie des enfants SS.
\end{aretenir}

\end{exoresolu}

\courssub{Bilan de l'activité}

Cette activité t'a permis d'intégrer l'ensemble des savoir-faire de la leçon dans une démarche authentique de \cle{conseil génétique}. Tu as d'abord \textbf{construit des pedigrees} à partir de récits familiaux, en appliquant les symboles conventionnels (carré/cercle, plein/vide, unions, générations). Tu as ensuite \textbf{diagnostiqué les modes de transmission} en mobilisant les règles d'argumentation rigoureuses :
\begin{itemize}
\item Caractère \textbf{récessif} révélé par des parents sains ayant un enfant atteint (albinisme, drépanocytose).
\item Caractère \textbf{lié au chromosome $X$} révélé par la transmission du père à toutes ses filles (conductrices obligées) et l'atteinte préférentielle des garçons (daltonisme).
\item \textbf{Codominance} (système ABO, drépanocytose) mise en évidence par l'électrophorèse (phénotype AS = deux bandes).
\end{itemize}
Tu as enfin \textbf{évalué des risques chiffrés} (25 \% drépanocytose SS, 50 \% daltonisme chez les garçons, risque albinisme quasi nul) et compris que ces probabilités s'appliquent \emph{à chaque grossesse, indépendamment des naissances précédentes} (indépendance des événements). Sur le plan du savoir-être, tu as constaté que le conseil génétique ne s'oppose jamais à une union : il éclaire le choix du couple et oriente vers les \cle{examens prénuptiaux} (électrophorèse, groupage, caryotype) et \cle{prénataux} (PVC, amniocentèse, diagnostic moléculaire), outils essentiels de prévention des maladies héréditaires comme la drépanocytose, dont la prévalence du trait porteur (AS) atteint \textbf{20 à 25 \% de la population au Cameroun}, justifiant pleinement la politique nationale de dépistage néonatal et de conseil génétique.

\newpage

\coursec{Méthodes et exercices résolus}

\courssub{Méthodes et exercices résolus}

\begin{methode}[Construire un arbre généalogique (pedigree) à partir d'un énoncé]
    \tikzset{
        male/.style={circle, minimum size=0.7cm, thick, draw=popdark},
        female/.style={rectangle, minimum size=0.7cm, thick, draw=popdark},
        affected/.style={fill=popdark},
        unaffected/.style={fill=white},
        carrier/.style={pattern=north east lines, pattern color=popdark},
        deceased/.style={path picture={\draw[thick, popdark] (path picture bounding box.south west) -- (path picture bounding box.north east);}}
    }
    \textbf{Objectif :} Représenter graphiquement l'histoire d'une famille pour y suivre un caractère héréditaire.
    \begin{enumerate}
        \item \textbf{Identifier les individus :} Numéroter les générations en chiffres romains (I, II, III...) et les individus dans chaque génération en chiffres arabes (1, 2, 3... de gauche à droite). Noter I-1, II-3, etc.
        \item \textbf{Utiliser les symboles conventionnels (standard international) :}
              \begin{center}
                  \begin{tabular}{ccl}
                      \tikz[baseline=-0.5ex]\node[male, unaffected] {}; & \hspace{2mm} & Homme sain (non atteint) \\
                      \tikz[baseline=-0.5ex]\node[male, affected] {}; & & Homme atteint \\
                      \tikz[baseline=-0.5ex]\node[female, unaffected] {}; & & Femme saine (non atteinte) \\
                      \tikz[baseline=-0.5ex]\node[female, affected] {}; & & Femme atteinte \\
                      \tikz[baseline=-0.5ex]\node[diamond, unaffected, minimum size=0.6cm] {}; & & Sexe indéterminé / fœtus \\
                      \tikz[baseline=-0.5ex]\node[male, unaffected] {\textcolor{popmuted}{?}}; & & Phénotype inconnu
                  \end{tabular}
              \end{center}
              \textit{Portage (hétérozygotie récessive) :} symbole mi-rempli (hachuré) ou génotype noté sous le symbole. \textit{Décédé :} barre oblique \tikz[baseline=-0.5ex]\draw[popdark, thick] (-0.15,-0.15) -- (0.15,0.15); sur le symbole. \textit{Jumeaux monozygotes :} trait diagonal reliant les symboles. \textit{Jumeaux dizygotes :} symboles côte à côte sur même ligne de fratrie.
        \item \textbf{Relier les individus :}
              \begin{itemize}
                  \item Trait horizontal \tikz[baseline=-0.5ex]\draw[popdark, thick] (0,0) -- (0.5,0); = union (mariage/accouplement). Double trait pour consanguinité.
                  \item Trait vertical descendant depuis le milieu de l'union = filiation.
                  \item Fratrie : trait horizontal reliant le bas des traits verticaux des enfants.
              \end{itemize}
        \item \textbf{Annoter :} Indiquer le génotype (si connu ou déduit) sous le symbole, l'âge, le statut (décédé), les jumeaux.
        \item \textbf{Vérifier :} Tous les individus cités sont présents ? Les liens de parenté sont respectés ? La cohérence génétique est-elle validée ?
    \end{enumerate}
\end{methode}

\begin{exoresolu}[Construction d'un pedigree : Drépanocytose dans une famille de Yaoundé]
    \textbf{Énoncé :} Dans une famille, le grand-père paternel (I-1) est drépanocytaire (SS) et la grand-mère paternelle (I-2) est saine homozygote (AA). Ils ont eu trois enfants : un garçon (II-1) décédé jeune, une fille (II-2) mariée à un homme sain (II-3) dont ils ont une fille saine (III-1) et un garçon drépanocytaire (III-2), et un dernier garçon (II-4) marié à une femme hétérozygote (II-5) dont ils attendent un enfant (III-3). Construire l'arbre généalogique de cette famille en faisant apparaître les génotypes.
    \tcblower
    \textbf{Solution détaillée :}
    \begin{enumerate}
        \item \textbf{Analyse génétique préalable (déduction des génotypes) :}
              \begin{itemize}
                  \item \textbf{Gén I :} I-1 (SS) $\times$ I-2 (AA) $\rightarrow$ 100\% hétérozygotes AS (porteurs sains). Donc II-1, II-2, II-4 sont \textbf{AS}.
                  \item \textbf{Famille II-2 $\times$ II-3 :} II-2 (AS) $\times$ II-3 (phénotype sain). Ils ont un fils III-2 atteint (SS).
                  \item \textbf{Piège / Incohérence :} Un parent sain homozygote (AA) ne peut transmettre l'allèle $S$. Or III-2 est SS. \textbf{Conclusion : II-3 n'est pas AA, il est obligatoirement AS (porteur sain).} L'énoncé « homme sain » est compatible avec AS.
                  \item \textbf{Famille II-4 $\times$ II-5 :} II-4 (AS) $\times$ II-5 (AS) $\rightarrow$ Croisement $AS \times AS$. Risque pour III-3 : 1/4 AA, 1/2 AS, 1/4 SS.
              \end{itemize}
        \item \textbf{Construction du pedigree (TikZ) :}
    \end{enumerate}
    \begin{popfigure}
        \begin{tikzpicture}[scale=0.75, thick, popdark,
            male/.style={circle, minimum size=0.8cm, draw, fill=white},
            female/.style={rectangle, minimum size=0.8cm, draw, fill=white},
            affected/.style={fill=poporange},
            carrier/.style={pattern=north east lines, pattern color=popdark},
            deceased/.style={path picture={\draw[thick] (path picture bounding box.south west) -- (path picture bounding box.north east);}}
        ]
            % Styles nodes
            \tikzset{person/.style={minimum size=0.8cm, thick, draw=popdark, inner sep=2pt, font=\small}}
            
            % Generation I (y=3)
            \node[person, circle, affected, align=center] (I1) at (0,3){I-1\\SS};
            \node[person, rectangle, align=center] (I2) at (3,3){I-2\\AA};
            \draw (I1) -- (I2) node[midway, above] {Union};
            \draw (1.5,3) -- (1.5,2.2);
            
            % Generation II (y=1.5)
            % Sibship line at y=2.2
            \draw (-2,2.2) -- (6.5,2.2);
            % Vertical drops
            \draw (-2,2.2) -- (-2,1.5);
            \draw (0,2.2) -- (0,1.5);
            \draw (3,2.2) -- (3,1.5);
            \draw (6.5,2.2) -- (6.5,1.5);
            
            % Individuals Gen II
            \node[person, circle, deceased, align=center] (II1) at (-2,1.5){II-1\\AS\\$\dagger$};
            \node[person, rectangle, align=center] (II2) at (0,1.5){II-2\\AS};
            \node[person, circle, align=center] (II3) at (3,1.5){II-3\\AS};
            \node[person, circle, align=center] (II4) at (5,1.5){II-4\\AS};
            \node[person, rectangle, align=center] (II5) at (7,1.5){II-5\\AS};
            
            % Unions Gen II
            \draw (II2) -- (II3) node[midway, above] {Union};
            \draw (II4) -- (II5) node[midway, above] {Union};
            
            % Vertical drops from unions (midpoints)
            \draw (1.5,1.5) -- (1.5,0.5); % II2-II3
            \draw (6,1.5) -- (6,0.5);   % II4-II5
            
            % Generation III (y=0)
            % Sibship lines
            \draw (0.5,0.5) -- (2.5,0.5); % Children II2-II3
            \draw (5.5,0.5) -- (6.5,0.5); % Child II4-II5 (single)
            % Vertical drops
            \draw (1,0.5) -- (1,0);
            \draw (2,0.5) -- (2,0);
            \draw (6,0.5) -- (6,0);
            
            % Individuals Gen III
            \node[person, rectangle, align=center] (III1) at (1,0){III-1\\AA/AS};
            \node[person, circle, affected, align=center] (III2) at (2,0){III-2\\SS};
            \node[person, diamond, minimum size=0.8cm, align=center] (III3) at (6,0){III-3\\?};
            
            % Legend
            \node[anchor=west, font=\small] at (-3.5, -0.5) {$\dagger$ = décédé ; \tikz\draw[popdark, thick, pattern=north east lines, pattern color=popdark] (0,0) circle (0.15); = porteur (AS) ; \tikz\filldraw[poporange] (0,0) circle (0.15); = atteint (SS)};
        \end{tikzpicture}
        \legende{Pedigree de la famille pour la drépanocytose (récessif autosomique). Symboles pleins orange = atteints (SS). Symboles vides = phénotype sain (AA ou AS). Génotypes déduits indiqués. III-3 a un risque de 25\% d'être SS.}
    \end{popfigure}
    \textbf{Conclusion :} L'arbre est construit. L'analyse a révélé une incohérence dans l'énoncé initial (père II-3 dit « sain » sous-entendu AA mais fils SS), corrigée par la déduction génétique rigoureuse : II-3 est AS. \textbf{Règle d'or : la cohérence du pedigree prime sur l'énoncé littéral.}
\end{exoresolu}

\begin{methode}[Exploiter un pedigree autosomique récessif (albinisme, drépanocytose)]
    \textbf{Objectif :} Déterminer le mode de transmission, les génotypes probables et le risque de récidive.
    \begin{enumerate}
        \item \textbf{Repérer les indices de récessivité autosomique :}
              \begin{itemize}
                  \item Le caractère \textbf{saute des générations} (parents sains $\rightarrow$ enfant atteint).
                  \item Atteint $\times$ Atteint $\rightarrow$ 100\% atteints.
                  \item Atteint $\times$ Sain $\rightarrow$ 100\% sains (si sain = AA) OU 50\% atteints / 50\% porteurs (si sain = AS).
                  \item \textbf{Égalité des sexes :} Hommes et femmes atteints dans les mêmes proportions (pas de prédominance sexuelle).
              \end{itemize}
        \item \textbf{Déduire les génotypes (du haut vers le bas) :}
              \begin{itemize}
                  \item \textbf{Atteint} = homozygote mutant ($aa$ ou $SS$). \textbf{Certitude 100\%.}
                  \item \textbf{Sain ayant un parent atteint} = hétérozygote obligé ($Aa$ ou $AS$). \textbf{Certitude 100\%.}
                  \item \textbf{Sain ayant un enfant atteint} = hétérozygote obligé ($Aa$ ou $AS$). \textbf{Certitude 100\%.}
                  \item \textbf{Sain sans parent ni enfant atteint} = génotype indéterminé ($AA$ ou $Aa$). On note $A-$ ou on calcule la probabilité d'être porteur selon la fréquence allélique dans la population (si donnée) ou le croisement parental connu.
              \end{itemize}
        \item \textbf{Calculer un risque :} Identifier le croisement concerné (ex: $Aa \times Aa$) et appliquer les lois de Mendel (carré de Punnett). Exprimer en fraction (1/4), pourcentage (25\%) ou probabilité (0,25).
    \end{enumerate}
\end{methode}

\begin{exoresolu}[Analyse d'un pedigree autosomique récessif : Albinisme]
    \textbf{Énoncé :} On étudie la transmission de l'albinisme (caractère récessif autosomique, allèle $a$) dans la famille décrite. Les individus atteints sont représentés en noir. Déterminer les génotypes de tous les individus. Calculer la probabilité que le couple III-1 $\times$ III-2 ait un enfant albinos.
    \begin{center}
        \textit{Description textuelle (support d'examen) :} 
        Gén I : I-1 (homme sain) $\times$ I-2 (femme albinos). 
        Gén II : II-1 (homme sain), II-2 (femme saine), II-3 (homme sain), II-4 (femme albinos). 
        Unions : II-1 $\times$ femme saine (II-1ép) $\rightarrow$ III-1 (homme sain). 
        II-2 $\times$ II-3 (consanguins, frère/sœur) $\rightarrow$ III-2 (femme albinos). 
        II-4 $\times$ homme sain (II-4ép) $\rightarrow$ III-3 (homme sain).
    \end{center}
    \tcblower
    \textbf{Solution détaillée :}
    \begin{enumerate}
        \item \textbf{Identification du mode de transmission :}
              \begin{itemize}
                  \item I-1 sain $\times$ I-2 atteinte $\rightarrow$ enfants sains (II-1, II-2, II-3) ET atteinte (II-4). Le caractère saute une génération.
                  \item Atteints des deux sexes (I-2 femme, II-4 femme, III-2 femme).
                  \item \textbf{Conclusion : Hérédité autosomique récessive.} Allèle $a$ (albinisme), allèle $A$ (normal). $aa$ = atteint, $A-$ = sain.
              \end{itemize}
        \item \textbf{Déduction des génotypes (du haut vers le bas) :}
              \begin{itemize}
                  \item \textbf{I-2 (atteinte) :} \textbf{$aa$}.
                  \item \textbf{I-1 (sain, parent d'atteint II-4) :} Transmet obligatoirement $a$ à II-4. Donc I-1 = \textbf{$Aa$}.
                  \item \textbf{II-4 (atteinte) :} \textbf{$aa$}.
                  \item \textbf{II-1, II-2, II-3 (sains, parents $Aa \times aa$) :} Reçoivent $a$ de la mère. Phénotype sain $\rightarrow$ ont reçu $A$ du père. Donc \textbf{tous $Aa$ (porteurs obligés)}.
                  \item \textbf{II-1 ép (époux de II-1, sain, enfant III-1 sain) :} Pas d'enfant atteint $\rightarrow$ génotype indéterminé ($AA$ ou $Aa$). Noté \textbf{$A-$}.
                  \item \textbf{III-1 (sain, parents $Aa \times A-$) :} A reçu $a$ de II-1. Pour être sain, a reçu $A$ de l'autre parent. Donc III-1 = \textbf{$Aa$ (porteur)}.
                  \item \textbf{II-2 ($Aa$) $\times$ II-3 ($Aa$) (consanguins frère/sœur) :} Croisement $Aa \times Aa$.
                  \item \textbf{III-2 (atteinte) :} \textbf{$aa$}. Cohérent (probabilité 1/4).
                  \item \textbf{II-4 ($aa$) $\times$ II-4 ép (sain, enfant III-3 sain) :} II-4 transmet $a$. III-3 sain $\rightarrow$ a reçu $A$ du père. Donc II-4 ép = \textbf{$Aa$ (porteur obligé)}. III-3 = \textbf{$Aa$}.
              \end{itemize}
        \item \textbf{Risque pour le couple III-1 $\times$ III-2 :}
              \begin{itemize}
                  \item III-1 = $Aa$ (déduit). III-2 = $aa$ (atteinte).
                  \item Croisement : $Aa \times aa$.
                  \item Gamètes : III-1 $\rightarrow$ $A$ (1/2) ou $a$ (1/2). III-2 $\rightarrow$ $a$ (1).
                  \item Zygotes : $Aa$ (sain porteur) 1/2, $aa$ (albinos) 1/2.
              \end{itemize}
    \end{enumerate}
    \begin{center}
        \fbox{\textbf{Probabilité d'avoir un enfant albinos = 1/2 (soit 50\%).}}
    \end{center}
    \textbf{Conclusion :} Le couple a un risque sur deux à chaque grossesse. L'analyse rigoureuse des \textbf{porteurs obligés} (parents d'atteints, enfants sains d'atteints) est la clé de la résolution.
\end{exoresolu}

\begin{methode}[Exploiter un pedigree autosomique dominant (achondroplasie) ou codominant (groupes sanguins ABO)]
    \textbf{Objectif :} Différencier dominance et codominance, gérer l'allélisme multiple (ABO).
    \begin{enumerate}
        \item \textbf{Indices de dominance autosomique (ex: Achondroplasie) :}
              \begin{itemize}
                  \item Le caractère \textbf{ne saute pas de génération} (parent atteint $\rightarrow$ enfants atteints présents).
                  \item Atteint $\times$ Sain $\rightarrow$ ~50\% atteints (si atteint hétérozygote).
                  \item Atteint $\times$ Atteint $\rightarrow$ 75\% atteints (zygotes) si tous deux hétérozygotes ; \textbf{attention : homozygote dominant souvent létal} (achondroplasie $AA$ létal \emph{in utero}).
                  \item Égalité des sexes.
              \end{itemize}
        \item \textbf{Indices de codominance (Groupes sanguins ABO) :}
              \begin{itemize}
                  \item 3 allèles : $I^A$, $I^B$ (codominants), $i$ (récessif).
                  \item 4 phénotypes : A ($I^A I^A$ ou $I^A i$), B ($I^B I^B$ ou $I^B i$), AB ($I^A I^B$), O ($ii$).
                  \item Règle clé : Un parent O ($ii$) transmet obligatoirement $i$. Un parent AB transmet $I^A$ ou $I^B$ (jamais $i$).
              \end{itemize}
        \item \textbf{Déduction des génotypes :}
              \begin{itemize}
                  \item Atteint (dominant) = $Aa$ (souvent, car $AA$ létal ou rare). Sain = $aa$.
                  \item Groupes sanguins : Phénotype A $\rightarrow$ $I^A I^A$ ou $I^A i$. Si parent O ou enfant O $\rightarrow$ \textbf{obligatoirement $I^A i$}.
              \end{itemize}
        \item \textbf{Cas particulier Achondroplasie (létalité $AA$) :} Croisement $Aa \times Aa$ $\rightarrow$ Zygotes : 1/4 $AA$ (létaux), 1/2 $Aa$ (atteints viables), 1/4 $aa$ (sains). \textbf{Parmi les naissances vivantes :} 2/3 atteints, 1/3 sains.
    \end{enumerate}
\end{methode}

\begin{exoresolu}[Achondroplasie et létalité de l'homozygote]
    \textbf{Énoncé :} L'achondroplasie (nanisme disproportionné) est une maladie autosomique dominante. L'allèle $A$ (muté, FGFR3) est dominant sur $a$ (normal). Le génotype $AA$ est létal \emph{in utero} (homozygotie létale). Un couple d'achondroplasiques (hétérozygotes) consulte pour connaître le risque d'avoir un enfant atteint viable.
    \tcblower
    \textbf{Solution détaillée :}
    \begin{enumerate}
        \item \textbf{Génotypes des parents :} Tous deux atteints et viables $\rightarrow$ obligatoirement \textbf{$Aa$} (car $AA$ meurt avant la naissance).
        \item \textbf{Croisement théorique (méiose) :} $Aa \times Aa$.
        \item \textbf{Carré de Punnett (zygotes à la fécondation) :}
              \[
              \begin{array}{c|cc}
              \text{Gamètes} & A & a \\ \hline
              A & AA & Aa \\
              a & Aa & aa \\
              \end{array}
              \]
              Fréquences zygotiques : 1/4 $AA$, 1/2 $Aa$, 1/4 $aa$.
        \item \textbf{Sélection létale :} Les zygotes $AA$ (1/4) ne survivent pas (avortement précoce). Population viable = $Aa$ + $aa$ = 3/4 du total initial.
        \item \textbf{Fréquences phénotypiques \textbf{parmi les naissances vivantes} :}
              \begin{align*}
                  P(\text{Achondroplasie } | \text{ viable}) &= \frac{P(Aa)}{P(Aa) + P(aa)} = \frac{1/2}{1/2 + 1/4} = \frac{1/2}{3/4} = \frac{2}{3} \approx 66,7\% \\
                  P(\text{Normal } | \text{ viable}) &= \frac{P(aa)}{3/4} = \frac{1/4}{3/4} = \frac{1}{3} \approx 33,3\%
              \end{align*}
        \item \textbf{Réponse au couple :} 
              \begin{itemize}
                  \item Risque d'enfant atteint viable = \textbf{2/3}.
                  \item Risque d'enfant normal = \textbf{1/3}.
                  \item Risque d'avortement spontané précoce (zygote $AA$) = \textbf{1/4}.
              \end{itemize}
    \end{enumerate}
    \begin{attention}[Piège Bac fréquent]
        Ne jamais écrire « 3/4 d'enfants atteints » pour un couple d'achondroplasiques. Préciser toujours « \textbf{parmi les enfants nés vivants} » ou « \textbf{à la naissance} ». Le ratio 3:1 est zygotique (à la fécondation), le ratio 2:1 est phénotypique à la naissance.
    \end{attention}
    \begin{savaistu}[Achondroplasie et « nanisme thyroïdien »]
        Le programme officiel mentionne « achondroplasie (nanisme thyroïdien) ». C'est une \textbf{erreur terminologique historique}. L'achondroplasie est une dysplasie osseuse constitutionnelle (mutation FGFR3), sans lien avec la thyroïde. Le « nanisme thyroïdien » désigne le crétinisme (hypothyroïdie congénitale non traitée), maladie \textbf{non génétique} (carence en iode ou dysgénésie thyroïdienne). Au Bac, traiter l'achondroplasie comme une maladie autosomique dominante à homozygote létal.
    \end{savaistu}
\end{exoresolu}

\begin{methode}[Exploiter un pedigree gonosomique récessif (hémophilie, daltonisme) ou dominant (rachitisme vitamino-résistant)]
    \textbf{Objectif :} Mettre en évidence la transmission liée au chromosome X.
    \begin{enumerate}
        \item \textbf{Notation rigoureuse (obligatoire) :} Allèles en exposant du chromosome X. $X^H$ (allèle normal), $X^h$ (allèle mutant récessif), $Y$. Génotype homme : $X^H Y$ ou $X^h Y$. Femme : $X^H X^H$, $X^H X^h$, $X^h X^h$.
        \item \textbf{Indices de récessivité liée à l'X (Hémophilie $X^h$, Daltonisme $X^d$) :}
              \begin{itemize}
                  \item \textbf{Prédominance masculine :} Beaucoup plus d'hommes atteints que de femmes.
                  \item \textbf{Pas de transmission père $\rightarrow$ fils :} Le père donne son Y à son fils. Un père atteint a des fils sains (sauf mutation \emph{de novo}).
                  \item \textbf{Transmission père $\rightarrow$ filles :} Père atteint ($X^h Y$) $\rightarrow$ \textbf{Toutes les filles} reçoivent son $X^h$ $\rightarrow$ porteuses ($X^H X^h$) ou atteintes si mère porteuse.
                  \item \textbf{Femme atteinte rare :} Nécessite père atteint ET mère porteuse/atteinte.
              \end{itemize}
        \item \textbf{Indices de dominância liée à l'X (Rachitisme vitamino-résistant $X^R$) :}
              \begin{itemize}
                  \item \textbf{Prédominance féminine :} Plus de femmes atteintes (hétérozygotes $X^R X^r$).
                  \item \textbf{Transmission père $\rightarrow$ filles :} Père atteint $\rightarrow$ \textbf{Toutes les filles} atteintes (reçoivent $X^R$).
                  \item \textbf{Pas de transmission père $\rightarrow$ fils.}
                  \item Homme atteint ($X^R Y$) $\times$ Femme saine ($X^r X^r$) $\rightarrow$ Filles 100\% atteintes ($X^R X^r$), Fils 100\% sains ($X^r Y$).
              \end{itemize}
        \item \textbf{Calcul de risque :} Séparer selon le sexe de l'enfant à naître.
              \begin{itemize}
                  \item Ex: Mère porteuse $X^H X^h$ $\times$ Père sain $X^H Y$.
                  \item Filles : 1/2 $X^H X^H$ (saines), 1/2 $X^H X^h$ (porteuses). \textbf{0\% atteintes}.
                  \item Garçons : 1/2 $X^H Y$ (sains), 1/2 $X^h Y$ (\textbf{atteints}).
                  \item Risque global enfant atteint = 1/4 (seuls les garçons). Risque \textbf{pour un garçon} = 1/2.
              \end{itemize}
    \end{enumerate}
\end{methode}

\begin{exoresolu}[Hémophilie : Risque pour un garçon vs risque global]
    \textbf{Énoncé :} Une femme (II-2), dont le père (I-1) est hémophile et la mère (I-2) saine non porteuse, épouse un homme sain (II-1). Ils ont déjà une fille saine (III-1). La femme est à nouveau enceinte. Calculer :
    \begin{enumerate}
        \item La probabilité que cet enfant soit un garçon hémophile.
        \item La probabilité qu'un \textbf{garçon} né de ce couple soit hémophile.
    \end{enumerate}
    \tcblower
    \textbf{Solution détaillée :}
    \begin{enumerate}
        \item \textbf{Génotypes des grands-parents (Gén I) :}
              \item I-1 (père hémophile) : $X^h Y$.
              \item I-2 (mère saine non porteuse) : $X^H X^H$.
        \item \textbf{Génotype de la mère (II-2) :} Elle reçoit $X^h$ de son père (obligatoire) et $X^H$ de sa mère. Donc II-2 = \textbf{$X^H X^h$ (porteuse hétérozygote)}. Certitude = 1.
        \item \textbf{Génotype du père (II-1) :} Homme sain $\rightarrow$ $X^H Y$.
        \item \textbf{Croisement :} $X^H X^h \times X^H Y$.
        \item \textbf{Gamètes :} Mère : $X^H$ (1/2), $X^h$ (1/2). Père : $X^H$ (1/2), $Y$ (1/2).
        \item \textbf{Tableau des descendants (zygotes) :}
              \[
              \begin{array}{c|cc}
              \text{Mère} \backslash \text{Père} & X^H (1/2) & Y (1/2) \\ \hline
              X^H (1/2) & X^H X^H \text{ (Fille saine) } 1/4 & X^H Y \text{ (Garçon sain) } 1/4 \\
              X^h (1/2) & X^H X^h \text{ (Fille porteuse) } 1/4 & X^h Y \text{ (Garçon hémophile) } 1/4 \\
              \end{array}
              \]
        \item \textbf{Réponses :}
              \begin{enumerate}
                  \item \textbf{Probabilité (Enfant = Garçon hémophile) = 1/4 (25\%).} (Case en bas à droite : $P(X^h) \times P(Y) = 1/2 \times 1/2$).
                  \item \textbf{Probabilité (Garçon hémophile | Enfant est un garçon) = } $\frac{P(\text{Garçon hémophile})}{P(\text{Garçon})} = \frac{1/4}{1/2} = \textbf{1/2 (50\%)}$.
              \end{enumerate}
    \end{enumerate}
    \begin{aretenir}[Distinction cruciale pour le Bac]
        \begin{itemize}
            \item \textbf{Risque global (à la conception, sexe inconnu) :} $P(\text{garçon malade}) = 1/4$.
            \item \textbf{Risque conditionnel (sexe connu = garçon) :} $P(\text{malade} | \text{garçon}) = 1/2$.
            \item \textbf{Risque conditionnel (sexe connu = fille) :} $P(\text{malade} | \text{fille}) = 0$ (pour récessif X).
            \item \textbf{Consigne :} Lire attentivement : « probabilité qu'un enfant soit un garçon malade » $\neq$ « probabilité qu'un garçon soit malade ».
        \end{itemize}
    \end{aretenir}
\end{exoresolu}

\begin{methode}[Évaluer un risque génétique (synthèse : autosomique + gonosomique + consanguinité)]
    \textbf{Objectif :} Calculer la probabilité qu'un couple à risque engendre un enfant atteint, en intégrant l'incertitude sur le génotype des parents (probabilité d'être porteur).
    \begin{enumerate}
        \item \textbf{Étape 1 : Déterminer le mode de transmission} (pedigree ou énoncé : autosomique récessif/dominant, gonosomique récessif/dominant).
        \item \textbf{Étape 2 : Calculer la probabilité d'être porteur pour chaque parent consultant.}
              \begin{itemize}
                  \item \textbf{Parent sain, issu d'un croisement connu (ex: parents $Aa \times Aa$) :} Conditionner par le phénotype sain.
                      $P(AA | \text{sain}) = \frac{1/4}{3/4} = 1/3$ ; $P(Aa | \text{sain}) = \frac{1/2}{3/4} = 2/3$.
                  \item \textbf{Parent sain, apparenté à un atteint (oncle/tante, cousin) :} Remonter le pedigree. Probabilité de transmettre l'allèle mutant = $(1/2)^n$ où $n$ = nombre de méioses entre l'ancêtre porteur connu (ou atteint) et le parent. \textbf{Conditionner par le phénotype sain des intermédiaires.}
                  \item \textbf{Parent sain, population générale (pas d'antécédent connu) :} Utiliser la fréquence des porteurs $2pq \approx 2q$ (si $q$ faible, maladie rare). Donnée fournie au Bac.
              \end{itemize}
        \item \textbf{Étape 3 : Croiser les probabilités (indépendance des allèles).}
              \begin{itemize}
                  \item Récessif autosomique : $P(\text{Enfant atteint}) = P(\text{Mère porteuse}) \times P(\text{Père porteur}) \times \frac{1}{4}$.
                  \item Récessif X : $P(\text{Enfant atteint}) = P(\text{Mère porteuse}) \times \frac{1}{2} \times \frac{1}{2}$ (seuls les garçons atteints, ou 1/4 global). Pour une fille : $P(\text{Mère porteuse}) \times P(\text{Père atteint}) \times 1/2$ (rare).
                  \item Dominant autosomique (pénétrance complète) : $P(\text{Enfant atteint}) = P(\text{Parent porteur allèle dominant}) \times 1/2$.
              \end{itemize}
        \item \textbf{Étape 4 : Exprimer le résultat} en fraction irréductible, pourcentage ou décimal (ex: 1/12, 8,3\%, 0,083).
    \end{enumerate}
\end{methode}

\begin{exoresolu}[Risque génétique avec incertitude parentale (Cousins germains, Drépanocytose)]
    \textbf{Énoncé :} La drépanocytose est récessive autosomique (allèle $S$, fréquence allélique $q = 0,1$ en zone endémique camerounaise). Un couple de cousins germains (III-1 et III-2) consulte. Leur grand-père commun (I-1) était drépanocytaire (SS). Leur grand-mère commune (I-2) était saine non porteuse (AA). Les parents respectifs (II-1 et II-2) sont sains. Calculer le risque que leur enfant (IV-1) soit drépanocytaire.
    \tcblower
    \textbf{Solution détaillée :}
    \begin{enumerate}
        \item \textbf{Analyse du pedigree (Gén I $\rightarrow$ Gén II) :}
              \item I-1 ($SS$) $\times$ I-2 ($AA$) $\rightarrow$ Tous les enfants (II-1, II-2, ...) sont \textbf{obligatoirement $AS$ (porteurs obligés)}.
        \item \textbf{Génotype des parents de la génération III (II-1 et II-2) :} Tous deux $AS$ (certitude 1).
        \item \textbf{Calcul de $P(\text{III-1 est } AS | \text{III-1 sain})$ :}
              \begin{itemize}
                  \item II-1 ($AS$) épouse un conjoint de la population générale.
                  \item Fréquence porteur population $P(AS) = 2q = 0,2$ ; $P(AA) = 0,8$.
                  \item \textbf{Cas 1 : Conjoint = $AA$ (prob 0,8).} Croisement $AS \times AA \rightarrow$ 1/2 $AA$, 1/2 $AS$. Tous sains. $P(AS | \text{sain}) = 1/2$.
                  \item \textbf{Cas 2 : Conjoint = $AS$ (prob 0,2).} Croisement $AS \times AS \rightarrow$ 1/4 $AA$, 1/2 $AS$, 1/4 $SS$. Condition « sain » retire $SS$. $P(AS | \text{sain}) = \frac{1/2}{3/4} = 2/3$.
              \end{itemize}
              \textbf{Loi des probabilités totales :}
              \[
              P(\text{III-1 } AS) = 0,8 \times \frac{1}{2} + 0,2 \times \frac{2}{3} = 0,4 + \frac{0,4}{3} = \frac{1,2 + 0,4}{3} = \frac{1,6}{3} = \frac{8}{15}
              \]
        \item \textbf{Par symétrie} (même structure pour III-2 via II-2) : $P(\text{III-2 } AS) = \frac{8}{15}$.
        \item \textbf{Risque pour l'enfant IV-1 (SS) :}
              \item Les deux parents (III-1, III-2) doivent être porteurs ($AS$) pour avoir un enfant $SS$.
              \item $P(\text{IV-1 } SS) = P(\text{III-1 } AS) \times P(\text{III-2 } AS) \times \frac{1}{4}$.
              \item $P = \frac{8}{15} \times \frac{8}{15} \times \frac{1}{4} = \frac{64}{900} = \frac{16}{225}$.
    \end{enumerate}
    \begin{center}
        \fbox{\textbf{Risque génétique = $\frac{16}{225} \approx 0,071$ (soit 7,1\%).}}
    \end{center}
    \textbf{Comparaison :} Sans consanguinité (deux individus pris au hasard dans la population), risque = $(2q)^2 \times 1/4 = q^2 = 0,01$ (1\%). La consanguinité (cousins germains) multiplie le risque par $\approx 7$ ici.
\end{exoresolu}

\begin{methode}[Expliquer la nécessité des examens prénuptiaux et prénataux]
    \textbf{Objectif :} Argumenter scientifiquement et éthiquement la prévention des maladies héréditaires graves (conseil génétique).
    \begin{enumerate}
        \item \textbf{Examen prénuptial (Conseil génétique préconceptionnel - Prévention primaire) :}
              \begin{itemize}
                  \item \textbf{Cible :} Couples à risque (antécédents familiaux connus, consanguinité, population à forte prévalence ex: drépanocytose au Cameroun $q \approx 0,1$).
                  \item \textbf{Actes biologiques :} Électrophorèse de l'hémoglobine (drépanocytose), caryotype (anomalies chromosomiques), recherche mutations connues (fibrose kystique, thalassémie, amyotrophie spinale), sérologies MST (VIH, Hépatites B/C, Syphilis).
                  \item \textbf{Intérêt :} Informer le couple \textbf{AVANT} la grossesse. Permettre un choix éclairé : union en connaissance de cause, adoption, PMA avec DPI (Diagnostic Préimplantatoire sur embryon FIV), don de gamètes.
                  \item \textbf{Contexte Cameroun :} Dépistage systématique de la drépanocytose avant le mariage fortement recommandé (forte prévalence). Coût faible (quelques milliers FCFA) vs coût de prise en charge vie entière (hospitalisations, hydroxyurée, transfusions).
              \end{itemize}
        \item \textbf{Examen prénatal (Diagnostic Prénatal - DP / Dépistage Prénatal - Prévention secondaire) :}
              \begin{itemize}
                  \item \textbf{Cible :} Grossesse en cours chez couple à risque connu ou dépistage populationnel (ex: Trisomie 21).
                  \item \textbf{Techniques invasives (Diagnostiques - Caryotype / Biologie moléculaire) :}
                        \begin{itemize}
                            \item \textbf{Amniocentèse (15-18 SA) :} Prélèvement liquide amniotique (amniocytes). Risque fausse couche $\approx 0,5 - 1\%$. Référence pour caryotype / ADN fœtal.
                            \item \textbf{Trophobiopsie / Biopsie des villosités choriales (11-13 SA) :} Prélèvement trophoblaste. Plus précoce. Risque fausse couche $\approx 1-2\%$. Risque mosaïcisme placentaire (faux positifs/négatifs).
                        \end{itemize}
                  \item \textbf{Techniques non invasives (Dépistage - Calcul de risque) :}
                        \begin{itemize}
                            \item \textbf{Dépistage combiné 1er trimestre (11-13 SA) :} Échographie (clarté nucale) + Dosage sérique maternel (PAPP-A, $\beta$-hCG libre). Calcule un risque (ex: Trisomie 21). \textbf{Ne donne pas le diagnostic.}
                            \item \textbf{ADN fœtal circulant (NIPT / DPNI) :} Prise de sang maternel (dès 10 SA). Analyse ADN fœtal plasmatique (cfDNA). Sensibilité/Spécificité $> 99\%$ pour T21, T18, T13. \textbf{Test de dépistage très performant, pas diagnostique.} Confirmation par amniocentèse si positif.
                        \end{itemize}
                  \item \textbf{Intérêt :} Préparer la prise en charge néonatale (ex: chirurgie cardiaque T21), organiser l'accouchement en centre spécialisé (niveau 3), ou \textbf{demande d'Interruption Médicale de Grossesse (IMG)} si pathologie grave incurable (légal au Cameroun : Loi 90/013 art 339, « grave danger pour la santé de la mère » ou « malformation grave incurable » confirmée par 2 médecins).
              \end{itemize}
        \item \textbf{Argumentation type Bac (structure attendue) :}
              \begin{quote}
              « L'examen prénuptial permet la \textbf{prévention primaire} (éviter la conception à risque par choix éclairé, DPI). L'examen prénatal permet la \textbf{prévention secondaire} (détection précoce \emph{in utero} pour prise en charge adaptée à la naissance ou IMG). Ils réduisent l'incidence et la mortalité des maladies graves (drépanocytose, trisomies, malformations) et la souffrance des familles, dans le respect de l'autonomie du couple (consentement libre et éclairé) et de la législation en vigueur. »
              \end{quote}
    \end{enumerate}
\end{methode}

\begin{exoresolu}[Argumentation : Dépistage néonatal de la drépanocytose au Cameroun]
    \textbf{Énoncé :} Au Cameroun, la prévalence de l'hétérozygotie drépanocytaire (AS) est d'environ 20\% à 25\% ($2q$). La drépanocytose homozygote (SS) touche environ 1 à 2\% des naissances. Expliquer pourquoi le dépistage néonatal systématique (test de Guthrie / électrophorèse hémoglobine à la naissance) est une mesure de santé publique prioritaire, alors qu'il ne prévient pas la naissance des enfants malades.
    \tcblower
    \textbf{Solution détaillée (Structure argumentée pour dissertation) :}
    \begin{enumerate}
        \item \textbf{Distinction des niveaux de prévention :}
              \begin{itemize}
                  \item \textbf{Primaire (Prénuptiale) :} Évite la conception (conseil génétique, DPI). \textit{Idéal mais couverture encore insuffisante au Cameroun.}
                  \item \textbf{Secondaire (Prénatale) :} Détecte \emph{in utero} $\rightarrow$ IMG possible. \textit{Accès limité (coût, technique, délai légal, acceptabilité culturelle/religieuse).}
                  \item \textbf{Tertiaire (Néonatale) :} Détecte \textbf{à la naissance} avant les symptômes. \textbf{Ne prévient pas la maladie, mais prévient la MORTALITÉ et la MORBIDITÉ SÉVÈRE.}
              \end{itemize}
        \item \textbf{Histoire naturelle de la drépanocytose non dépistée (données OMS/Cameroun) :}
              \begin{itemize}
                  \item Enfant apparemment sain à la naissance (HbF fœtale protectrice $\approx$ 80\%).
                  \item Premiers symptômes vers 6-12 mois : anémie hémolytique chronique, crises vaso-occlusives (douleurs osseuses/abdominales), \textbf{sepsis à \emph{Streptococcus pneumoniae} (pneumocoque) : risque majeur de décès avant 5 ans}, syndrome main-pied, accidents vasculaires cérébraux (AVC) silencieux ou déclarés.
                  \item Mortalité infantile non dépistée/non prise en charge : $\approx 50\%$ avant 5 ans en zone tropicale.
              \end{itemize}
        \item \textbf{Apport majeur du dépistage néonatal (DNN) précoce (résultat avant 2 mois de vie) :}
              \begin{itemize}
                  \item \textbf{Prophylaxie pénicilline V orale (dès 2 mois jusqu'à 5 ans minimum) :} Réduit le risque de sepsis pneumococcique invasif de \textbf{84\%} (Étude PROPS, NEJM 1986). \textbf{Mesure salvatrice n°1, coût dérisoire.}
                  \item \textbf{Vaccination anti-pneumococcique conjuguée (PCV13/PCV10) et anti-\emph{Haemophilus influenzae} b :} Renforce la protection (calendrier EPI élargi).
                  \item \textbf{Suivi spécialisé en Centre de Drépanocytose (ex: Hôpital Mère-Enfant Fondation Chantal Biya, Hôpital Laquintinie, Hôpital Général de Douala) :} Éducation parents (reconnaître fièvre $>38,5^\circ$C = urgence, palpation rate, déshydratation), hydroxyurée (dès 9 mois si forme sévère) $\rightarrow$ augmentation HbF, réduction crises, AVC, transfusions.
                  \item \textbf{Détection des porteurs (AS) :} Information pour la fratrie future (conseil génétique différé).
              \end{itemize}
        \item \textbf{Rentabilité (Analyse Coût-Efficacité) :}
              \item Coût test (électrophorèse sur cellulose / HPLC) $\approx$ 1 000 - 2 000 FCFA.
              \item Coût prise en charge d'une crise sévère / sepsis / AVC $\gg$ 100 000 FCFA (hospitalisation, transfusions, réanimation).
              \item Années de vie gagnées (AVG) et années de vie ajustées sur l'incapacité (DALYs évités) énormes.
        \item \textbf{Conclusion :} Le DNN est la \textbf{stratégie la plus équitable, accessible et efficace} actuellement au Cameroun pour réduire drastiquement la mortalité drépanocytaire infantile, en attendant la généralisation du conseil prénuptial et de l'accès au DPI/DPN. Il transforme une maladie mortelle silencieuse en pathologie chronique gérée.
    \end{enumerate}
\end{exoresolu}

\begin{attention}[Les 5 erreurs qui coûtent des points au Bac (Check-list finale)]
    \begin{enumerate}
        \item \textbf{Confondre « Risque global » et « Risque conditionnel au sexe » (Gonosomique).}
              \textit{Exemple :} « Risque d'avoir un garçon hémophile = 1/2 » \textbf{FAUX}. C'est 1/4 (global) ou 1/2 (SI c'est un garçon). \textbf{Lire la question : « un enfant » vs « un garçon ».}
        \item \textbf{Oublier la létalité de l'homozygote dominant (Achondroplasie).}
              \textit{Écrire :} « Croisement $Aa \times Aa \rightarrow$ 3/4 atteints ». \textbf{FAUX}. C'est 2/3 parmi les nés vivants. Préciser toujours « \textbf{zygotes} » ou « \textbf{nés vivants} ».
        \item \textbf{Ne pas déduire le génotype « porteur obligé » (Récessif autosomique).}
              \textit{Exemple :} Parent sain d'un enfant atteint $\rightarrow$ Génotype noté $A-$ ou $AA$ possible. \textbf{FAUX}. C'est \textbf{obligatoirement hétérozygote $Aa$}. Écrire « Porteur obligé (100\%) ».
        \item \textbf{Erreur de notation allélique (Gonosomique / Groupes sanguins).}
              \begin{itemize}
                  \item Écrire $Xh$ au lieu de $X^h$ (exposant obligatoire sur le chromosome).
                  \item Écrire $I^A I^B$ pour groupe A. \textbf{FAUX}. Groupe A = $I^A I^A$ ou $I^A i$. $I^A I^B$ = Groupe AB (codominance).
                  \item Oublier le chromosome Y chez l'homme ($X^H Y$, pas $X^H$ seul).
              \end{itemize}
        \item \textbf{Calcul de risque consanguin : Ne pas conditionner par le phénotype « sain » des intermédiaires.}
              \textit{Exemple :} Grand-père $SS$, petit-fils sain. Calcul direct $(1/2)^2 = 1/4$ pour être porteur. \textbf{FAUX}. Le petit-fils est sain, il ne peut pas être $SS$. La probabilité d'être $AS$ sachant qu'il est sain est plus élevée (ex: $2/3$ si parents $AS \times AS$). \textbf{Toujours conditionner par le phénotype connu (sain/atteint) à chaque génération.}
    \end{enumerate}
\end{attention}

\newpage

\coursec{Exercices}

\courssub{Je vérifie mes connaissances}

\begin{exercice}[QCM : modes de transmission]
Pour chaque affirmation, choisir la (ou les) bonne(s) réponse(s) et justifier brièvement.
\begin{enumerate}
\item Une maladie génétique est dite autosomique récessive lorsque :
\begin{enumerate}
\item l'allèle responsable est porté par un chromosome sexuel ;
\item l'allèle responsable ne s'exprime qu'à l'état homozygote ;
\item les parents sains peuvent avoir un enfant atteint ;
\item les deux sexes sont atteints avec la même fréquence.
\end{enumerate}
\item Dans une hérédité gonosomique récessive liée à $X$ (hémophilie) :
\begin{enumerate}
\item un père atteint transmet l'allèle à tous ses fils ;
\item un père atteint transmet l'allèle à toutes ses filles ;
\item une femme conductrice a un risque de 1/2 de transmettre l'allèle à chacun de ses fils ;
\item une fille ne peut être atteinte que si son père est atteint et sa mère conductrice.
\end{enumerate}
\item Dans le système ABO, les allèles $A$ et $B$ sont :
\begin{enumerate}
\item récessifs l'un par rapport à l'autre ;
\item codominants entre eux et dominants sur $O$ ;
\item portés par le chromosome $X$ ;
\item responsables d'une hérédité gonosomique.
\end{enumerate}
\item Le rachitisme vitamino-résistant se transmet selon un mode :
\begin{enumerate}
\item autosomique récessif ;
\item autosomique dominant ;
\item gonosomique dominant lié à $X$ ;
\item gonosomique récessif lié à $Y$.
\end{enumerate}
\end{enumerate}
\end{exercice}

\begin{exercice}[Vrai ou faux justifié]
Répondre par \textbf{vrai} ou \textbf{faux}, puis justifier chaque réponse en une ou deux phrases.
\begin{enumerate}
\item Deux parents sains ne peuvent jamais avoir un enfant albinos.
\item Une femme hémophile a obligatoirement un père hémophile.
\item Un homme daltonien a reçu son allèle daltonien de son père.
\item Dans l'achondroplasie, l'allèle responsable s'exprime dès qu'il est présent en un seul exemplaire.
\item Un individu de groupe sanguin AB possède les allèles $A$ et $B$ sur ses deux chromosomes homologues.
\item Tous les fils d'une femme conductrice de l'hémophilie seront hémophiles.
\end{enumerate}
\end{exercice}

\begin{exercice}[Définitions et symboles du pedigree]
\begin{enumerate}
\item Définir : arbre généalogique (pedigree) ; maladie génique ; maladie chromosomique ; risque génétique.
\item Représenter les symboles conventionnels suivants : homme sain, femme atteinte, union, descendance, jumeaux monozygotes.
\item Expliquer la différence entre une hérédité autosomique et une hérédité gonosomique, en citant un exemple de maladie pour chaque cas.
\item Qu'appelle-t-on femme \cle{conductrice} ? Pourquoi cette notion ne s'applique-t-elle qu'aux hérédités récessives ?
\end{enumerate}
\end{exercice}

\begin{exercice}[Calculs directs de risques]
\begin{enumerate}
\item Un couple dont les deux membres sont hétérozygotes $[A//a]$ pour un allèle récessif $a$ responsable de l'albinisme. Calculer la probabilité que leur premier enfant soit albinos, puis la probabilité qu'il soit sain.
\item Une femme conductrice de l'hémophilie ($X^{H}/X^{h}$) épouse un homme sain. Calculer la probabilité d'avoir : un garçon hémophile ; une fille hémophile ; un enfant (sexes confondus) hémophile.
\item Un homme de groupe sanguin O épouse une femme de groupe AB. Déterminer les groupes sanguins possibles de leurs enfants et leurs probabilités.
\end{enumerate}
\end{exercice}

\courssub{Je m'entraîne}

\begin{exercice}[Pedigree de l'albinisme]
Dans une famille de la région de l'Ouest du Cameroun, deux parents sains (II-3 et II-4) ont quatre enfants : deux filles saines, un garçon sain et une fille albinos (III-2).
\begin{enumerate}
\item Construire l'arbre généalogique de cette famille sur trois générations (les parents de II-3 sont I-1 et I-2, tous deux sains).
\item Déterminer le mode de transmission de l'albinisme dans cette famille. Justifier.
\item Donner les génotypes certains de II-3, II-4 et III-2 (on notera $A$ l'allèle normal et $a$ l'allèle albinos).
\item III-1, fille saine, épouse un homme albinos. Quel est le risque que leur premier enfant soit albinos ? (On considérera les deux génotypes possibles de III-1.)
\end{enumerate}
\end{exercice}

\begin{exercice}[Groupes sanguins ABO]
Un couple, tous deux de groupe sanguin A, a un premier enfant de groupe O.
\begin{enumerate}
\item Expliquer, à l'aide d'un échiquier de croisement, comment cela est possible.
\item Déterminer les génotypes certains des deux parents.
\item Établir la liste des groupes sanguins possibles pour les enfants suivants de ce couple, avec leurs probabilités respectives.
\item Le couple appartient aussi au système Rhésus : le père est Rh$^{+}$ homozygote et la mère Rh$^{-}$. Quel sera le phénotype Rhésus de tous leurs enfants ? Quel problème médical peut se poser lors d'une seconde grossesse ?
\end{enumerate}
\end{exercice}

\begin{exercice}[Transmission du daltonisme]
Une femme à la vision des couleurs normale, dont le père était daltonien, épouse un homme à la vision normale.
\begin{enumerate}
\item Rappeler le mode de transmission du daltonisme.
\item Déterminer les génotypes de la femme, de son père et de son mari (on notera $X^{D}$ l'allèle normal et $X^{d}$ l'allèle daltonien).
\item À l'aide d'un échiquier de croisement, déterminer les proportions phénotypiques attendues dans la descendance de ce couple, en distinguant garçons et filles.
\item Une fille de ce couple peut-elle être daltonienne ? Justifier.
\end{enumerate}
\end{exercice}

\begin{exercice}[Rachitisme vitamino-résistant]
Un homme atteint de rachitisme vitamino-résistant épouse une femme saine. Ils ont trois filles, toutes atteintes, et deux fils, tous sains.
\begin{enumerate}
\item Construire le pedigree de cette famille.
\item Montrer que ces observations sont compatibles avec une transmission gonosomique dominante liée à $X$. Justifier chaque affirmation.
\item Donner les génotypes de tous les membres de la famille (on notera $X^{R}$ l'allèle responsable et $X^{r}$ l'allèle normal).
\item Une des filles atteintes épouse un homme sain. Prédire, à l'aide d'un échiquier de croisement, les proportions d'enfants atteints selon le sexe.
\end{enumerate}
\end{exercice}

\courssub{Je me prépare au Bac}

\begin{exercice}[Type Bac : albinisme sur quatre générations]
Au cours d'une consultation de génétique dans un hôpital de Yaoundé, on établit le pedigree d'une famille sur quatre générations. On note $A$ l'allèle normal (pigmentation normale) et $a$ l'allèle responsable de l'albinisme. Les individus I-1 (homme) et I-2 (femme) sont sains ; leur fils II-3 est albinos ; II-3 épouse une femme saine II-4 dont le père était albinos. Le couple III-2 (femme saine, fille de II-3) et III-3 (homme sain, sans antécédent familial) attend un enfant.
\begin{enumerate}
\item Construire le pedigree complet de cette famille en respectant les conventions.
\item Déterminer le mode de transmission de l'albinisme dans cette famille. Justifier par deux arguments tirés du pedigree.
\item Justifier les génotypes des individus I-1, I-2, II-3 et II-4.
\item Calculer le risque pour le couple III-2 $\times$ III-3 d'avoir un enfant albinos, en distinguant les hypothèses sur le génotype de III-3.
\item Rédiger un court conseil génétique à l'attention de ce couple, en expliquant l'intérêt d'un examen prénuptial et, le cas échéant, prénatal.
\end{enumerate}
\end{exercice}

\begin{exercice}[Type Bac : électrophorèse de l'hémoglobine et drépanocytose]
Dans un centre de santé de Douala, on réalise l'électrophorèse de l'hémoglobine de quatre membres d'une famille : le père (I-1), la mère (I-2), leur fille (II-1) et leur fils (II-2). Les résultats sont les suivants : I-1 présente deux bandes (HbA et HbS) ; I-2 présente deux bandes (HbA et HbS) ; II-1 présente une seule bande (HbA) ; II-2 présente une seule bande (HbS). On rappelle que l'allèle $HbS$ est responsable de la drépanocytose et que les hétérozygotes $[HbA//HbS]$ sont cliniquement sains (porteurs du trait drépanocytaire).
\begin{enumerate}
\item Identifier les génotypes de chacun des quatre membres de la famille à partir des profils d'électrophorèse. Justifier.
\item Préciser le mode de transmission de la drépanocytose et expliquer pourquoi les hétérozygotes sont dits « sains » alors qu'ils possèdent l'allèle $HbS$.
\item Le couple I-1 $\times$ I-2 attend un troisième enfant. À l'aide d'un échiquier de croisement, calculer le risque que cet enfant soit drépanocytaire, puis le risque qu'il soit porteur du trait.
\item On estime qu'au Cameroun environ 20 \% de la population est porteuse du trait drépanocytaire. Expliquer pourquoi l'allèle $HbS$ se maintient à une fréquence aussi élevée dans les zones d'endémie palustre.
\item Expliquer en quoi l'examen prénuptial (électrophorèse de l'hémoglobine) permet de limiter la fréquence de la drépanocytose dans les familles, et rédiger un conseil génétique pour un couple dont les deux membres seraient porteurs du trait.
\end{enumerate}
\end{exercice}

\begin{exercice}[Type Bac : hémophilie et conseil génétique]
Un homme hémophile épouse une femme saine dont le père était hémophile. On note $X^{H}$ le chromosome $X$ portant l'allèle normal et $X^{h}$ celui portant l'allèle responsable de l'hémophilie.
\begin{enumerate}
\item Rappeler le mode de transmission de l'hémophilie et justifier le fait que les hommes sont beaucoup plus fréquemment atteints que les femmes.
\item Déterminer les génotypes du mari, de la femme et du père de la femme. Justifier chaque génotype.
\item À l'aide d'un échiquier de croisement, déterminer pour ce couple : la probabilité d'avoir un garçon hémophile ; la probabilité d'avoir une fille hémophile ; la probabilité d'avoir une fille conductrice ; la probabilité globale d'avoir un enfant hémophile.
\item Le couple souhaite un deuxième enfant. Les risques calculés à la question 3 sont-ils modifiés ? Justifier.
\item Expliquer l'intérêt des examens prénataux (caryotype, recherche de l'allèle) pour cette famille, et dégager les limites éthiques du conseil génétique.
\end{enumerate}
\end{exercice}

\newpage

\coursec{Corrigés des exercices}

\begin{corrige}[Exercice 1 — QCM : modes de transmission]
\begin{enumerate}
\item \textbf{Bonnes réponses : b, c, d.}
\begin{itemize}
\item a) Faux : « autosomique » signifie précisément que l'allèle est porté par un \cle{autosome} (chromosome non sexuel), et non par un chromosome sexuel.
\item b) Vrai : par définition, un allèle \cle{récessif} ne s'exprime qu'à l'état \cle{homozygote} $[a//a]$.
\item c) Vrai : deux parents sains mais \cle{conducteurs} $[A//a]$ peuvent engendrer un enfant $[a//a]$ atteint (probabilité $1/4$).
\item d) Vrai : les autosomes étant identiques chez les deux sexes, la maladie frappe indifféremment garçons et filles.
\end{itemize}
\item \textbf{Bonnes réponses : b, c, d.}
\begin{itemize}
\item a) Faux : un père transmet son chromosome $Y$ à ses fils, jamais son $X$ ; il ne peut donc pas leur transmettre un allèle lié à $X$.
\item b) Vrai : un père atteint $X^{h}Y$ donne obligatoirement son unique $X^{h}$ à toutes ses filles.
\item c) Vrai : une conductrice $X^{H}/X^{h}$ produit deux types d'ovules équiprobables ($X^{H}$ et $X^{h}$) ; chaque fils, qui reçoit son $X$ de sa mère, a donc un risque de $1/2$ d'être hémophile.
\item d) Vrai : une fille atteinte doit être $X^{h}/X^{h}$, ce qui exige un père atteint ($X^{h}Y$) et une mère au moins conductrice.
\end{itemize}
\item \textbf{Bonne réponse : b.} $A$ et $B$ sont \cle{codominants} (ils s'expriment tous deux chez l'hétérozygote $[A//B]$, de groupe AB) et tous deux \cle{dominants} sur $O$. Le gène ABO est porté par un autosome (chromosome 9) : il ne s'agit pas d'une hérédité gonosomique (c et d faux, a faux car il n'y a pas de relation de récessivité entre $A$ et $B$).
\item \textbf{Bonne réponse : c.} Le rachitisme vitamino-résistant se transmet selon un mode \cle{gonosomique dominant lié à $X$} : un père atteint transmet la maladie à toutes ses filles et à aucun de ses fils.
\end{enumerate}
\end{corrige}

\begin{corrige}[Exercice 2 — Vrai ou faux justifié]
\begin{enumerate}
\item \textbf{Faux.} Deux parents sains peuvent être tous deux conducteurs $[A//a]$ ; leur enfant a alors une probabilité de $1/4$ d'être albinos $[a//a]$.
\item \textbf{Vrai.} Une femme hémophile est $X^{h}/X^{h}$ : elle a reçu un $X^{h}$ de chaque parent, donc son père est obligatoirement $X^{h}Y$, c'est-à-dire hémophile.
\item \textbf{Faux.} Un garçon reçoit le chromosome $Y$ de son père et son unique $X$ de sa mère : l'allèle daltonien d'un homme daltonien provient donc de sa \cle{mère} (conductrice ou daltonienne).
\item \textbf{Vrai.} L'achondroplasie est une hérédité autosomique \cle{dominante} : l'allèle muté s'exprime à l'état hétérozygote, dès qu'il est présent en un seul exemplaire.
\item \textbf{Vrai.} Le groupe AB correspond au génotype $[A//B]$ : l'allèle $A$ est porté par l'un des chromosomes homologues et l'allèle $B$ par l'autre ; $A$ et $B$ sont codominants.
\item \textbf{Faux.} Chaque fils d'une conductrice $X^{H}/X^{h}$ a une probabilité de $1/2$ de recevoir $X^{h}$ et d'être hémophile, et de $1/2$ de recevoir $X^{H}$ et d'être sain : ce n'est pas une certitude.
\end{enumerate}
\end{corrige}

\begin{corrige}[Exercice 3 — Définitions et symboles du pedigree]
\begin{enumerate}
\item \textbf{Définitions.}
\begin{itemize}
\item \textbf{Arbre généalogique (pedigree)} : représentation schématique, à l'aide de symboles conventionnels, des liens de parenté et de la répartition d'un caractère (ou d'une maladie) au sein d'une famille sur plusieurs générations.
\item \textbf{Maladie génique} : maladie due à la mutation d'un seul gène (ex. : drépanocytose, hémophilie).
\item \textbf{Maladie chromosomique} : maladie due à une anomalie du nombre ou de la structure des chromosomes (ex. : trisomie 21).
\item \textbf{Risque génétique} : probabilité, calculée à partir du pedigree et du mode de transmission, qu'un individu (souvent un enfant à naître) soit atteint d'une maladie héréditaire.
\end{itemize}
\item \textbf{Symboles conventionnels.}
\begin{popfigure}
\begin{tikzpicture}[scale=0.9]
% homme sain
\draw[thick,popblue] (0,0) rectangle (0.7,0.7);
\node[below] at (0.35,0) {\small homme sain};
% femme atteinte
\draw[thick,popblue,fill=popblue] (2.2,0) circle (0.35);
\node[below] at (2.2,-0.35) {\small femme atteinte};
% union + descendance
\draw[thick,popblue] (4.4,0) rectangle (5.1,0.7);
\draw[thick,popblue] (6.1,0.35) circle (0.35);
\draw[thick,popblue] (5.1,0.35) -- (5.75,0.35);
\node[below] at (5.3,0) {\small union};
\draw[thick,popblue] (5.45,0.35) -- (5.45,-0.9);
\draw[thick,popblue] (4.9,-0.9) -- (6.0,-0.9);
\draw[thick,popblue] (4.9,-0.9) -- (4.9,-1.3);
\draw[thick,popblue] (6.0,-0.9) -- (6.0,-1.3);
\draw[thick,popblue] (4.7,-1.3) rectangle (5.1,-0.95);
\draw[thick,popblue] (6.0,-1.65) circle (0.35);
\node[below] at (5.45,-1.95) {\small descendance};
% jumeaux monozygotes
\draw[thick,popblue] (8.6,0.35) -- (8.6,-0.5);
\draw[thick,popblue] (8.6,-0.5) -- (8.1,-1.1);
\draw[thick,popblue] (8.6,-0.5) -- (9.1,-1.1);
\draw[thick,popblue] (8.35,-0.8) -- (8.85,-0.8);
\draw[thick,popblue] (7.9,-1.1) rectangle (8.3,-0.75);
\draw[thick,popblue] (8.9,-1.1) rectangle (9.3,-0.75);
\node[below] at (8.6,-1.4) {\small jumeaux monozygotes};
\end{tikzpicture}
\legende{Symboles du pedigree : carré = homme, cercle = femme, symbole rempli = individu atteint ; trait horizontal = union, traits verticaux = descendance ; les jumeaux monozygotes sont reliés par un trait entre les deux branches.}
\end{popfigure}
\item \textbf{Hérédité autosomique} : le gène responsable est porté par un autosome ; les deux sexes sont atteints avec la même fréquence (ex. : albinisme, drépanocytose). \textbf{Hérédité gonosomique} : le gène est porté par un chromosome sexuel (le plus souvent $X$) ; la répartition de la maladie dépend du sexe (ex. : hémophilie, daltonisme).
\item Une femme \cle{conductrice} est une femme hétérozygote $X^{H}/X^{h}$, phénotypiquement saine, qui porte un allèle récessif responsable d'une maladie et peut le transmettre à sa descendance. Cette notion ne s'applique qu'aux hérédités \cle{récessives} : dans une hérédité dominante, tout porteur de l'allèle est atteint, il ne peut donc exister de porteur sain.
\end{enumerate}
\end{corrige}

\begin{corrige}[Exercice 4 — Calculs directs de risques]
\begin{enumerate}
\item Croisement $[A//a] \times [A//a]$. Gamètes de chaque parent : $1/2\,A$, $1/2\,a$.
\begin{center}
\begin{tabular}{|c|c|c|}\hline
\rowcolor{popblueL} & $A$ ($1/2$) & $a$ ($1/2$) \\ \hline
$A$ ($1/2$) & $[A//A]$ sain & $[A//a]$ sain \\ \hline
$a$ ($1/2$) & $[A//a]$ sain & $[a//a]$ \textbf{albinos} \\ \hline
\end{tabular}
\end{center}
Probabilité que l'enfant soit albinos : $\dfrac{1}{2}\times\dfrac{1}{2}=\dfrac{1}{4}=0,25$ (soit $25\,\%$). Probabilité qu'il soit sain : $1-\dfrac{1}{4}=\dfrac{3}{4}=0,75$ ($75\,\%$).
\item Croisement $X^{H}/X^{h} \times X^{H}Y$.
\begin{center}
\begin{tabular}{|c|c|c|}\hline
\rowcolor{popblueL} & $X^{H}$ ($1/2$) & $Y$ ($1/2$) \\ \hline
$X^{H}$ ($1/2$) & $X^{H}/X^{H}$ fille saine & $X^{H}Y$ garçon sain \\ \hline
$X^{h}$ ($1/2$) & $X^{H}/X^{h}$ fille conductrice & $X^{h}Y$ \textbf{garçon hémophile} \\ \hline
\end{tabular}
\end{center}
\begin{itemize}
\item Probabilité d'un garçon hémophile : $P(\text{garçon})\times P(X^{h}\text{ maternel})=\dfrac{1}{2}\times\dfrac{1}{2}=\dfrac{1}{4}$.
\item Probabilité d'une fille hémophile : $0$ (le père sain transmet obligatoirement $X^{H}$ à ses filles).
\item Probabilité d'un enfant hémophile, sexes confondus : $\dfrac{1}{4}=0,25$ ($25\,\%$).
\end{itemize}
\item Croisement $[O//O] \times [A//B]$. Gamètes du père : $100\,\%\,O$ ; gamètes de la mère : $1/2\,A$, $1/2\,B$.
\begin{itemize}
\item Enfants $[A//O]$ : groupe \textbf{A}, probabilité $1/2$ ;
\item Enfants $[B//O]$ : groupe \textbf{B}, probabilité $1/2$.
\end{itemize}
Aucun enfant ne peut être de groupe AB ni de groupe O.
\end{enumerate}
\end{corrige}

\begin{corrige}[Exercice 5 — Pedigree de l'albinisme]
\begin{enumerate}
\item Arbre généalogique :
\begin{popfigure}
\begin{tikzpicture}[scale=0.85]
% Génération I
\draw[thick,popblue] (0,3) rectangle (0.7,3.7); \node[above] at (0.35,3.7) {\small I-1};
\draw[thick,popblue] (1.7,3.35) circle (0.35); \node[above] at (1.7,3.7) {\small I-2};
\draw[thick,popblue] (0.7,3.35) -- (1.35,3.35);
% descendance vers II-3
\draw[thick,popblue] (1.05,3.35) -- (1.05,2.4);
\draw[thick,popblue] (1.05,2.4) -- (2.4,2.4);
\draw[thick,popblue] (2.4,2.4) -- (2.4,2.0);
% Génération II
\draw[thick,popblue] (2.05,1.3) rectangle (2.75,2.0); \node[left] at (2.05,1.65) {\small II-3};
\draw[thick,popblue] (3.75,1.65) circle (0.35); \node[right] at (4.1,1.65) {\small II-4};
\draw[thick,popblue] (2.75,1.65) -- (3.4,1.65);
% descendance III
\draw[thick,popblue] (3.05,1.65) -- (3.05,0.7);
\draw[thick,popblue] (1.0,0.7) -- (5.1,0.7);
\foreach \x in {1.0,2.35,3.75,5.1}{\draw[thick,popblue] (\x,0.7) -- (\x,0.3);}
\draw[thick,popblue] (0.8,-0.05) circle (0.35); \node[below] at (0.8,-0.4) {\small III-1};
\draw[thick,popblue,fill=popblue] (2.35,-0.05) circle (0.35); \node[below] at (2.35,-0.4) {\small III-2};
\draw[thick,popblue] (3.55,-0.05) circle (0.35); \node[below] at (3.55,-0.4) {\small III-3};
\draw[thick,popblue] (4.9,-0.25) rectangle (5.5,0.15); \node[below] at (5.2,-0.4) {\small III-4};
\end{tikzpicture}
\legende{Pedigree de la famille : I-1 et I-2 (sains) ; II-3 et II-4 (sains) ; III-1, III-3 (filles saines), III-4 (garçon sain), III-2 (fille albinos, symbole rempli).}
\end{popfigure}
\item \textbf{Mode de transmission : autosomique récessif.}
\begin{itemize}
\item II-3 et II-4 sont sains mais ont une fille atteinte : l'allèle était donc présent chez eux sans s'exprimer $\Rightarrow$ l'allèle est \cle{récessif} et les parents sont hétérozygotes (conducteurs).
\item Une fille (III-2) est atteinte alors que son père est sain : si le gène était lié à $X$, une fille atteinte $[a//a]$ aurait un père obligatoirement atteint ; le gène est donc porté par un \cle{autosome}.
\end{itemize}
\item Génotypes certains : II-3 : $[A//a]$ ; II-4 : $[A//a]$ (obligatoirement conducteurs puisqu'ils ont un enfant $[a//a]$) ; III-2 : $[a//a]$ (atteinte).
\item III-1 est saine, issue du croisement $[A//a]\times[A//a]$ : parmi les individus sains, $P([A//A])=\dfrac{1}{3}$ et $P([A//a])=\dfrac{2}{3}$. Son mari est albinos $[a//a]$.
\begin{itemize}
\item Si III-1 est $[A//A]$ : tous les enfants sont $[A//a]$, sains ; risque $=0$.
\item Si III-1 est $[A//a]$ : croisement $[A//a]\times[a//a]$ $\Rightarrow$ risque d'enfant albinos $=\dfrac{1}{2}$.
\end{itemize}
Risque global : $P=\dfrac{2}{3}\times\dfrac{1}{2}=\dfrac{1}{3}\approx 0,33$ (environ $33\,\%$).
\end{enumerate}
\end{corrige}

\begin{corrige}[Exercice 6 — Groupes sanguins ABO]
\begin{enumerate}
\item Un enfant de groupe O a le génotype $[O//O]$ : il a reçu un allèle $O$ de chaque parent. Deux parents de groupe A peuvent donc engendrer un enfant O s'ils sont tous deux hétérozygotes $[A//O]$.
\item Génotypes certains des parents : père $[A//O]$, mère $[A//O]$.
\item Échiquier de croisement :
\begin{center}
\begin{tabular}{|c|c|c|}\hline
\rowcolor{popblueL} & $A$ ($1/2$) & $O$ ($1/2$) \\ \hline
$A$ ($1/2$) & $[A//A]$ groupe A & $[A//O]$ groupe A \\ \hline
$O$ ($1/2$) & $[A//O]$ groupe A & $[O//O]$ groupe O \\ \hline
\end{tabular}
\end{center}
Groupes possibles pour les enfants suivants : groupe \textbf{A} avec probabilité $\dfrac{3}{4}$ ($75\,\%$) ; groupe \textbf{O} avec probabilité $\dfrac{1}{4}$ ($25\,\%$). Les groupes B et AB sont impossibles.
\item Le père est Rh$^{+}$ homozygote $[Rh^{+}//Rh^{+}]$, la mère Rh$^{-}$ $[Rh^{-}//Rh^{-}]$. Tous les enfants sont $[Rh^{+}//Rh^{-}]$, donc de phénotype \textbf{Rh$^{+}$}. Problème médical : lors d'une première grossesse, une mère Rh$^{-}$ portant un fœtus Rh$^{+}$ peut produire des anticorps anti-Rh ; lors d'une \textbf{seconde grossesse} avec un fœtus Rh$^{+}$, ces anticorps traversent le placenta et détruisent les hématies fœtales : c'est la \cle{maladie hémolytique du nouveau-né} (incompatibilité Rhésus fœto-maternelle).
\end{enumerate}
\end{corrige}

\begin{corrige}[Exercice 7 — Transmission du daltonisme]
\begin{enumerate}
\item Le daltonisme est une maladie à transmission \cle{gonosomique récessive liée au chromosome $X$}.
\item Génotypes :
\begin{itemize}
\item la femme est saine mais son père était daltonien ($X^{d}Y$) : elle a obligatoirement reçu $X^{d}$ de lui ; étant saine, elle est $X^{D}/X^{d}$ (\cle{conductrice}) ;
\item son père : $X^{d}Y$ (daltonien) ;
\item son mari, à vision normale : $X^{D}Y$.
\end{itemize}
\item Échiquier de croisement $X^{D}/X^{d} \times X^{D}Y$ :
\begin{center}
\begin{tabular}{|c|c|c|}\hline
\rowcolor{popblueL} & $X^{D}$ (père) & $Y$ (père) \\ \hline
$X^{D}$ (mère) & $X^{D}/X^{D}$ fille saine & $X^{D}Y$ garçon sain \\ \hline
$X^{d}$ (mère) & $X^{D}/X^{d}$ fille saine conductrice & $X^{d}Y$ \textbf{garçon daltonien} \\ \hline
\end{tabular}
\end{center}
Proportions phénotypiques : $1/4$ filles saines non conductrices, $1/4$ filles saines conductrices, $1/4$ garçons sains, $1/4$ garçons daltoniens. Autrement dit : toutes les filles sont saines ; parmi les garçons, $1/2$ sont daltoniens.
\item Une fille de ce couple \textbf{ne peut pas} être daltonienne : elle reçoit obligatoirement le $X^{D}$ de son père sain ; elle serait au pire conductrice $X^{D}/X^{d}$, donc de vision normale. Une fille daltonienne $X^{d}/X^{d}$ exigerait un père daltonien.
\end{enumerate}
\end{corrige}

\begin{corrige}[Exercice 8 — Rachitisme vitamino-résistant]
\begin{enumerate}
\item Pedigree :
\begin{popfigure}
\begin{tikzpicture}[scale=0.85]
% Génération I
\draw[thick,popblue,fill=popblue] (0,2.3) rectangle (0.7,3.0); \node[above] at (0.35,3.0) {\small I-1};
\draw[thick,popblue] (1.7,2.65) circle (0.35); \node[above] at (1.7,3.0) {\small I-2};
\draw[thick,popblue] (0.7,2.65) -- (1.35,2.65);
% descendance
\draw[thick,popblue] (1.05,2.65) -- (1.05,1.7);
\draw[thick,popblue] (-0.9,1.7) -- (3.0,1.7);
\foreach \x in {-0.9,0.075,1.05,2.025,3.0}{\draw[thick,popblue] (\x,1.7) -- (\x,1.3);}
\draw[thick,popblue,fill=popblue] (-0.9,0.95) circle (0.35); \node[below] at (-0.9,0.6) {\small II-1};
\draw[thick,popblue,fill=popblue] (0.075,0.95) circle (0.35); \node[below] at (0.075,0.6) {\small II-2};
\draw[thick,popblue,fill=popblue] (1.05,0.95) circle (0.35); \node[below] at (1.05,0.6) {\small II-3};
\draw[thick,popblue] (1.825,0.75) rectangle (2.225,1.15); \node[below] at (2.025,0.6) {\small II-4};
\draw[thick,popblue] (2.8,0.75) rectangle (3.2,1.15); \node[below] at (3.0,0.6) {\small II-5};
\end{tikzpicture}
\legende{Pedigree : I-1 homme atteint, I-2 femme saine ; II-1, II-2, II-3 filles atteintes ; II-4, II-5 fils sains.}
\end{popfigure}
\item \textbf{Compatibilité avec une transmission gonosomique dominante liée à $X$ :}
\begin{itemize}
\item \textbf{Dominante} : tous les enfants qui portent l'allèle sont atteints ; la maladie apparaît à chaque génération dès qu'un parent la transmet (ici, les filles atteintes ont un parent atteint).
\item \textbf{Liée à $X$} : le père atteint transmet son unique chromosome $X$ (porteur de $X^{R}$) à toutes ses filles, qui sont donc toutes atteintes ; il transmet son $Y$ à ses fils, qui ne reçoivent jamais son $X$ et sont donc tous sains. Cette répartition « toutes les filles atteintes, tous les fils sains » est la signature d'un gène dominant porté par le $X$ d'un père atteint.
\end{itemize}
\item Génotypes : I-1 : $X^{R}Y$ ; I-2 : $X^{r}/X^{r}$ ; filles II-1, II-2, II-3 : $X^{R}/X^{r}$ (hétérozygotes atteintes) ; fils II-4, II-5 : $X^{r}Y$ (sains).
\item Croisement d'une fille atteinte $X^{R}/X^{r}$ avec un homme sain $X^{r}Y$ :
\begin{center}
\begin{tabular}{|c|c|c|}\hline
\rowcolor{popblueL} & $X^{r}$ (père) & $Y$ (père) \\ \hline
$X^{R}$ (mère) & $X^{R}/X^{r}$ fille atteinte & $X^{R}Y$ garçon atteint \\ \hline
$X^{r}$ (mère) & $X^{r}/X^{r}$ fille saine & $X^{r}Y$ garçon sain \\ \hline
\end{tabular}
\end{center}
Proportions : $1/2$ des filles atteintes, $1/2$ des fils atteints ; globalement $1/2$ des enfants atteints, indépendamment du sexe (contrairement à la génération précédente, car c'est ici la mère qui transmet l'allèle).
\end{enumerate}
\end{corrige}

\begin{corrige}[Exercice 9 — Type Bac : albinisme sur quatre générations]
\begin{enumerate}
\item Pedigree :
\begin{popfigure}
\begin{tikzpicture}[scale=0.85]
% I
\draw[thick,popblue] (0,3.3) rectangle (0.7,4.0); \node[above] at (0.35,4.0) {\small I-1};
\draw[thick,popblue] (1.7,3.65) circle (0.35); \node[above] at (1.7,4.0) {\small I-2};
\draw[thick,popblue] (0.7,3.65) -- (1.35,3.65);
\draw[thick,popblue] (1.05,3.65) -- (1.05,2.7);
\draw[thick,popblue] (1.05,2.7) -- (2.4,2.7);
\draw[thick,popblue] (2.4,2.7) -- (2.4,2.3);
% II
\draw[thick,popblue,fill=popblue] (2.05,1.6) rectangle (2.75,2.3); \node[left] at (2.05,1.95) {\small II-3};
\draw[thick,popblue] (3.75,1.95) circle (0.35); \node[right] at (4.1,1.95) {\small II-4};
\draw[thick,popblue] (2.75,1.95) -- (3.4,1.95);
\draw[thick,popblue] (3.05,1.95) -- (3.05,1.0);
\draw[thick,popblue] (3.05,1.0) -- (3.05,0.6);
% III
\draw[thick,popblue] (3.05,0.25) circle (0.35); \node[left] at (2.7,0.25) {\small III-2};
\draw[thick,popblue] (4.4,-0.1) rectangle (5.1,0.6); \node[right] at (5.1,0.25) {\small III-3};
\draw[thick,popblue] (3.4,0.25) -- (4.4,0.25);
\draw[thick,popblue] (3.9,0.25) -- (3.9,-0.7);
\draw[thick,popblue] (3.9,-1.05) circle (0.35);
\node[below] at (3.9,-1.4) {\small ?};
\end{tikzpicture}
\legende{Pedigree : I-1 et I-2 sains ; II-3 albinos ; II-4 saine (père albinos) ; III-2 saine, III-3 sain ; enfant à naître (?).}
\end{popfigure}
\item \textbf{Mode de transmission : autosomique récessif.} Deux arguments :
\begin{itemize}
\item I-1 et I-2 sont sains et ont un fils atteint (II-3) : l'allèle ne s'exprime pas chez les parents qui le portent $\Rightarrow$ allèle \cle{récessif}.
\item Si le gène était lié à $X$, II-3 ($X^{a}Y$) aurait reçu $X^{a}$ de sa mère (possible), mais surtout le fait que des hommes et des femmes soient concernés sans liaison au sexe, et que II-3 atteint ait une fille saine, est compatible avec un autosome ; l'argument décisif classique : des parents sains des deux sexes engendrent un enfant atteint, et la maladie frappe indifféremment les sexes $\Rightarrow$ gène \cle{autosomique}.
\end{itemize}
\item Génotypes :
\begin{itemize}
\item II-3 est albinos : $[a//a]$ (génotype certain).
\item I-1 et I-2 sont sains mais ont un fils $[a//a]$ : chacun a fourni un allèle $a$ $\Rightarrow$ I-1 : $[A//a]$, I-2 : $[A//a]$ (conducteurs).
\item II-4 est saine mais son père était albinos $[a//a]$ : elle a obligatoirement reçu $a$ de lui $\Rightarrow$ II-4 : $[A//a]$ (conductrice).
\end{itemize}
\item Risque pour le couple III-2 $\times$ III-3. III-2 est saine, fille de II-3 $[a//a]$ : elle a reçu $a$ de son père, donc III-2 est certainement $[A//a]$. Deux hypothèses pour III-3 :
\begin{itemize}
\item \textbf{Hypothèse 1 : III-3 est $[A//A]$} (sans antécédent familial, hypothèse la plus probable) : croisement $[A//a]\times[A//A]$ $\Rightarrow$ tous les enfants sains ; risque d'enfant albinos $=0$ (mais $1/2$ des enfants seront conducteurs).
\item \textbf{Hypothèse 2 : III-3 est conducteur $[A//a]$} : croisement $[A//a]\times[A//a]$ $\Rightarrow$ risque $=\dfrac{1}{4}=25\,\%$.
\end{itemize}
\item \textbf{Conseil génétique} : « Votre famille présente un cas d'albinisme : Madame est conductrice certaine de l'allèle. Un \cle{examen prénuptial} (analyse génétique) permettrait de déterminer si Monsieur est lui-même conducteur. Si tel est le cas, chaque enfant aurait un risque de $25\,\%$ d'être albinos ; un \cle{diagnostic prénatal} pourrait alors être proposé pendant la grossesse pour connaître le statut du fœtus et préparer sa prise en charge. L'albinisme n'empêche pas une vie normale mais impose une protection solaire rigoureuse et un suivi ophtalmologique et dermatologique. »
\end{enumerate}
\textbf{Barème indicatif (sur 10)} : pedigree correct (2 pts) ; mode de transmission avec 2 arguments (2 pts) ; génotypes justifiés (2 pts) ; calcul du risque avec distinction des hypothèses (2 pts) ; conseil génétique rédigé (2 pts).\par
\textbf{Points de vigilance} : justifier le caractère récessif par « parents sains, enfant atteint » ; justifier le caractère autosomique (et non lié à $X$) ; ne pas oublier que III-2 est conductrice \emph{certaine} car son père est $[a//a]$ ; distinguer clairement les deux hypothèses sur III-3 dans le calcul du risque.
\end{corrige}

\begin{corrige}[Exercice 10 — Type Bac : électrophorèse de l'hémoglobine et drépanocytose]
\begin{enumerate}
\item Génotypes déduits des profils d'électrophorèse :
\begin{itemize}
\item I-1 : deux bandes (HbA + HbS) $\Rightarrow$ $[HbA//HbS]$ (porteur du trait, sain) ;
\item I-2 : deux bandes (HbA + HbS) $\Rightarrow$ $[HbA//HbS]$ ;
\item II-1 : une seule bande HbA $\Rightarrow$ $[HbA//HbA]$ (saine, non porteuse) ;
\item II-2 : une seule bande HbS $\Rightarrow$ $[HbS//HbS]$ (drépanocytaire).
\end{itemize}
Justification : chaque bande révèle un type d'hémoglobine, donc un allèle exprimé ; un individu homozygote ne produit qu'un seul type d'hémoglobine (une bande), l'hétérozygote en produit deux (deux bandes), car les allèles $HbA$ et $HbS$ sont \cle{codominants} au niveau biochimique.
\item \textbf{Mode de transmission : autosomique, avec récessivité au niveau clinique} (la maladie ne se déclare que chez $[HbS//HbS]$) et codominance au niveau moléculaire. Les hétérozygotes sont dits « sains » car la quantité d'HbA produite suffit à assurer une oxygénation normale des tissus : l'allèle $HbS$ ne s'exprime cliniquement qu'à l'état homozygote ; ils ne présentent pas les crises drépanocytaires, bien qu'ils possèdent et transmettent l'allèle.
\item Croisement $[HbA//HbS]\times[HbA//HbS]$ :
\begin{center}
\begin{tabular}{|c|c|c|}\hline
\rowcolor{popblueL} & $HbA$ ($1/2$) & $HbS$ ($1/2$) \\ \hline
$HbA$ ($1/2$) & $[HbA//HbA]$ sain & $[HbA//HbS]$ porteur \\ \hline
$HbS$ ($1/2$) & $[HbA//HbS]$ porteur & $[HbS//HbS]$ \textbf{drépanocytaire} \\ \hline
\end{tabular}
\end{center}
Risque que l'enfant soit drépanocytaire : $\dfrac{1}{4}=25\,\%$ ; risque qu'il soit porteur du trait : $\dfrac{1}{2}=50\,\%$.
\item \textbf{Maintien de l'allèle $HbS$ en zone d'endémie palustre} : les hétérozygotes $[HbA//HbS]$ possèdent une \cle{résistance accrue au paludisme} (le plasmodium se développe mal dans leurs hématies, qui déforment facilement). Dans les zones où sévit le paludisme, comme au Cameroun, les hétérozygotes survivent mieux que les homozygotes $[HbA//HbA]$ et transmettent l'allèle $HbS$ : c'est un \cle{avantage sélectif de l'hétérozygote} (polymorphisme équilibré), qui compense l'élimination des homozygotes $[HbS//HbS]$.
\item \textbf{Examen prénuptial} : l'électrophorèse de l'hémoglobine, réalisée avant le mariage, permet d'identifier les porteurs du trait. Un couple dont les deux membres sont $[HbA//HbS]$ sait alors que chaque enfant a $25\,\%$ de risque d'être drépanocytaire ; informé, il peut faire un choix éclairé (diagnostic prénatal, suivi médical renforcé). \textbf{Conseil génétique} : « Vous êtes tous deux porteurs sains du trait drépanocytaire : vous ne serez jamais malades, mais à chaque grossesse, l'enfant a $25\,\%$ de risque d'être drépanocytaire, $50\,\%$ d'être porteur comme vous et $25\,\%$ d'être non porteur. Un diagnostic prénatal peut être envisagé, et tout enfant atteint devra bénéficier d'un suivi médical régulier (vaccinations, hydratation, prise en charge des crises) dans un centre spécialisé. »
\end{enumerate}
\textbf{Barème indicatif (sur 10)} : génotypes justifiés (2 pts) ; mode de transmission et statut des hétérozygotes (2 pts) ; échiquier et calculs de risques (2 pts) ; explication de l'avantage sélectif (2 pts) ; rôle de l'examen prénuptial et conseil génétique (2 pts).\par
\textbf{Points de vigilance} : bien distinguer « porteur du trait » (sain) et « drépanocytaire » (malade) ; citer la codominance moléculaire pour justifier les deux bandes ; relier explicitement la fréquence de $HbS$ à la pression sélective du paludisme ; ne pas affirmer que l'examen prénuptial « guérit » la maladie : il informe et permet de prévoir.
\end{corrige}

\begin{corrige}[Exercice 11 — Type Bac : hémophilie et conseil génétique]
\begin{enumerate}
\item \textbf{Mode de transmission : gonosomique récessif lié au chromosome $X$.} Les hommes sont beaucoup plus souvent atteints car ils ne possèdent qu'\textbf{un seul} chromosome $X$ : il leur suffit de recevoir un $X^{h}$ pour être hémophiles ($X^{h}Y$). La femme, possédant deux $X$, n'est atteinte que si elle est $X^{h}/X^{h}$, situation rare qui exige un père hémophile et une mère au moins conductrice ; l'hétérozygote $X^{H}/X^{h}$ est saine (conductrice).
\item Génotypes :
\begin{itemize}
\item le mari est hémophile : $X^{h}Y$ (génotype certain, un seul $X$) ;
\item le père de la femme était hémophile : $X^{h}Y$ ;
\item la femme est saine mais a obligatoirement reçu le $X^{h}$ de son père : elle est donc $X^{H}/X^{h}$ (conductrice certaine).
\end{itemize}
\item Échiquier de croisement $X^{H}/X^{h} \times X^{h}Y$ :
\begin{center}
\begin{tabular}{|c|c|c|}\hline
\rowcolor{popblueL} & $X^{h}$ (père) & $Y$ (père) \\ \hline
$X^{H}$ (mère) & $X^{H}/X^{h}$ fille conductrice saine & $X^{H}Y$ garçon sain \\ \hline
$X^{h}$ (mère) & $X^{h}/X^{h}$ \textbf{fille hémophile} & $X^{h}Y$ \textbf{garçon hémophile} \\ \hline
\end{tabular}
\end{center}
\begin{itemize}
\item Probabilité d'un garçon hémophile : $\dfrac{1}{2}\times\dfrac{1}{2}=\dfrac{1}{4}=25\,\%$ ;
\item Probabilité d'une fille hémophile : $\dfrac{1}{2}\times\dfrac{1}{2}=\dfrac{1}{4}=25\,\%$ ;
\item Probabilité d'une fille conductrice : $\dfrac{1}{4}=25\,\%$ ;
\item Probabilité globale d'un enfant hémophile : $\dfrac{1}{4}+\dfrac{1}{4}=\dfrac{1}{2}=50\,\%$.
\end{itemize}
\item \textbf{Non, les risques ne sont pas modifiés} pour le deuxième enfant : chaque grossesse est un événement indépendant (les gamètes se forment à nouveau par méiose, avec les mêmes probabilités de $1/2$). Le résultat de la première naissance n'influence pas la seconde.
\item \textbf{Intérêt des examens prénataux} : le caryotype permet de connaître le sexe chromosomique du fœtus et la recherche moléculaire de l'allèle $X^{h}$ permet de déterminer si le fœtus est porteur ; le couple peut ainsi anticiper la prise en charge médicale d'un enfant hémophile (prévention des hémorragies, apports en facteur anti-hémophilique). \textbf{Limites éthiques} : le conseil génétique doit rester \emph{informatif et non directif} : il appartient au couple seul de décider ; ces examens posent des questions délicates (interruption médicale de grossesse, risque d'eugénisme, secret médical, stigmatisation des porteurs) qui imposent un accompagnement respectueux de la dignité des personnes.
\end{enumerate}
\textbf{Barème indicatif (sur 10)} : mode de transmission et explication de la prédominance masculine (2 pts) ; génotypes justifiés (2 pts) ; échiquier et quatre probabilités (3 pts) ; indépendance des grossesses (1 pt) ; intérêt des examens prénataux et limites éthiques (2 pts).\par
\textbf{Points de vigilance} : ne pas confondre « probabilité d'un garçon hémophile » ($1/4$, sexes confondus) et « probabilité qu'un garçon soit hémophile » ($1/2$, conditionnelle au sexe) ; justifier que la femme est conductrice \emph{certaine} grâce à son père hémophile ; rappeler que les tirages successifs sont indépendants ; pour la question éthique, exiger un propos nuancé (information, liberté de choix, respect de la personne).
\end{corrige}

\newpage

\coursec{Fiche bilan}

\begin{popfigure}
\begin{tikzpicture}[mindmap, grow cyclic, every node/.style={concept, font=\small},
  root concept/.append style={concept color=popblue, font=\bfseries},
  level 1/.append style={level distance=3.4cm, sibling angle=60, font=\scriptsize},
  level 2/.append style={level distance=2.2cm, sibling angle=45, font=\tiny}]
\node[root concept, text=white]{Maladies géniques et chromosomiques : prévision}
  child[concept color=popgreen]{ node {Arbre généalogique}
    child { node {cercle = femme, carré = homme} }
    child { node {rempli = atteint} }
  }
  child[concept color=poporange]{ node {Autosomique récessive}
    child { node {albinisme, drépanocytose} }
    child { node {saute des générations ; parents sains hétérozygotes} }
  }
  child[concept color=poppurple]{ node {Autosomique dominante}
    child { node {achondroplasie} }
    child { node {toutes les générations ; 1 parent atteint} }
  }
  child[concept color=popgold]{ node {Codominance}
    child { node {groupes ABO : $I^A$, $I^B$ codominants, $i$ récessif} }
    child { node {Rhésus : Rh$^+$/Rh$^-$} }
  }
  child[concept color=poppink]{ node {Gonosomique (X)}
    child { node {daltonisme, hémophilie (récessifs)} }
    child { node {rachitisme vitamino-résistant (dominant)} }
  }
  child[concept color=popblue!60]{ node {Risque génétique}
    child { node {échiquier de croisement} }
    child { node {probabilités : 1/4, 1/2, 3/4} }
  };
\end{tikzpicture}
\legende{Carte mentale de la leçon : les six branches à maîtriser pour prévoir la transmission d'une maladie héréditaire dans une famille.}
\end{popfigure}

\begin{bilan}
\textbf{Définitions clés}
\begin{itemize}
  \item \cle{Arbre généalogique} (pedigree) : représentation schématique, sur plusieurs générations, de la transmission d'un caractère dans une famille. Symboles : \textbf{cercle} = femme, \textbf{carré} = homme, symbole \textbf{rempli} = individu atteint, trait horizontal = union, trait vertical = descendance.
  \item \cle{Maladie génique} : due à la mutation d'un gène (albinisme, hémophilie). \cle{Maladie chromosomique} : due à une anomalie de nombre ou de structure des chromosomes (trisomie 21).
  \item \cle{Autosomique} : gène porté par un autosome (chromosome non sexuel). \cle{Gonosomique} : gène porté par un chromosome sexuel (presque toujours X).
  \item \cle{Récessif} : l'allèle morbide ne s'exprime qu'à l'état homozygote. \cle{Dominant} : il s'exprime dès l'état hétérozygote. \cle{Codominant} : deux allèles s'expriment simultanément chez l'hétérozygote.
  \item \cle{Risque génétique} : probabilité, pour un couple, de donner naissance à un enfant atteint.
\end{itemize}

\textbf{Maladies et modes de transmission à connaître par cœur}
\begin{center}
\begin{tabularx}{\linewidth}{|l|X|X|}\hline
\rowcolor{popblueL} \textbf{Maladie} & \textbf{Transmission} & \textbf{Indices sur le pedigree} \\ \hline
Albinisme & autosomique récessive & parents sains, enfant atteint ; saute des générations ; touche les deux sexes \\ \hline
Drépanocytose & autosomique récessive ($Hb^S/Hb^S$) & idem ; $Hb^A/Hb^S$ = porteur sain (résistant au paludisme) \\ \hline
Achondroplasie & autosomique dominante & présente à chaque génération ; un parent atteint au moins \\ \hline
Daltonisme & gonosomique récessive (X) & surtout des garçons ; père atteint $\to$ filles conductrices \\ \hline
Hémophilie & gonosomique récessive (X) & idem ; transmission par les femmes conductrices \\ \hline
Rachitisme vitamino-résistant & gonosomique dominante (X) & père atteint $\to$ toutes les filles atteintes, aucun fils \\ \hline
Groupes ABO & codominance $I^A$/$I^B$, $i$ récessif & [AB] = $I^A/I^B$ ; [O] = $i/i$ \\ \hline
Rhésus & Rh$^+$ dominant sur Rh$^-$ & Rh$^-$ = homozygote récessif \\ \hline
\end{tabularx}
\end{center}

\textbf{Méthodes express}
\begin{enumerate}
  \item \textbf{Mode de transmission} : des parents sains ont un enfant atteint $\Rightarrow$ allèle \emph{récessif} ; la maladie apparaît à chaque génération $\Rightarrow$ \emph{dominante} ; seuls (ou presque) les garçons sont atteints $\Rightarrow$ \emph{gonosomique récessive liée à X}.
  \item \textbf{Test autosome/gonosome} : si le père atteint transmet à toutes ses filles mais à aucun fils $\Rightarrow$ gène lié à X ; si les deux sexes sont touchés également $\Rightarrow$ autosomique.
  \item \textbf{Risque génétique} : écrire les génotypes, dresser l'échiquier de croisement, compter les cases favorables : parents hétérozygotes $A/a \times A/a$ $\Rightarrow$ 1/4 atteint, 1/2 porteur, 1/4 sain ; femme conductrice $X^H/X^h \times$ homme sain $\Rightarrow$ 1/2 des garçons atteints.
  \item \textbf{Prévention} : examen \cle{prénuptial} (avant mariage : groupes sanguins, dépistage drépanocytose, VIH) et examen \cle{prénatal} (pendant la grossesse : échographie, amniocentèse, caryotype).
\end{enumerate}
\end{bilan}

\begin{aretenir}[Checklist avant le Bac]
\begin{itemize}
  \item[$\square$] Je sais dessiner et lire un pedigree (cercle, carré, symbole rempli, traits d'union et de descendance).
  \item[$\square$] Je sais distinguer maladie génique et maladie chromosomique.
  \item[$\square$] Je sais reconnaître une transmission autosomique récessive (albinisme, drépanocytose).
  \item[$\square$] Je sais reconnaître une transmission autosomique dominante (achondroplasie).
  \item[$\square$] Je connais les génotypes des groupes ABO et la codominance $I^A/I^B$.
  \item[$\square$] Je sais que Rh$^+$ domine Rh$^-$ et je sais expliquer l'incompatibilité Rhésus mère--fœtus.
  \item[$\square$] Je sais reconnaître une transmission gonosomique récessive (daltonisme, hémophilie) et le rôle des femmes conductrices.
  \item[$\square$] Je sais que le rachitisme vitamino-résistant est dominant lié à X : père atteint $\Rightarrow$ toutes les filles atteintes.
  \item[$\square$] Je sais construire un échiquier de croisement et calculer un risque génétique (1/4, 1/2, 3/4).
  \item[$\square$] Je justifie toujours le mode de transmission par des arguments tirés du pedigree.
  \item[$\square$] Je sais expliquer l'intérêt des examens prénuptiaux et prénataux (amniocentèse, caryotype).
  \item[$\square$] J'écris les génotypes correctement : $X^H/X^h$, $I^A/i$, $Hb^A/Hb^S$, avec l'allèle morbide noté clairement.
\end{itemize}
\end{aretenir}

\findecours

\end{document}
